%!TEX TS-program = xelatex
%!TEX TS-options = -shell-escape
%!TEX root = ../Topo2_WS15/topologie_2.tex
\RequirePackage{fix-cm} 
\documentclass[a4paper, twoside, headsepline, index=totoc,toc=listof, fontsize=10pt, cleardoublepage=empty, headinclude, DIV=12, BCOR=5mm, titlepage,draft]{scrartcl}
\usepackage{scrtime} % KOMA, Uhrzeit ermoeglicht

\usepackage{etoolbox}
\usepackage{letltxmacro}
\usepackage{ifthen}

%--Farbdefinitionen
%-- muss vor tikz geladen werden
\usepackage[usenames, table, x11names]{xcolor}
\definecolor{dark_gray}{gray}{0.45}
\definecolor{light_gray}{gray}{0.6}
\usepackage[final]{graphicx}

%--Zum Zeichnen
%-- muss vor polyglossia bzw. babel geladen werden
\usepackage{tikz}
\usepackage{tikz-cd}
\usetikzlibrary{external}
\tikzset{>=latex}
\usetikzlibrary{shapes,arrows.meta,intersections}
\usetikzlibrary{calc,3d}
\usetikzlibrary{decorations.pathreplacing,decorations.markings, decorations.pathmorphing}
\usetikzlibrary{angles}

%-- Konfiguration von tikzexternalize
\tikzexternalize[prefix=tikz/,up to date check=diff]
\pgfkeys{/pgf/images/include external/.code=\includegraphics{#1}}
\tikzset{external/system call={xelatex \tikzexternalcheckshellescape -halt-on-error -interaction=batchmode --shell-escape -jobname "\image" "\texsource"}}


%-- tikzexternalize fuer tikzcd deaktivieren, da inkompatibel
\AtBeginEnvironment{tikzcd}{\tikzexternaldisable}
\AtEndEnvironment{tikzcd}{\tikzexternalenable}

%-- um Inkompatibilitaeten von quotes und polyglossia bzw. babel zu vermeiden
\tikzset{
  every picture/.append style={
    execute at begin picture={\shorthandoff{"}},
    execute at end picture={\shorthandon{"}}
  }
}
\usetikzlibrary{quotes}
\usepackage{pgfplots}
\usepgfplotslibrary{colormaps}
% \pgfplotsset{compat=1.12}



%-- Mathesymbole etc.
\usepackage{mathtools} % beinhaltet amsmath
\mathtoolsset{showonlyrefs, centercolon,showmanualtags}
\newtagform{brackets}[\textbf]{[}{]}
\usetagform{brackets}
\usepackage{fix-cm}
% \usepackage{wasysym} % zusätzliche Symbole
\usepackage{amssymb} % zusätzliche Symbole
% \usepackage{stmaryrd} % für Blitz
\usepackage{nicefrac} % schräge Brüche
\usepackage{faktor}
\newcommand{\Faktor}[1]{\faktor[\textstyle]{#1}}
\usepackage{xfrac}
\usepackage{cancel}
\usepackage{mathdots} % Verbesserung von Punkten wie zB \ldots
\usepackage[bb=px]{mathalfa}

%-- einzelne Symbole aus dem mathx-package
% \DeclareFontFamily{U}{mathx}{\hyphenchar\font45}
% \DeclareFontShape{U}{mathx}{m}{n}{<-> mathx10}{}
% \DeclareSymbolFont{mathx}{U}{mathx}{m}{n}
% \DeclareMathSymbol{\bigplus}{1}{mathx}{"90}

%-- 
\usepackage{silence}
\WarningFilter{latexfont}{Size substitutions with differences}
\WarningFilter{latexfont}{Font shape `U/bbold/m/n' in size}

\DeclareSymbolFont{bbold}{U}{bbold}{m}{n}
\DeclareSymbolFontAlphabet{\mathbbold}{bbold}
\newcommand{\ind}{\mathbbold{1}} % charakteristische-Funktion-Eins

\def\mathul#1#2{\color{#1}\underline{{\color{black}#2}}\color{black}} %farbiges Untersteichen im Mathe-Modus
\renewcommand{\le}{\leqslant}
\renewcommand{\ge}{\geqslant}

%-- Underbrace als Befehl in LaTeX-Syntax (und ohne Spacingprobleme mit nachfolgenden Operatoren...)
\newcommand{\Underbrace}[2]{{\underbrace{#1}_{#2}}}
\newcommand{\Underbracket}[2]{{\underbracket[0.7pt][2pt]{#1}_{#2}}}

%-- Alles was mit Schrift und XeTeX zu tun hat
\usepackage{palatino}
\usepackage[euler-digits]{eulervm}
\usepackage[no-math]{fontspec}
\usepackage{polyglossia} % moderner babel-ersatz
\setmainlanguage[spelling=new,babelshorthands=true]{german}
\shorthandoff{"}
\setotherlanguage{english}
\defaultfontfeatures{Mapping=tex-text, WordSpace={1.2}, Ligatures={Required,Common,Contextual}} %

%-- Alle hier aufgeführten Schriftarten sind in jeder TeXlive-Installation vorhanden!!
%-- Hauptschriftart festlegen
% \setmainfont{texgyrepagella}[Extension=.otf,UprightFont=*-regular,BoldFont=*-bold,ItalicFont=*-italic,BoldItalicFont=*-bolditalic,ItalicFeatures={Style=Historic}]
\setmainfont{TeXGyrePagellaX}[Extension=.otf,UprightFont=*-Regular,BoldFont=*-Bold,ItalicFont=*-Italic,BoldItalicFont=*-BoldItalic,ItalicFeatures={Style=Historic},Ligatures={Required,Common,Contextual,Historic}]
% \setmainfont[BoldFont={* Bold}, ItalicFont={* Light Italic},SmallCapsFont={Linux Libertine O}, SmallCapsFeatures={Letters=SmallCaps}]{Source Sans Pro}


%-- Serifenlose Schriftart festlegen
% \setsansfont{Roboto}[Scale=MatchUppercase,Extension=.ttf,UprightFont=*-Regular,BoldFont=*-Bold,ItalicFont=*-RegularItalic,BoldItalicFont=*-MediumItalic]
% \setsansfont{RobotoSlab}[Scale=MatchUppercase,Extension=.ttf,UprightFont=*-Regular,BoldFont=*-Bold,ItalicFont=*-Thin,SmallCapsFont=*-Regular]
% \setsansfont{FiraSans}[Scale=MatchUppercase,Extension=.otf, UprightFont=*-Regular, BoldFont=*-SemiBold, ItalicFont=*-RegularItalic, BoldItalicFont=*-MediumItalic]
% \setsansfont{SourceSansPro}[Scale=MatchUppercase,Extension=.otf,UprightFont=*-Regular,BoldFont=*-Semibold,ItalicFont=*-LightIt,BoldItalicFont=*-SemiboldIt]
% \setsansfont{texgyreadventor}[Scale=MatchUppercase, Extension=.otf, UprightFont=*-regular, BoldFont=*-bold, ItalicFont=*-italic,BoldItalicFont=*-bolditalic]
% \setsansfont{Kurier}[Scale=MatchUppercase, Extension=.otf, UprightFont=*Medium-Regular, BoldFont=*Heavy-Regular, ItalicFont=*Medium-Italic,BoldItalicFont=*Heavy-Italic]
% \setsansfont{Iwona}[Scale=MatchUppercase, Extension=.otf, UprightFont=*Medium-Regular, BoldFont=*Heavy-Regular, ItalicFont=*Medium-Italic,BoldItalicFont=*Heavy-Italic]
% \setsansfont{LinBiolinum}[Scale=MatchUppercase, Extension=.otf, UprightFont=*_R, BoldFont=*_RB, ItalicFont=*_RI,BoldItalicFont=*_RBO]
% \setsansfont{texgyrebonum}[Scale=MatchUppercase,Extension=.otf, UprightFont=*-regular, BoldFont=*-bold, ItalicFont=*-italic, BoldItalicFont=*-bolditalic]
\setsansfont{texgyreadventor}[Scale=MatchUppercase,Extension=.otf, UprightFont=*-regular, BoldFont=*-bold, ItalicFont=*-italic, BoldItalicFont=*-bolditalic]



% \newfontfamily{\geometric}{texgyreheros}[Scale=MatchLowercase, Extension=.otf, UprightFont=*-regular, BoldFont=*-bold, ItalicFont=*-italic]
% \addtokomafont{disposition}{\geometric}
\setmonofont{SourceCodePro}[Scale=0.9,Extension=.otf,UprightFont=*-Regular, BoldFont=*-Semibold, ItalicFont=*-Light]
\usepackage{xltxtra}
\usepackage{fontawesome}
\usepackage[final]{microtype}


\usepackage[draft=false]{scrlayer-scrpage} % wie fancyhdr, nur optimiert auf KOMA-Skript, leicht andere Syntax

\flushbottom


%--Mixed
\usepackage[shortlabels,inline]{enumitem}
\setlist[itemize,1]{label=\faCaretRight}
\setlist[enumerate]{font=\bfseries}
\setlist[description]{font=\normalfont\bfseries}
\usepackage[german=quotes]{csquotes}
\usepackage{makeidx}
\usepackage{wrapfig}
\usepackage{float}
\usepackage[margin=10pt, font=small, labelfont={sf, bf}, format=plain, indention=1em]{caption}
\captionsetup[wrapfigure]{name=Abb. }
\usepackage{stackrel}


%--Unterstreichung
\usepackage[normalem]{ulem}
\setlength{\ULdepth}{1.8pt}

%--Indexverarbeitung
\newcommand{\bet}[1]{\emph{#1}}
\newcommand{\Index}[1]{\bet{#1}\index{#1}}
\makeindex
\setindexpreamble{{\noindent\sffamily\small Die \emph{Seitenzahlen} sind mit Hyperlinks versehen und somit anklickbar}  \par \bigskip}
\renewcommand{\indexpagestyle}{scrheadings}


%--Marginnote & todonotes
\deffootnote[1.5em]{1.5em}{1.5em}{\textsuperscript{\thefootnotemark}\ }
\usepackage[fulladjust]{marginnote}
\renewcommand*{\marginfont}{\color{dark_gray} \itshape \footnotesize}
\usepackage[textsize=small]{todonotes}
\usepackage{ragged2e}
\renewcommand*{\raggedleftmarginnote}{\RaggedLeft}
\renewcommand*{\raggedrightmarginnote}{\RaggedRight}

%--Todonotes mit tikz/externalize kompatibel machen
\LetLtxMacro{\oldtodo}{\todo}
\renewcommand{\todo}[2][]{\tikzexternaldisable\oldtodo[#1]{#2}\tikzexternalenable}
\LetLtxMacro{\oldmissingfigure}{\missingfigure}
\renewcommand{\missingfigure}[2][]{\tikzexternaldisable\oldmissingfigure[{#1}]{#2}\tikzexternalenable}

%--Konfiguration von Hyperref pdfstartview=FitH, 
\usepackage[hidelinks, pdfpagelabels,  bookmarksopen=true, bookmarksnumbered=true, linkcolor=black, urlcolor=SkyBlue2, plainpages=false,pagebackref, citecolor=black, hypertexnames=true, pdfauthor={Jannes Bantje}, pdfborderstyle={/S/U}, linkbordercolor=SkyBlue2, colorlinks=false,final]{hyperref}


\newcommand{\appendLink}[1]{#1\,\faExternalLink}
\newcommand{\hrefsym}[2]{\href{#1}{\texttt{\appendLink{#2}}}}
\newcommand{\hrefsymmail}[2]{\href{#1}{\texttt{\faEnvelopeO\,#2}}}
\renewcommand{\url}[1]{\hrefsym{#1}{\nolinkurl{#1}}}



% -- QR-Codes
\usepackage{qrcode} %-- hinter hyperref laden!

%--Römische Zahlen
\newcommand{\RM}[1]{\MakeUppercase{\romannumeral #1{}}}

%%--Abkürzungen etc.
\newcommand{\light}{\text{\Large $\lightning$}}

%-- Definitionen von weiteren Mathe-Befehlen

\DeclareMathOperator{\im}{im} %Bild
\DeclareMathOperator{\id}{id} %identische Abbildung
\DeclareMathOperator{\conj}{conj} %Konjugation
\DeclareMathOperator{\sgn}{sgn} %Signum
\DeclareMathOperator{\End}{End} %Endomorphismen
\DeclareMathOperator{\Hom}{Hom} %Homomorphismen
\DeclareMathOperator{\Iso}{Iso} %Isomorphismen
\DeclareMathOperator{\Span}{span} %Span

\DeclareMathOperator{\grad}{grad} %Gradient
\DeclareMathOperator{\dive}{div} %Gradient
\DeclareMathOperator{\rot}{rot} %Rotation


\DeclareMathOperator{\Char}{char} %Charakteristik
\DeclareMathOperator{\Aut}{Aut} %Automorphismen

% -- Zum Finetuning von Befehlen
\makeatletter
\newcommand{\raisemath}[1]{\mathpalette{\raisem@th{#1}}}
\newcommand{\raisem@th}[3]{\raisebox{#1}{$#2#3$}}
\makeatother

\newcommand{\sing}{{\raisemath{1.1pt}{\scriptscriptstyle\mathrm{sing}}}}
\newcommand{\pt}{\mathrm{pt}}
\newcommand{\op}{\mathrm{op}}
\DeclareMathOperator{\Sp}{Sp}
\DeclareMathOperator{\Keg}{Keg}
\newcommand{\slashedi}{i\hspace{-3.5pt}/}
\newcommand{\cupp}{\smallsmile}
\newcommand{\capp}{\smallfrown}
\DeclareMathOperator*{\colim}{colim}


\newcommand{\D}{\ensuremath{\mathrm{D}\mkern-1.5mu}}
\newcommand{\mathd}{\ensuremath{\mathrm{d}\mkern-1.0mu}}
\newcommand{\Tmap}{\ensuremath{\mathrm{T}\mkern-0.85mu}}
\newcommand{\opL}{\ensuremath{\mathrm{L}\mkern-0.6mu}}


\DeclareMathOperator{\supp}{supp} %Träger


% -- Befehle für FunkAna und Operatoralgebren
\newcommand{\tr}{\mathrm{tr}}
\newcommand{\w}{\mathrm{w}}
\newcommand{\sa}{\mathrm{sa}}
\newcommand{\so}{\mathrm{s.o.}}
\newcommand{\weakT}[1]{\ensuremath{\mathcal{T}_{#1}^{\mkern+1.0mu\text{\raisebox{0.4ex}{$\mathrm{w}$}}}}}
\newcommand{\weakTstar}[1]{\ensuremath{\mathcal{T}_{#1}^{\mkern+1.0mu\text{\raisebox{0.4ex}{$\mathrm{w}$}}^*}}}
\newcommand{\finSub}{\subset\mkern-0.7mu \subset}
\DeclareMathOperator{\Inv}{Inv}
\DeclareMathOperator{\dist}{dist}
\newcommand{\simm}{{\hspace{-1.6pt}\raisemath{0.5pt}{\sim}}}
\DeclareMathOperator{\ev}{ev}
\DeclareMathOperator{\Alg}{Alg}


% -- Kategorien
\newcommand{\TOP}{\textsc{Top}}
\newcommand{\HTOP}{\textsc{HTop}}
\newcommand{\VR}{\textsc{VR}}
\newcommand{\MOD}{\textsc{Mod}}
\newcommand{\MONOIDE}{\textsc{Monoide}}
\newcommand{\MENGEN}{\textsc{Mengen}}
\newcommand{\MAN}{\textsc{Man}}
\newcommand{\GRUPPEN}{\textsc{Gruppen}}
\newcommand{\ABELGRUPPEN}{\textsc{Abel.Gruppen}}
\newcommand{\KAT}{\textsc{Kat}}
\newcommand{\FUN}{\textsc{Fun}}
\newcommand{\SIMP}{\textsc{Simp}}
\newcommand{\VEKT}{\textsc{Vekt}}

\DeclareMathOperator{\cell}{cell}
\DeclarePairedDelimiter{\homologieklasse}{\llbracket}{\rrbracket}
\newcommand{\rand}[1]{\ensuremath{\partial^{\scriptscriptstyle #1}}}


% -- theorem packages
\usepackage{amsthm}
\usepackage{thmtools,thm-restate}
\renewcommand{\listtheoremname}{Übersicht aller Aussagen}

% -- Theoreme als PDF-Lesezeichen
\usepackage{bookmark}
\bookmarksetup{open,numbered}
\makeatletter
\newcommand*{\theorembookmark}{%
  \bookmark[
    dest=\@currentHref,
    rellevel=1,
    keeplevel,
  ]{%
    \thmt@thmname\space\csname the\thmt@envname\endcsname
    \ifx\thmt@shortoptarg\@empty
    \else
      \space(\thmt@shortoptarg)%
    \fi
  }%
}   
\makeatother

% -- Definition der einzelnen Umgebungen
\declaretheoremstyle[%
	headfont=\sffamily\bfseries,
	notefont=\normalfont\sffamily\scshape,
	bodyfont=\normalfont,
	headformat=\NUMBER\ \NAME\NOTE,
	headpunct=.,
	postheadspace=1em,
	spaceabove=15pt,spacebelow=10pt,
	postheadhook=\theorembookmark]%
{mainstyle}
\declaretheoremstyle[%
	headfont=\sffamily\bfseries,
	notefont=\normalfont\sffamily\scshape,
	bodyfont=\normalfont,
	headformat=\NUMBER\NAME\NOTE,
	headpunct=.,
	postheadspace=1em,
	spaceabove=15pt,spacebelow=10pt,
	postheadhook=\theorembookmark]%
{mainstyle_unnumbered}
\declaretheoremstyle[%
	headfont=\sffamily\bfseries,
	notefont=\normalfont\sffamily\scshape,
	bodyfont=\normalfont,
	headformat=swapnumber,
	headpunct=.,
	postheadspace=1em,
	spaceabove=15pt,spacebelow=10pt,
	postheadhook=\theorembookmark,
	qed=\qedsymbol]%
{mainstyleB}
\declaretheorem[name=Definition,parent=section,style=mainstyle]{definition}
\declaretheorem[name=Definition,numbered=no,style=mainstyle_unnumbered]{definition*}
\declaretheorem[name=Theorem,sharenumber=definition,style=mainstyle]{theorem}
\declaretheorem[name=Theorem,numbered=no,style=mainstyle_unnumbered]{theorem*}
\declaretheorem[name=Proposition,sharenumber=definition,style=mainstyle]{proposition}
\declaretheorem[name=Lemma,sharenumber=definition,style=mainstyle]{lemma}
\declaretheorem[name=Satz,sharenumber=definition,style=mainstyle]{satz}
\declaretheorem[name=Korollar,sharenumber=definition,style=mainstyle]{korollar}
\declaretheorem[name=Korollar,sharenumber=definition,style=mainstyleB]{korollarB}
\declaretheorem[name=Frage,numbered=no,style=mainstyle_unnumbered]{frage}
\declaretheorem[name=Erinnerung,sharenumber=definition,style=mainstyle]{erinnerungA}


% -- Beweise
\declaretheoremstyle[headfont=\bfseries\scshape,bodyfont=\normalfont,headpunct=:,postheadspace=1em,spacebelow=12pt,spaceabove=2pt,qed=\qedsymbol]{beweise}
\declaretheoremstyle[headfont=\sffamily\bfseries,bodyfont=\normalfont,headpunct=:,postheadspace=1em,spacebelow=10pt,spaceabove=10pt]{bemerkungen}
\declaretheorem[name=Beweis,numbered=no,style=beweise]{beweis}
\declaretheorem[name=Bemerkung,sharenumber=definition,style=mainstyle]{bemerkung}
\declaretheorem[name=Beispiel,sharenumber=definition,style=mainstyle]{beispiel}
\declaretheorem[name=Übung,numbered=no,style=bemerkungen]{uebung}
\declaretheorem[name=Erinnerung,numbered=no,style=bemerkungen]{erinnerung}


%--Skalarprodukt
\DeclarePairedDelimiterX\sprod[2]{\langle}{\rangle}{#1\,\delimsize\vert\,#2}
\DeclarePairedDelimiterX\skal[2]{\langle}{\rangle}{#1\,,\,#2}

%--Betrag, Gaußklammer
\DeclarePairedDelimiter{\abs}{\lvert}{\rvert}
\DeclarePairedDelimiter{\floor}{\lfloor}{\rfloor}
\DeclarePairedDelimiter{\ceil}{\lceil}{\rceil}

%--Norm
\DeclarePairedDelimiter\norm{\Vert}{\Vert}


%--Umklammern
\DeclarePairedDelimiter\enbrace{(}{)}
\DeclarePairedDelimiter\benbrace{[}{]}
\DeclarePairedDelimiter\lenbrace{<}{>}
\newcommand{\ssbrace}[1]{{\scriptscriptstyle\enbrace{#1}}}

%--Mengen
\newcommand\SetSymbol[1][]{\nonscript\:#1\vert\allowbreak\nonscript\:\mathopen{}}
\providecommand\given{} % to make it exist
\DeclarePairedDelimiterX\set[1]\{\}{\renewcommand\given{\SetSymbol[\delimsize]}#1}

%--Abbildungsdefinition
\newcommand{\mapdef}[5]{%
	\[
		\begin{array}{rcl}
			\textstyle #1 &\xrightarrow{\minwidthbox{#5}{2em}} & \textstyle #2 \\[0.5ex]
			\textstyle #3 &\xmapsto{\minwidthbox{\mbox{ }}{2em}} & \textstyle #4
		\end{array}
	\]
}

%--modifiziertes Stackrel
\newcommand{\StackText}[2]{\stackrel{\mbox{\scriptsize #1}}{#2}}

%--Befehl für eine Box mit Mindestweite
\DeclareRobustCommand{\minwidthbox}[2]{%
  \ifmmode
    \expandafter\mathmakebox
  \else
    \expandafter\makebox
  \fi
  [\ifdim#2<\width\width\else#2\fi]{#1}%
}

%--Befehl zum Einkreisen
\newcommand*\circled[1]{\tikzexternaldisable\tikz[baseline=(char.base)]{\node[shape=circle,draw,inner sep=2pt] (char) {#1};}\tikzexternalenable}


%--Differential
\newcommand{\diff}[2]{\ensuremath{\frac{{\partial #1}}{{\partial #2}} }}
\newcommand{\diffd}[2]{\ensuremath{\frac{\mathd #1}{\mathd #2} }}

%--Inhaltsverzeichnis
\usepackage[tocindentauto]{tocstyle}
\usetocstyle{KOMAlike}	
\shorthandon{"}
		
%--Punkte (als letztes laden)
\usepackage{ellipsis}

\newcommand{\fach}{Topologie \RM{2}.}
\newcommand{\semester}{WiSe 2015}
\newcommand{\homepage}{https://wwwmath.uni-muenster.de/reine/u/topos/lehre/WS2015-2016/Topologie2/}

\newcommand{\prof}{Prof.\ Dr.\ Arthur Bartels}
% \publishers{\scalebox{14}{$K$}}
\providecommand{\verfasser}{Jannes Bantje}

%--Konfiguration von scrheadings
\setheadsepline{1pt}[\color{light_gray}]
\pagestyle{scrheadings}
%\fancyhf[H, F]{}
\clearscrheadfoot

\providecommand{\shortFach}{\fach}
\lehead{ \includegraphics[height=0.5 cm,keepaspectratio]{../!config/Bilder/Logo_WWU_Muenster_light_gray.pdf}}
\rehead{\rule{0cm}{0.5cm}\footnotesize \sffamily \color{light_gray} \verfasser -- Mitschrift \shortFach}
\lohead{\rule{0cm}{0.5cm} \footnotesize \sffamily \color{light_gray} Stand: \today \; \thistime[:]}
\rohead{\includegraphics[height=0.5 cm,keepaspectratio]{../!config/Bilder/fb10logo_gray.pdf}}


\ofoot[{ \color{dark_gray} \LARGE \sffamily \thepage}]{{ \color{dark_gray} \LARGE \sffamily \thepage}} %hier wir auch der plain Stil bearbeitet!
\automark{section}
\ifoot{ \color{dark_gray} \small \leftmark}

%--Metadaten
\author{\verfasser}
\titlehead{\includegraphics[height=1.5cm, keepaspectratio]{../!config/Bilder/Logo_WWU_Muenster.pdf}%
\hfill \includegraphics[height=1.3cm, keepaspectratio]{../!config/Bilder/fb10logo.pdf}}
\title{Skript \fach}
\subtitle{Mitschrift der Vorlesung  \enquote{\fach} von \prof}

\publishers{
Aktuelle Version verfügbar bei: \bigskip\\
\normalsize
\begin{minipage}[t]{0.45\textwidth}
	\centering{
	\qrcode[height=3.5cm, version=4]{https://github.com/JaMeZ-B/latex-wwu}
	\medskip\\
	\includegraphics[height=0.4cm, keepaspectratio]{../!config/Bilder/github_octo.pdf} 
	\includegraphics[height=0.4cm, keepaspectratio]{../!config/Bilder/GitHub_Logo.pdf} (inklusive Sourcecode)\\
	{\small\url{https://github.com/JaMeZ-B/latex-wwu}}}
\end{minipage}
\hfill
\begin{minipage}[t]{0.45\textwidth}
	\centering{
	\qrcode[height=3.5cm, version=4]{btsync://B6WH2DISQ5QVYIRYIEZSF4ZR2IDVKPN3I?n=latex_share}
	\medskip\\
	\includegraphics[height=0.4cm, keepaspectratio]{../!config/Bilder/bt_sync_logo.pdf} 
	{\large \textbf{Bittorrent} Sync} \\
	\texttt{B6WH2DISQ5QVYIRYIEZSF4ZR2IDVKPN3I}
	}
\end{minipage}
}


\begin{document}
\pagenumbering{Roman}
\maketitle
\begin{abstract}
\section*{Vorwort --- Mitarbeit am Skript}
Dieses Dokument ist eine Mitschrift aus der Vorlesung \enquote{\fach, \semester}, gelesen von \prof. Der Inhalt entspricht weitestgehend dem Tafelanschrieb. Für die
Korrektheit des Inhalts übernehme ich keinerlei Garantie! Für Bemerkungen und Korrekturen -- und seien es nur Rechtschreibfehler -- bin ich sehr dankbar. 
Korrekturen lassen sich prinzipiell auf drei Wegen einreichen: 
\begin{itemize}
	\item Persönliches Ansprechen in der Uni, Mails an \hrefsymmail{mailto:\mail}{\mail} (gerne auch mit annotieren PDFs) 
	\item \emph{Direktes} Mitarbeiten am Skript: Den Quellcode poste ich auf GitHub (siehe oben), also stehen vielfältige Möglichkeiten der Zusammenarbeit zur Verfügung:
	Zum Beispiel durch Kommentare am Code über die Website und die Kombination Fork + Pull Request. Wer sich verdient macht oder ein Skript zu einer Vorlesung, die 
	ich nicht besuche, beisteuern will, dem gewähre ich gerne auch Schreibzugriff.
	
	Beachten sollte man dabei, dass dazu ein Account bei \url{github.com} notwendig ist, der allerdings ohne Angabe von persönlichen Daten angelegt werden kann. 
	Wer bei GitHub (bzw. dem zugrunde liegenden Open-Source-Programm \enquote{\texttt{git}}) -- verständlicherweise -- Hilfe beim Einstieg braucht, dem helfe ich gerne 
	weiter. Es gibt aber auch zahlreiche empfehlenswerte Tutorials im Internet.\footnote{zB. \url{https://try.github.io/levels/1/challenges/1}, ist auf Englisch, aber dafür 
	interaktives LearningByDoing}
	\item \emph{Indirektes} Mitarbeiten: \TeX-Dateien per Mail verschicken. 
	
	Dies ist nur dann sinnvoll, wenn man einen ganzen Abschnitt ändern möchte (zB. einen alternativen Beweis geben), da ich die Änderungen dann per Hand einbauen muss! Ich freue mich aber auch über solche Beiträge!
\end{itemize}
\section*{Vorlesungshomepage}
{\centering 
\begin{minipage}[c][][c]{\textwidth}
	\centering \qrcode[height=3cm]{\homepage} \medskip\\
	\footnotesize \url{\homepage}
\end{minipage}
\par}

\end{abstract}

\tableofcontents
\cleardoubleoddemptypage

\pagenumbering{arabic}
\setcounter{page}{1}
\setcounter{footnote}{0}

\section{Kohomologie} % (fold)
\label{sec:1}

\begin{definition}[{name=[Kokettenkomplex und Kohomologie]}]
	Sei $R$ ein Ring. Ein \bet{$R$-Kokettenkomplex}\index{Kokettenkomplex} $(C^*,d^*)$ ist eine Folge von $R$-Moduln $(C^n)_{n}$ zusammen mit $R$-linearen Abbildungen 
	$d^n \colon C^n \to C^{n+1}$, sodass $d^{n+1} \circ d^n =0$ gilt. Der \bet{$n$-te Kohomologiemodul}\index{Kohomologiemodul} von $(C^*,d^*)$ ist definiert als 
	\[
		H^n(C^*,d^*) = \frac{\ker d^n \colon C^n \to C^{n+1}}{\im d^{n-1} \colon C^{n-1} \to C^n} 
	\]
	Sei $(D^*,d^{*})$ ein weiterer Kokettenkomplex.
	Eine \Index{Kokettenabbildung} ist eine Folge von $R$-linearen Abbildungen $f^n\colon C^n\to D^n$, sodass $ d^n_D \circ f^n = f^{n+1}\circ d^n_C$ für alle $n$ gilt. Ist 
	auch $g^n \colon C^n \to D^n$ eine Kokettenabbildung, so nennen wir eine $R$-lineare Abbildung $h^n \colon C^n \to D^{n-1}$ mit
	\[
		f^n - g^n  = h^{n+1} \circ d_C^n + d_D^{n-1} \circ h^n 
	\]
	eine \Index{Kokettenhomotopie} zwischen $f$ und $g$. Zu jeder Kokettenabbildung $f^n\colon C^n\to D^n$ gibt es eine 
	\bet{induzierte Abbildung}\index{induzierte Abbildung!Kohomologie} auf Kohomologie genau wie bei Homologie.
	
\end{definition}

\begin{bemerkung}[{name=[Dualer Kokettenkomplex und Umwandlung von Kettenkomplexe in Kokettenkomplexe]}] \leavevmode
	\begin{enumerate}[i)]
		\item Sei $(C_*,d_*)$ ein $R$-Kettenkomplex und $V$ ein $R$-Modul. Dann erhalten wir einen $R$"=Kokettenkomplex $(C^*,d^*)$ durch 
		\[
			C^n \coloneqq \Hom_R(C_n,V) 
		\]
		und $d^n \colon C^n \to C^{n+1}$ definiert durch $\alpha \mapsto \alpha \circ d_{n+1}$. Dieser Kokettenkomplex $(C^*,d^*)$ heißt der 
		\bet{$V$-duale}\index{Kokettenkomplex!$V$-dualer Kokettenkomplex} $R$-Kokettenkomplex zu $(C_*,d_*)$. Ist $V=R$, so nennen wir ihn auch kürzer 
		\Index{dualer Kokettenkomplex} 
		\item Benutzen wir $\mathbb{Z}$ statt $\mathbb{N}$ als Indexmenge, so können wir durch $(C^n,d^n) \leadsto (C_n \coloneqq C^{-n}, d_n \coloneqq d^{-n})$ jeden 
		Kokettenkomplex einem Kettenkomplex zuordnen. Dieser Prozess ist offensichtlich umkehrbar.
	\end{enumerate}
\end{bemerkung}

\begin{beispiel}[{name=[Kohomologie ist nicht Dualisieren von Homologie]}]
	Es sei $(C_*,d_*)= \big(\begin{tikzcd}[cramped] \mathbb{Z} & \mathbb{Z} \lar["\cdot 2","d_1"'] & 0 \lar & \ldots \lar \end{tikzcd}\big)$ ein Kettenkomplex. 
	Dann ist
	\[
		H_k(C_*,d_*) = \begin{cases}
			\nicefrac{\mathbb{Z}}{2 \mathbb{Z}}, &\text{ falls } k =0\\
			0 , & \text{ sonst} 
		\end{cases}
	\]
	Der $\mathbb{Z}$-duale Kokettenkomplex hat dann folgende Gestalt
	\[
		\begin{tikzcd}
			 \Hom_{\mathbb{Z}}(\mathbb{Z},\mathbb{Z}) \rar["d^0"] \dar["\cong" description] & \Hom_{\mathbb{Z}}(\mathbb{Z},\mathbb{Z}) \rar \dar["\cong" description] 
			& 0 \rar & \ldots \\
			\mathbb{Z} \rar["\cdot 2"] & \mathbb{Z}
		\end{tikzcd}
	\]
	Damit ist die Kohomologie $H^k(C^*,d^*)=0$ für $k \neq 1$ und isomorph zu $\sfrac{\mathbb{Z}}{2 \mathbb{Z}}$ für $k=1$.
	Es gilt also nicht immer $H^*(\Hom(C_*;R),d^*) =  \Hom \enbrace{H_*(C_*,d_*),R}$. 
\end{beispiel}

\begin{definition}[{name=[singulärer Kokettenkomplex]}]
	Sei $(X,A)$ ein Paar von topologischen Räumen und $V$ eine abelsche Gruppe. Der \Index{singuläre Kokettenkomplex} von $(X,A)$ mit Koeffizienten in $V$ ist definiert durch
	\[
		C_\sing^*(X,A;V) \coloneqq \Hom_\mathbb{Z} \enbrace[\big]{C_*^\sing(X,A),V}
	\]
	und $d^*_\sing(\alpha) \coloneqq - (-1)^{\abs{\alpha}} \cdot \alpha \circ d_{*+1}^\sing$. Dabei ist $\abs{\alpha}=n$ für $\alpha \in C^n_\sing(X,A;R)$. Die 
	Kohomologie von $\enbrace[\big]{C^*_\sing(X,A;V), d^*_\sing}$ heißt die \Index{singuläre Kohomologie} von $(X,A)$ mit Koeffizienten in $R$.
\end{definition}

\begin{bemerkung}[{name=[Isomorphie des singulären Kettenkomplexes]}]
	Sei $R$ ein kommutativer Ring und $V$ ein $R$-Modul. Dann ist der singuläre Kokettenkomplex
	\(
		\enbrace*{C^*_\sing(X,A;V), d^*_\sing}
	\)
	isomorph zum $V$-dualen $R$-Kokettenkomplex des $R$-Kettenkomplexes $\enbrace*{C^\sing_*(X,A;R), d^\sing_*}$, da $C^\sing_*(X,A;R) \cong R \otimes C^\sing_*(X,A)$ und
	$\hom_\mathbb{Z}(W,V) \cong \hom_R(R \otimes W,V)$ für jeden $\mathbb{Z}$-Modul $W$.
\end{bemerkung}


\begin{definition}[{name=[induzierte Abbildung auf Kokettenkomplexen]}]
	Sei $f \colon (X,A) \to (Y,B)$ eine stetige Abbildung von Paaren. Dann erhalten wir eine Kokettenabbildung $f^* \colon C^*_\sing(Y,B;V) \to C^*_\sing(X,A;V)$ durch
	\index{induzierte Abbildung!singuläre Kokettenkomplexe}
	\[
		f^*(\alpha) \coloneqq \alpha \circ f_*
	\]
\end{definition}

\begin{bemerkung}[{name=[Induzierte Abbildung einer Komposition]}]
	Ist $g \colon (Y,B) \to (Z,C)$ eine weitere Abbildung von Paaren, so gilt $(g \circ f)^* = f^* \circ g^*$.
\end{bemerkung}

\begin{definition}[{name=[kontravarianter Funktor]}]
	Seien $\mathcal{C}$ und $\mathcal{D}$ Kategorien. Ein \Index{kontravarianter Funktor}  $F \colon \mathcal{C} \to \mathcal{D}$ ordnet jedem Objekt $C$ in $\mathcal{C}$ ein Objekt
	$D$ in $\mathcal{D}$ zu und jedem Morphismus $f \colon C \to C'$ einem Morphismus $F(f) \colon F(C') \to F(C)$ in $\mathcal{D}$ zu. Dabei muss gelten:
	\begin{enumerate}[i),noitemsep]
		\item $F(\id_C)= \id_{F(C)}$
		\item Für $\begin{tikzcd}[cramped] C \rar["f"] & C' \rar["f'"] & C'' \end{tikzcd}$ gilt $F(f' \circ f) = F(f) \circ F(f')$. 
	\end{enumerate}
	Kürzer ist ein \emph{kontravarianter} Funktor $F \colon \mathcal{C} \to \mathcal{D}$ das selbe wie ein \emph{kovarianter} Funktor $\mathcal{C}^\op \to \mathcal{D}$.
\end{definition}

\begin{beispiel}[{name=[kontravariante Funktoren]}] Es gibt zahlreiche Beispiele für kontravariante Funktoren:
	\begin{enumerate}[i),itemsep=0pt]
		\item $\id \colon \mathcal{C} \to \mathcal{C}^\op$ ist kontravariant.
		\item Sei $V$ eine abelsche Gruppe. Wir erhalten einen kontravarianten Funktor
		\[
			\Hom(-,V) \colon \mathbb{Z}\text{-}\MOD  \longrightarrow \mathbb{Z}\text{-}\MOD
		\]
		\item Sei $V$ eine abelsche Gruppe. Dann sind 
		\begin{align}
			C_\sing^*(-,V) &\colon \TOP^2 \longrightarrow \mathbb{Z}\text{-}\textsc{Kokettenkomplexe}  \\
			H^*_\sing \enbrace[\big]{C_\sing^*(-,V),d_\sing^*} &\colon \TOP^2 \longrightarrow \textsc{Gr}\text{-}\mathbb{Z}\text{-}\MOD   
		\end{align}
		kontravariante Funktoren.
	\end{enumerate}
\end{beispiel}

\begin{satz}[{name=[Eigenschaften singulärer Kohomologie]}]
	Singuläre Kohomologie hat die folgenden Eigenschaften:
	\begin{enumerate}[i),itemsep=1.5pt]
		\item \textbf{Dimensionsaxiom:} Es gilt $H^n_\sing(\set{\pt};V) = V$, falls $n=0$ ist und sonst $0$.\index{Dimensionsaxiom} 
		\item \textbf{Paarfolge:} Es gibt eine natürliche Transformation $\partial^*\colon H^*(A;V) \to H^{*+1}(X,A;V)$ sodass für jedes Paar\index{Paarfolge} 
		\[
			\begin{tikzcd}
				0 \rar & H^0(X,A;V) \rar & H^0(X;V) \rar & H^0(A;V) \rar["\partial"] & H^1(X,A;V) \rar & \ldots 
			\end{tikzcd}
		\]
		eine lange exakte Folge ist. $\partial$ bezeichnet man auch als \Index{verbindende Abbildung}.
		\item \textbf{Ausschneidung:} Sei $L \subseteq A$ mit $\overline{L} \subseteq \mathring{A}$. Dann induziert die Inklusion \index{Ausschneidung}
		$i \colon (X \setminus L, A \setminus L) \hookrightarrow (X,A)$ einen Isomorphismus $i^* \colon H^*(X,A;V) \to H^*(X \setminus L, A \setminus L;V)$.
		\item \textbf{Homotopieinvarianz:} Sind $f,g \colon (X,A) \to (Y,B)$ homotope Abbildungen von Paaren, so gilt $f^* = g^*$ für die induzierten Abbildungen in singulärer 
		Kohomologie.\index{Homotopieinvarianz}
	\end{enumerate}
\end{satz}
\begin{beweis}
	Für singuläre Homologie haben wir die entsprechenden Aussagen schon bewiesen. In allen vier Fällen folgt die Aussage für Kohomologie aus schon bewiesenen Aussagen über 
	den singulären Kettenkomplex.
	
	Wir führen dies an dieser Stelle nur für iv) aus, die anderen Punkte wurden teilweise in den Übungen bewiesen. Seien $f, g \colon X \to Y$ homotop. Dann gibt es eine 
	Kettenhomotopie $H \colon C^\sing_*(X) \to C^\sing_{*+1}(Y)$ zwischen den auf dem singulären Kettenkomplex induzierten Abbildungen $f_*$ und $g_*$. Es gilt also
	\[
		d_{n+1} \circ H + H \circ d_n = f_* - g_*
	\]
	$H$ induziert $H^\# \colon C^*_\sing(Y;V) \to C_\sing^{*-1}(X;V)$ mit $H^\#(\alpha) \coloneqq (-1)^{\abs{\alpha}} \cdot \alpha \circ H$. Es gilt nun für $\alpha \in C_\sing^n(Y;V)$
	\begin{align}
		\enbrace*{d^{n-1} \circ H^\# + H^\# \circ d^n}(\alpha) &= d^{n-1} \circ H^\#(\alpha) + H^\# \circ d^n (\alpha)\\
		&= d^{n-1} \enbrace[\big]{(-1)^n \cdot (\alpha \circ H)} - (-1)^n \cdot H^\# \enbrace*{\alpha \circ d_{n+1}} \\
		&= (-1)^n \cdot \enbrace*{(-1)^{n} \alpha \circ H  \circ d_n - (-1)^{n+1} \alpha \circ d_{n+1} \circ H} \\
		&= \alpha \circ H \circ d_n + \alpha \circ d_{n+1} \circ H \\
		&= \alpha \enbrace*{f_* - g_*} = f^*(\alpha) - g^*(\alpha) 
	\end{align}
	Damit ist $f^* - g^* =0$ in Kohomologie, da die linke Seite für $\alpha \in \ker d^n$ im Bild von $d^{n-1}$ liegt.
\end{beweis}

\begin{bemerkung}[{name=[Angabe des Verbindungshomomorphismus]}]
	Sei $(X,A)$ ein Paar von topologischen Räumen. Der Verbindungshomomorphismus $\partial \colon H^n(A;V) \to H^{n+1}(X,A;V)$ kann wie folgt beschrieben werden: Sei 
	$\alpha \colon C_n(A) \to V$ ein Kozykel. Setze $\alpha$  durch $\hat{\alpha}(\sigma)=0$ für $\sigma \colon \abs{\Delta^n} \to X$, $\im \sigma \nsubseteq A$ zu 
	$\hat{\alpha} \colon C_n(X) \to V$ fort. Dann ist 
	\[
		\partial \benbrace{\alpha} = \benbrace*{d^n \hat{\alpha}} = \benbrace*{\hat{\alpha} \circ d_{n+1}} \in H^{n+1}(X,A;V)
	\]
\end{bemerkung}

\begin{beispiel}[{name=[$0$-te Kohomologie]},label=0teKoho]
	Die Gruppe $H^0(X;V)$ ist die Gruppe aller Abbildungen $\xi \colon X \to V$, die konstant auf Wegzusammenhangskomponenten sind. Die Gruppe $H^0(X,A;V)$ besteht aus allen 
	solchen Abbildungen, die zusätzlich auf $A$ trivial sind.
\end{beispiel}

\begin{definition}[{name=[Produkt von $R$-Moduln]}]
	Seien $(V_i)_{i \in I}$ $R$-Moduln. Mit $V \coloneqq \prod_{i \in I} V_i$ bezeichnen wir das \Index{Produkt} der $V_i$. Element in $V$ sind $I$-Folgen $(v_i)_{i \in I}$ mit 
	$v_i \in V_i$. Die $R$-Modulstruktur ist erklärt durch
	\begin{align}
		(v_i)_{i \in I}+ (w_i)_{i \in I} &\coloneqq (v_i + w_i)_{i \in I} \\
		r \cdot (v_i)_{i \in I} & \coloneqq (r \cdot v_i)_{i \in I}
	\end{align}
	Für jedes $i_0 \in I$ erhalten wir eine $R$-lineare Abbildung $\pi_{i_0} \colon V \to V_{i_0}$, $(v_i)_{i \in I} \mapsto v_{i_0}$
\end{definition}

\begin{bemerkung}[name={Universelle Eigenschaft des Produktes}]
	Seien $V_i$ für $i \in I$ $R$-Moduln. Sei $W$ ein weiterer $R$-Modul. Dann gibt es zu jeder Folge $(f_i\colon W \to V_i)_{i \in I}$ von $R$-linearen Abbildungen eine eindeutige 
	$R$-lineare Abbildung $f \colon W \to \prod_{i \in I}V_i$ mit $f_i= \pi_i \circ f$. Diese ist gegeben durch
	$f(w) \coloneqq \enbrace*{f_i(w)}_{i \in I}$.
\end{bemerkung}

\begin{bemerkung}[{name=[endliche Produkte]}]
	Ist $I$ endlich, so gilt
	\(
		\bigoplus_{i \in I} V_i = \prod_{i \in I} V_i.
	\)
\end{bemerkung}

\begin{bemerkung}[{name=[{Eigenschaften von $\Hom(-,W)$ mit Produkten}]}]
	Es seien $V_i$ für $i \in I$ $R$-Moduln und $W$ ein weiterer $R$-Modul. Seien $j_{i_0} \colon V_{i_0} \to\bigoplus_{i \in I}V_i$ die Inklusionen $v_{i_0}\mapsto (v_i)_{i \in I}$
	mit $v_i =v_{i_0}$ für $i=i_0$ und $0$ sonst. Dann erhalten wir einen Isomorphismus 
	\mapdef{\Hom_R \enbrace[\big]{\bigoplus_i V_i,W}}{\prod_{i \in I} \Hom_R(V_i,W)}{f}{(f \circ j_i)_{i \in I}}{\cong}
\end{bemerkung}

\begin{satz}[{name=[Kohomologie und Summe von topologischen Räumen]}]
	Sei $X= \coprod_{i \in I} X_i$ die Summe von topologischen Räumen $X_i$. Dann induzieren die Inklusionen $j_i \colon X_i \to X$ einen Isomorphismus 
	\mapdef{H^* \enbrace*{X,V}}{\prod_{i \in I} H^*(X_i,V)}{\xi}{\enbrace[\big]{(j_i)^*(\xi)}_{i \in I}}{\cong}
	% \[
	% 	\begin{tikzcd}[row sep=0em]
	% 		H^* \enbrace*{X,V} \rar & \prod_{i \in I} H^*(X_i,V) \\
	% 		\xi \rar[mapsto] & \enbrace[\big]{(j_i)^*(\xi)}_{i \in I}
	% 	\end{tikzcd}
	%
	% \]
\end{satz}
\begin{beweis}
	Die $(j_i)_{i \in I}$ induzieren einen Isomorphismus von Kettenkomplexen $\varphi \colon \bigoplus_{i \in I}C_*(X_i) \to C_*(X)$ mittels 
	$(a_i)_{i \in I} \mapsto \sum_{i \in I} (j_i)_* (a_i)$. Wegen 
	\[
		\Hom \enbrace*{\bigoplus\nolimits_{i \in I} C_*(X_i),V} \cong \prod_{i \in I} \Hom \enbrace*{C_*(X_i),V}
	\]
	erhalten wir einen Isomorphismus von Kokettenkomplexen
	\mapdef{C^* (X;V)}{\prod_{i \in I} C^*(X_i;V)}{\alpha}{\enbrace[\big]{j_i^*(\alpha)}_{i \in I}}{\cong}
	Dieser induziert den behaupteten Isomorphismus in Kohomologie.  
\end{beweis}

\begin{definition}[{name=[reduzierte Kohomologie]}]
	Die reduzierte Kohomologie von $X$, $\tilde{H}^*(X;V)$ ist definiert als der Kokern von $p^* \colon H^* \enbrace*{\set{\pt};V} \to H^*(X;V)$, also als der Quotient
	$\sfrac{H^*(X;V)}{\im p^*}$.
\end{definition}

\begin{bemerkung}[{name=[Zusammenhang zwischen reduzierter und regulärer Homologie]}]
	Für reduzierte Kohomologie gilt analog zu reduzierter Homologie
	\[
		H^n(X;V) \cong \begin{cases}
			\tilde{H}^n(X;V), &\text{ falls } n \not= 0\\
			\tilde{H}^0(X;V) \oplus V , &\text{ falls } n=0 
		\end{cases}
	\]
\end{bemerkung}

\begin{beispiel}[{name=[Kohomologie der Sphäre]}]
	Viele Berechnungen für Homologiegruppen lassen sich problemlos auf Kohomologie übertragen. Zum Beispiel ist die reduzierte Kohomologie der Sphäre gegeben durch
	\[
		\tilde{H}^l (S^n;V) \cong H^l \enbrace*{D^n,S^{n-1};V} \cong \begin{cases}
			V, &\text{ falls } l=n\\
			0, &\text{ sonst}  
		\end{cases}
	\]
\end{beispiel}
% section 1 (end)
\newpage

\section{Die Paarung zwischen Kohomologie und Homologie} % (fold)
\label{sec:2}

\begin{definition}[{name=[Paarung]}]
	Sei $V$ ein $\mathbb{Z}$-Modul, $(X,A)$ ein Paar von topologischen Räumen. Wir definieren die \Index{Paarung} zwischen Kohomologie und Homologie 
	\[
		H^n(X,A;V) \times H_n(X,A) \longrightarrow V \quad , \quad (\xi,x) \longmapsto \xi(x)
	\]
	wie folgt: Wähle $\alpha \in C^n_\sing(X,A;V)$ mit $\benbrace*{\alpha}=\xi$ und $a \in C^\sing_n(X,A)$ mit $\benbrace*{a}=x$. Dann setze $\xi(x) \coloneqq \alpha(a)$.
\end{definition}

\begin{bemerkung}[{name=[Grundlegende Eigenschaften der Paarung]}]
	\leavevmode
	\begin{itemize}[itemsep=0pt]
		\item Sei $\beta \in C^{n-1}_\sing(X,A;V)$ und $b \in C^\sing_{n+1}(X,A)$. Für $\alpha \in \ker d^n$, $a \in\ker d_n$ folgt 
		\[
			\enbrace*{\alpha + d^{n-1}(\beta)}\enbrace[\big]{a + d_{n+1}(b)} = \alpha(a) \pm \Underbrace{\alpha \enbrace[\big]{d_{n+1}(b)}}{=0 } 
			+ \Underbrace{d^{n-1}(\beta) \enbrace[\big]{a + d_{n+1}(b)}}{=0}
		\]
		Damit folgt, dass die Paarung wohldefiniert ist.
		\item Für $r \in \mathbb{Z}$, $x,x' \in H_n(X,A)$, $\xi,\xi' \in H^n(X,A;V)$ gelten folgende Regeln
		\[
			(r \cdot \xi)(x) = r \cdot \xi(x) = \xi(r \cdot x) \qquad (\xi +\xi')(x) = \xi(x) + \xi'(x) \qquad \xi(x+x')= \xi(x) + \xi(x')
		\]
		\item 	Wir können $(\xi,x)\mapsto \xi(x)$ auch als interpretieren als Homomorphismus
		\mapdef{f \colon H^n(X,A;V)}{\Hom_\mathbb{Z} \enbrace[\big]{H_n(X,A),V}}{\xi}{\enbrace[\big]{x \mapsto \xi(x)}}{}
		\end{itemize}
\end{bemerkung}


\begin{satz}[label=eig_kohomo_to_hom_homo,{name=[Paarung von Kohomologie in Homomorphismen über Homologie ist surjektiv]}]
	Für die eben definierte Abbildung $f \colon H^n(X,A;V) \to \Hom_\mathbb{Z} \enbrace[\big]{H_n(X,A),V}$ gilt
	\begin{enumerate}[(i),itemsep=0pt]
		\item $f$ ist surjektiv
		\item Ist $V$ ein $\mathbb{Q}$-Vektorraum, so ist $f$ auch injektiv.
	\end{enumerate}
\end{satz}

\noindent Um den Satz beweisen zu können, benötigen wir zunächst zwei technische Aussagen:

\begin{lemma}[label=lem:untergrp_frei,{name=[Untergruppen freier abelscher Gruppen sind frei]}]
	Untergruppen freier abelscher Gruppen sind frei.
	% Sei $C$ eine Untergruppe der freien abelschen Gruppe $\mathbb{Z}[S]$. Dann gibt es $T \subseteq S$ und eine $T$-Folge $(n_t)_{t \in T}$ von positiven Zahlen $n_t$, so dass $C$
	% das Bild der injektiven Abbildung $M \colon \mathbb{Z}[T] \to \mathbb{Z}[S]$, $\sum z_t t \mapsto \sum z_t m_t \cdot t$ ist. Insbesondere ist $C$ frei.
\end{lemma}
\begin{beweis}
	Sei $C$ eine Untergruppe von $\mathbb{Z}[S]$. Sei $\mathcal{M}$ die Menge der Tripel $(T,R,\varphi)$ mit
	\begin{itemize}[itemsep=0pt,topsep=1ex]
		\item $T \subseteq R \subseteq S$.
		\item $\varphi \colon \mathbb{Z}[T] \to \mathbb{Z}[R] \cap C$ ein Isomorphismus.
	\end{itemize}
	Wir definieren eine partielle Ordnung auf $\mathcal{M}$ durch
	\[
		(T,R,\varphi) \le (T',R',\varphi') :\iff T \subseteq T', R \subseteq R', \varphi'|_{T}=\varphi
	\]
	$\mathcal{M}$ ist nicht leer ($T=R=\emptyset$) und jede aufsteigende Kette besitzt eine obere Schranke (Vereinigung der Kette) in $\mathcal{M}$.
	Das Lemma von Zorn liefert uns nun die Existenz eines maximalen Elements $(T,R,\varphi) \in \mathcal{M}$. Zu zeigen bleibt $R=S$. 
	
	Angenommen es existiert ein $s \in S \setminus R$.
	Ist $C \cap \mathbb{Z} \benbrace[\big]{R \cup \set{s}} = C \cap \mathbb{Z}[R]$, so ist $(T,R,\varphi) \lneqq (T,R \cup \set{s},\varphi) \in \mathcal{M}$ im Widerspruch zur 
	Maximalität von $(T,R,\varphi)$. Sei also $\mathbb{Z}[R] \cap C \subsetneq \mathbb{Z}\benbrace[\big]{R \cup \set{s}} \cap C$. Betrachte nun das Diagramm in Abbildung 
	\ref{diag:frei_untergrp}.
	\begin{figure}[t]
		\centering{
		\begin{tikzcd}[row sep=large]
			\mathbb{Z}[T] \dar["\varphi","\cong"'] \rar[hook] & \mathbb{Z} \benbrace[\big]{T \cup \set*{s}} \dar["\varphi^+"] \rar[two heads] 
			& \mathbb{Z}\benbrace[\big]{\set*{s}} {\color{dark_gray} \ni s} \dar["\cong"] 
			\ar[dd,dark_gray,controls={+(-17:1.5) and +(20:1.5)},mapsto,start anchor=east, end anchor=east] \\
			\mathbb{Z}[R] \cap C \rar[hook,"j"] \dar[hook,"i_0"]& \mathbb{Z}\benbrace[\big]{R \cup \set{s}} \cap C \rar[two heads,"p"] \dar[hook,"i_1"] 
			& \faktor{\mathbb{Z}\benbrace[\big]{R \cup \set{s}} \cap C}{\mathbb{Z}[R] \cap C} \dar["i_2"] \\
			\mathbb{Z}[R] \rar[hook,"j'"] & \mathbb{Z}\benbrace[\big]{R \cup\set{s}} \rar[two heads,"p'"] & \mathbb{Z}\benbrace[\big]{\set{s}}{\color{dark_gray} \ni m \cdot s} 
		\end{tikzcd}
		\caption{Diagramm zum Beweis von \autoref{lem:untergrp_frei} \label{diag:frei_untergrp}}
		}
	\end{figure}
	Es ist $\mathbb{Z}[R] \cap C  = \mathbb{Z}[R] \cap \enbrace[\big]{\mathbb{Z}[R \cup \set{s}] \cap C}$ und damit muss auch der von $i_0$ und $i_1$ induzierte Homomorphismus 
	$i_2 \colon \sfrac{\mathbb{Z}[R \cup \set{s}] \cap C}{\mathbb{Z}[R] \cap C} \to \mathbb{Z}[\set{s}]$ injektiv sein. Es folgt, dass 
	$\im i_2 = m \cdot \mathbb{Z}\benbrace*{\set{s}}$
	ist für ein $m >0$. Sei $c \in \mathbb{Z}[R \cup \set{s}] \cap C$ ein Urbild von $m \cdot s$ unter $p \circ i_1$. Nun können wir $\varphi$ durch $s \mapsto c$ zu
	$\varphi^+ \colon \mathbb{Z}\benbrace*{T \cup\set{s}} \to \mathbb{Z}\benbrace*{R \cup \set{s}}\cap C$ fortsetzen. Es folgt, dass $\varphi^+$ ein Isomorphismus ist im Widerspruch zur 
	Maximalität von $(T,R,\varphi)$.
	% Sei $\mathcal{M}$ die Menge aller Tripel $(T,R,m)$ mit
	% \begin{itemize}
	% 	\item $T \subseteq R \subseteq S$
	% 	\item $m=(m_t)_{t \in T}$ ist eine $T$-Folge von positiven Zahlen
	% 	\item das Bild von $M \colon \mathbb{Z}[T] \to \mathbb{Z}[S]$, $t \mapsto m_t \cdot t$ ist $C \cap \mathbb{Z}[R]$.
	% \end{itemize}
	% Wir definieren eine partielle Ordnung auf $\mathcal{M}$ durch
	% \[
	% 	(T,R,m) \le (T',R',m') :\iff T \subseteq T', R \subseteq R', m'|_{T}=m
	% \]
	% Es ist $\mathcal{M} \not= \emptyset$, da $(\emptyset,\emptyset, \cdot ) \in \mathcal{M}$. Jede aufsteigende Ketten in $\mathcal{M}$ hat eine obere Schranke. Nach Zorns Lemma
	% existiert ein maximales Element $(T,R,m) \in \mathcal{M}$. Es bleibt zu zeigen: $R=S$.
\end{beweis}

\begin{lemma}[label=lem:abel_qvek,{name=[{Fortsetzungen für rationale Vektorräume}]}]
	Sei $A_0$ eine Untergruppe der abelschen Gruppe $A$, $V$ ein $\mathbb{Q}$-Vektorraum und $\beta_0 \colon A_0 \to V$ ein Gruppenhomomorphismus. Dann gibt es eine Fortsetzung 
	$\beta \colon A \to V$ von $\beta_0$ zu einem Gruppenhomomorphismus.  
\end{lemma}
\begin{beweis}
	Die Inklusion $i \colon A_0 \to A$ induziert nach Aufgabe 1 vom ersten Übungsblatt einen injektiven $\mathbb{Q}$-Vektorraum-Homomorphismus 
	\[
		\mathbb{Q}\otimes i \colon \mathbb{Q} \otimes A_0 \longrightarrow \mathbb{Q} \otimes A , \qquad q \otimes a_0 \mapsto q \otimes i(a_0)
	\]
	Nun können wir die $\mathbb{Q}$-lineare Abbildung $q \otimes a_0 \mapsto q \cdot \beta_0(a_0) \in V$ von $\mathbb{Q} \otimes A_0$ zu 
	$\overline{\beta} \colon \mathbb{Q} \otimes A \to V$ fortsetzen. Dann ist $a \mapsto \overline{\beta}(1 \otimes a)$ die gesuchte Fortsetzung von $\beta_0$.
\end{beweis}

\begin{beweis}[name={von \autoref{eig_kohomo_to_hom_homo}}] \leavevmode
	\begin{enumerate}[(i)]
		\item Sei $\varphi \colon H_n(X,A) \to V$ gegeben. Sei $p \colon\ker d_n \twoheadrightarrow H_n(X,A)$ die Projektion. Betrachte die kurze exakte Folge 
	\[
		\begin{tikzcd}
			\ker  d_n \rar[hook,"i"] & C_n^\sing(X,A) \rar["d_n",two heads] & \im d_n
		\end{tikzcd}
	\]
	Als Untermodul des freien Moduls $C_{n-1}^\sing(X,A)$ ist $\im d_n$ nach \autoref{lem:untergrp_frei} frei, insbesondere spaltet die kurze exakte Sequenz und es gilt
	$C^\sing_n(X,A) \cong \ker d_n \oplus \im d_n$. Daher können wir $\varphi \circ p \colon \ker d_n \to V$ zu $\alpha \colon C_n^\sing(X,A) \to V$ fortsetzen. Genauer: Sei 
	$s \colon \im d_n \to C_n^\sing(X,A)$ ein Spalt. Dann können wir $\alpha \colon C_n^\sing(X,A) \to V$ definieren durch 
	\(
		\alpha(a) \coloneqq \varphi \circ p \enbrace*{a - s \enbrace[\big]{d_n(a)}}
	\).
	Es folgt 
	\begin{align}
		d^n (\alpha)(a)= \alpha \circ d_{n+1}(a)= \varphi \circ p \enbrace*{d_{n+1}(a) - d_{n+1} \enbrace[\big]{s \enbrace*{d_n a}}} =0
	\end{align}
	und damit $[\alpha] \in H^n(X,A;V)$ sowie $f \enbrace*{[\alpha]}=\varphi$, womit die Surjektivität gezeigt ist.
	\item Sei $\alpha \in C^n_\sing(X,A;V)$ mit $d^n(\alpha)=0$. Sei $\benbrace*{\alpha}\in\ker f$, also $\alpha(a)=0$ für alle $a \in \ker d_n$. Dann faktorisiert $\alpha$ über
	$\im d_n \subseteq C_{n-1}(X,A)$, denn $\alpha$ liegt im Kern von $d^n$ und induziert $\beta_0 \colon \im d_n \to V$ mit $\alpha = \beta_0 \circ d_n$ durch Einschränkung 
	auf den Summanden $\im d_n$. 
	
	Ist $V$ ein $\mathbb{Q}$-Vektorraum, so können wir $\beta_0$ nach \autoref{lem:abel_qvek} zu $\beta \colon C_{n-1}(X,A) \to V$ fortsetzen. Es folgt 
	$\alpha= \beta \circ d_n = \pm \enbrace*{d^{n-1} \beta}=0$ und damit $\benbrace*{\alpha}=0$ in $H^n(X,A;V)$. \qedhere
	\end{enumerate}
\end{beweis}

\begin{korollarB}[{name=[{Kohomologie mit rationalen Vektorräumen}]}]
	Für den Sonderfall der $\mathbb{Q}$-Vektorräume gilt $H^n(X,A;\mathbb{Q}) \cong \Hom_{\mathbb{Z}}\enbrace[\big]{H_n(X,A);\mathbb{Q}}$.
\end{korollarB}

\begin{bemerkung}
	Es gilt sogar $H^n \enbrace*{X,A;\mathbb{Q}} \cong \Hom_\mathbb{Q} \enbrace[\big]{H_n(X,A;\mathbb{Q}),\mathbb{Q}}$.
\end{bemerkung}
\newpage
% section 2 (end)

\section{Produkte auf Kohomologie} % (fold)
\label{sec:3}

\begin{definition}[{name=[Einschränkung von singulären Simplizes]}]
	Sei $\sigma \colon \abs{\Delta^n}  \to X$ ein singulärer Simplex in $X$. Für $0 \le p \le n$ definieren wir 
	\[
		\sigma|_{\benbrace*{0,\ldots ,p}} \colon \abs{\Delta^p} \longrightarrow X \quad , \qquad \sigma|_{\benbrace{p,\ldots ,n}} \colon \abs{\Delta^{n-p}} \longrightarrow X  
	\]
	durch $\sigma|_{\benbrace*{0,\ldots ,p}}(t_0, \ldots ,t_p) = \sigma \enbrace*{t_0, \ldots , t_p, 0, \ldots ,0}$ und $\sigma|_{p,\ldots ,n}(t_{p}, \ldots ,t_n)= \sigma(0,\ldots ,0,t_{p}, \ldots , t_n)$.
\end{definition}

\begin{figure}[h]
	\centering{
	\begin{tikzpicture}[scale=1.4]
		\draw[fill=DodgerBlue3,fill opacity=0.4] (0,0) -- ++(1,0) -- ++ (120:1) -- cycle;
		\draw[thick] (0,0) node[below left]{$0$} -- ++(1,0) node[below right]{$2$} -- ++ (120:1) node[above]{$1$} -- cycle;
		\draw[thick] (3,0) node[below left]{$0$} -- ++(1,0) node[below right]{$2$} -- ++ (120:1) node[above]{$1$} -- cycle;
		\draw[fill=Firebrick3,fill opacity=0.4] (6,0) -- ++(1,0) -- ++ (120:1) -- cycle;
		\draw[thick] (6,0) node[below left]{$0$} -- ++(1,0) node[below right]{$2$} -- ++ (120:1) node[above]{$1$} -- cycle;
		
		% linker Simplex
		\draw[DodgerBlue3,dashed,ultra thick] (0,0) -- ++(1,0) -- ++ (120:1) node[midway,right]{$\sigma|_{\benbrace*{0,1,2}}$} -- cycle;
		\draw[fill,Firebrick3] (0,0) circle[radius=0.05] node[above left]{$\sigma|_{\benbrace*{0}}$};
		
		% mittlerer  Simplex
		\draw[DodgerBlue3,dashed,ultra thick] (4,0) -- ++ (120:1) node[midway,right]{$\sigma|_{\benbrace*{1,2}}$};
		\draw[fill,Firebrick3, dashed,ultra thick] (3,0) -- ++ (60:1) node[midway, left]{$\sigma|_{\benbrace*{0,1}}$};
		
		% rechter Simplex
		\draw[Firebrick3,dashed,ultra thick] (6,0) -- ++(1,0) -- ++ (120:1) -- (6,0) node[midway,left]{$\sigma|_{\benbrace*{0,1,2}}$};
		\draw[fill,DodgerBlue3] (7,0) circle[radius=0.05] node[above right]{$\sigma|_{\benbrace*{2}}$};
	\end{tikzpicture}
	\caption{Einschränkung von singulären Simplizes am Beispiel des $2$-Simplex.}
	}
\end{figure}

\begin{bemerkung}[{name=[Schreibweise für Einschränkung auf eine Seite]}]
	Wir schreiben auch $\sigma|_{\benbrace*{0,\ldots , \slashedi, \ldots ,n}}$ für die $i$-te Seite von $\sigma$.
\end{bemerkung}

\begin{definition}[label=sec3:def_cup,{name=[Cup-Produkt]}]
	Sei $R$ ein Ring und $X$ ein topologischer Raum. 
	Wir definieren das \bet{Cup-Produkt} ($\cupp$"~Produkt)\index{Cup-Produkt, $\cupp$-Produkt} auf dem singulären Kokettenkomplex von $X$
	\[
		\cupp \colon C^p_\sing \enbrace*{X;R} \otimes C^q_\sing(X;R) \longrightarrow C^{p+q}_\sing \enbrace*{X;R}
	\]
	durch $(\alpha \cupp \beta)(\sigma) \coloneqq  (-1)^{p \cdot q} \cdot \alpha\enbrace*{\sigma|_{\benbrace{0,\ldots ,p}}} \cdot \beta \enbrace*{\sigma|_{\benbrace*{p, \ldots , p+q}}}$ für $\sigma \colon \abs{\Delta^{p+q}} \to X$.
\end{definition}

\begin{lemma}[{name=[Eigenschaften des Cup-Produktes]}]
	Das $\cupp$-Produkt hat folgende Eigenschaften:
	\begin{enumerate}[1)]
		\item $d^{p+q}(\alpha \cupp \beta) = d^p(\alpha) \cupp \beta + (-1)^p \cdot \alpha \cupp d^q(\beta)$
		\item Das $\cupp$-Produkt ist assoziativ.
		\item Für $f \colon X \to X'$ gilt $f^*(\alpha \cupp \beta)= f^*(\alpha) \cupp f^*(\beta)$ für $\alpha \in C^p_\sing(X';R)$, $\beta \in C^q_\sing(X';R)$.
		\item Sei $1_X \in C_\sing^0(X;R)$ mit $1_X(\sigma)=1_R$ für alle $\sigma \colon \abs{\Delta^0} \to X$. Dann gilt $1_X \cupp \alpha = \alpha = \alpha\cupp 1_X$.
	\end{enumerate}
\end{lemma}
\begin{beweis}
	Übung bzw. Notizen auf Homepage \todo{\TeX{}en wenn Zeit}
\end{beweis}

\begin{definition}[{name=[Cup-Produkt in Kohomologie]}]
	Die vom $\cupp$-Produkt auf dem singulären Kokettenkomplex induzierte Abbildung in Kohomologie
	\[
		\cupp \colon H^p(X;R) \otimes H^q(X;R) \longrightarrow H^{p+q}(X;R)
	\]
	ist das $\cupp$-Produkt in Kohomologie $\benbrace*{\alpha} \cupp \benbrace*{\beta} := \benbrace*{\alpha \cupp \beta}$. 
	\begin{itemize}
		\item Mit 1) folgt, dass das Cup-Produkt in Kohomologie wohldefiniert ist. 
		\item Nach 2) ist es assoziativ und somit ist $H^*(X) := \bigoplus_{p=0}^\infty H^p(X)$ ein Ring. 
		\item Nach 3) sind induzierte Abbildungen in Kohomologie multiplikativ und somit Ringhomomorphismen. 
		\item Nach 4) ist $\benbrace*{1_X} \in H^0(X;R)$ ein Einselement bezüglich $\cupp$ in $H^*$.
	\end{itemize}
\end{definition}

\begin{definition}[{name=[Kreuz-Produkt]}]
	Seien $X$ und $Y$ zwei topologische Räume. Seien $p_X \colon X \times Y \to X$ und $p_Y$ die Projektionen. 
	Das \bet{Kreuz-Produkt} ($\times$"~Produkt)\index{Kreuz-Produkt, $\times$-Produkt}
	\[
		\times \colon H^p(X;R) \otimes H^q(Y;R) \longrightarrow H^{p+q}(X \times Y;R)
	\]
	ist definiert durch $\xi \times \zeta := p_X^*(\xi) \cupp p_Y^*(\zeta)$. 
\end{definition}

\begin{bemerkung}[{name=[Zusammenhang von Kreuz- und Cup-Produkt]}]
	Das $\cupp$-Produkt kann wie folgt aus dem $\times$-Produkt zurückgewonnen werden: 
	Sei $\Delta \colon X \to X \times X$ die Diagonalabbildung $\Delta(x) := (x,x)$, dann gilt für $\xi \in H^p(X;R)$, $\zeta \in H^q(X;R)$
	\[
		\xi \cup \zeta = \Delta^*(\xi \times \zeta)
	\]
\end{bemerkung}

\begin{bemerkung}[{name=[relative Versionen des Cup- und Kreuz-Produktes]}]
	Es gibt auch relative Versionen des $\cupp$-Produkts und des $\times$-Produkts:
	\begin{enumerate}[(i)]
		\item Sei $A \subseteq X$ ein Unterraum. 
		Für $\alpha \in C^p(X,A;R)$ und $\beta \in C^q(X;R)$ sind dann $\alpha \cupp \beta, \beta \cupp \alpha \in C^{p+q}(X,A;R)$. 
		Wir erhalten relative $\cupp$-Produkte
		\[
			\begin{tikzcd}[row sep=0em]
				H^p(X,A;R) \otimes_R H^q(X;R) \rar & \minwidthbox{H^{p+q}(X,A;R)}{2cm} \\
				H^p(X;R) \otimes_R H^q(X,A;R) \rar & \minwidthbox{H^{p+q}(X,A;R)}{2cm} 
			\end{tikzcd}
		\]
		\item Seien $A \subseteq X$ und $Y \subseteq Y$ Unterräume.
		Dann erhalten wir relative $\times$-Produkte
		\[
			\begin{tikzcd}[row sep=0em]
				H^p(X,A;R) \otimes_R H^q(Y;R) \rar & H^{p+q}(X \times Y, A \times Y;R) \\
				H^p(X;R) \otimes_R H^q(Y,B;R) \rar & H^{p+q}(X \times Y,X \times B;R) 
			\end{tikzcd}
		\]
		\item Sind $A,B \subseteq X$ offen, so kann man mit dem Satz über kleine Simplizes ein relatives $\cupp$-Produkt
		\[
			\begin{tikzcd}
				H^p(X,A;R) \otimes_R H^q(X,B;R) \rar & H^{p+q}(X,A \cup B;R)
			\end{tikzcd}
		\]
		konstruieren. Siehe Aufgabe 3 von Blatt 4.
		\item Sind $A \subseteq X$ und $B \subseteq Y$ offen, so erhält man ein relatives $\times$-Produkt
		\[
			\begin{tikzcd}
				H^p(X,A;R) \otimes_R H^q(Y,B;R) \rar & H^{p+q}(X \times Y,X \times B \cup A \times Y;R)
			\end{tikzcd}
		\]
	\end{enumerate}
\end{bemerkung}

\begin{beispiel}[{name=[Cup-Produkt in der $0$-ten Kohomologie]}]
	Wir betrachten das $\cupp$-Produkt auf in der $0$-ten Kohomologie $\cupp \colon H^0(X;R) \otimes H^0(X;R) \to H^0(X;R)$. 
	Sei $A,B \subseteq X$ beide Vereinigungen von Wegzusammenhangskomponenten von $X$. 
	Seien $1_A, 1_B \in C^0(X;R)$ definiert durch 
	\[
		1_A(\sigma) = \begin{cases}
			1, &\text{ falls } \im \sigma \subseteq A\\
			0 , & \text{ sonst} 
		\end{cases}
	\]
	und $1_B$ analog. Dann ist $\benbrace*{1_A} \cupp \benbrace*{1_B} = \benbrace*{1_{A \cap B}}$.
\end{beispiel}

\begin{lemma}[label=sec3:lem1,{name=[Randabbildungen und Kreuz-Produkt]}]
	Sei $(X,A)$ ein Paar von topologischen Räumen und $Y$ ein weiterer topologischer Raum. Seien $\partial \colon H^p(A;R) \to H^{p+1}(X,A;R)$ und 
	$\partial^{\times} \colon H^{p+q}(A \times Y;R) \to H^{p+q+1}(X\times Y,A \times Y;R)$ die Verbindungshomomorphismen in den Paarfolgen. Dann gilt für $\xi \in H^p(A;R)$
	und $\zeta \in H^q(X;R)$
	\[
		\partial^\times (\xi \times \zeta) = \enbrace*{\partial \xi} \times \zeta
	\]
\end{lemma}
\begin{beweis}
	Sei $\xi= \benbrace*{\alpha}$ und $\zeta= \benbrace*{\beta}$. Setze $\alpha \colon C_p(A) \to R$ durch $\hat{\alpha}(\sigma)=0$ für $\sigma \colon \abs{\Delta^p} \to X$
	mit Bild nicht in $A$ zu $\hat{\alpha} \colon C_p(X) \to R$ fort. Dann ist
	\(
		\benbrace[\big]{d^p \hat{\alpha}}= \partial \benbrace*{\alpha}
	\)
	Weiter ist 
	\begin{align}
		\partial^\times(\xi \times \zeta)= \partial^\times \enbrace*{p_A^* \xi \cupp p_Y^* \zeta} = \partial^\times \benbrace[\big]{p_A^* \alpha \cupp p_Y^* \beta} &= 
		\enbrace[\Big]{d^{p+q} \enbrace[\big]{p^*_X (\hat{\alpha}) \cupp p^*_Y(\beta)}} \\
		\intertext{da $p^*_X(\hat{\alpha}) \cup p_Y^*(\beta)$ eine Fortsetzung von $p_Y^* (\alpha) \cup p_Y^*(\beta)$ ist. Weiter gilt}
		&= \benbrace[\big]{p_X^* \enbrace*{d^p(\hat{\alpha})} \cupp p_Y^*(\beta) \pm p_X^*(\hat{\alpha}) \cupp p_Y^* \enbrace*{d^q \beta}} \\
		&= p_X^* \benbrace[\big]{d^p(\hat{\alpha})} \cup p_Y^* \benbrace*{\beta} = \benbrace[\big]{d^p(\hat{\alpha})} \times \benbrace*{\beta} \\ 
		&= \partial \xi \times \zeta \qedhere
	\end{align}
\end{beweis}

\begin{lemma}[label=sec3:lem2,{name=[Isomorphismus für das Produkt mit $D^1$]}]
	Sei $\xi \in H^1 \enbrace*{D^1,\partial D^1;R}$ ein Erzeuger. Dann ist für jeden Raum $X$ die Abbildung 
	\mapdef{H^p(X;R)}{H^{p+1} \enbrace*{D^1 \times X, \partial D^1 \times X;R}}{\eta}{\xi \times \eta}{}
	ein Isomorphismus.
\end{lemma}
\begin{beweis} Wir betrachten folgendes Diagramm
	\[
		\begin{tikzcd}
			H^0(D^1;R) \otimes H^p(X;R) \dar \rar["\times"] & H^p(D^1 \times X;R) \dar\\
			H^0(\partial D^1;R)  \otimes H^p(X;R) \rar["\times"] \dar & H^p(\partial D^1 \times X;R) \dar \\
			H^1(D^1,\partial D^1;R) \otimes H^p(X;R) \rar["\times"] & H^{p+1} (D^1 \times X,\partial D^1 \times X;R)
		\end{tikzcd}
	\]
	Behauptung: Beide Spalten sind kurz exakt.
	Betrachte dazu folgenden Ausschnitt aus der Paarsequenz des Paares $(D^1, \partial D^1)$:
	\[
		\begin{tikzcd}[cramped]
			\Underbracket{H^0(D^1,\partial D^1;R)}{\StackText{\ref{0teKoho}}{=}0} \rar & H^0(D^1;R) \rar["j^*"] & H^0(\partial D^1;R) \rar & H^1(D^1,\partial D^1;R) \rar & 0
		\end{tikzcd}
	\]
	Seien $j \colon \partial D^1 \hookrightarrow D^1$ und $i_\pm \colon \set*{\pm 1} \hookrightarrow \partial D^1$ die Inklusionen. 
	Sei weiter $p \colon D^1 \to \set*{-1}$ die Projektion. 
	Dann gilt $p^* \circ i_-^* \circ j^* = \id_{H^0(D^1;R)}$. 
	Insbesondere spaltet die Sequenz und bleibt exakt nach Anwenden des Funktors ${-}\otimes H^p(X;R)$. 
	Genauso überlegt man sich, dass die Paarfolge
	\[
		\begin{tikzcd}
			\ldots \rar & H^p \enbrace*{D^1\times X, \partial D^1\times X; R}\rar["j^*"] & H^p \enbrace*{D^1 \times X;R}\rar["i^*"] & H^p \enbrace*{\partial D^1 \times X;R}\rar & \ldots 
		\end{tikzcd}
	\]
	in eine kurze exakte Folge 
	\[
		\begin{tikzcd}
			0 \rar & H^p \enbrace*{D^1\times X;R} \rar["i^*"] & H^p \enbrace*{\partial D^1 \times X;R} \rar["\partial"] \lar[bend left,"s^*"'] & H^{p+1} \enbrace*{D^1 \times X,\partial D^1 \times X;R} \rar & 0
		\end{tikzcd}
	\]
	zerfällt: Sei dazu $s \colon D^1 \times X \to \partial D^1 \times X$ gegeben durch $(t,x) \mapsto (0,x)$ und $i \colon \partial D^1 \times X \hookrightarrow D^1 \times X$ die Inklusion.
	Dann ist $i \circ s$ homotop zu $\id_{D^1 \times X}$, also $s^* \circ i^*=\id$ in Kohomologie und mit der Exaktheit folgt $j^*=0$. $s^*$ ist die gesuchte Spaltung.\todo{Revision 3}
	
	Die ersten zwei $\times$-Produkte sind Isomorphismen. 
	Dies folgt für das erste aus $H^0(D^1;R)\cong R$ und der Betrachtung des Erzeugers $1_R \in R$, denn für diesen gilt $p_{D^1}^*(1_R)=1_{D^1\times X}$ und somit
	\[
		p_{D^1}^*(1_R) \cupp p_{X}^*(\eta) = p_{X}^*(\eta)
	\]
	Für das zweite $\times$-Produkt betrachte \todo{Diagramm vervollständigen}
	\[
		\begin{tikzcd}
			H^0(\set*{-1}) \otimes H^p(X) \dar["(p_{-1})^* \otimes \id"] & \\
			H^0(\partial D^1) \otimes H^p(X) \rar["\times"] & H^p(\partial D^1 \times X) \dar["(i_{+1} \times \id_X)^*"] \\
			H^1(D^1,\partial D^1) \otimes H^p(X) & \mbox{ }
		\end{tikzcd}
	\]
	Nach \autoref{sec3:lem1} kommutiert das Diagramm und damit ist auch das dritte $\times$-Produkt ein Isomorphismus.
\end{beweis}

\begin{proposition}[label=sec3:prop_Hp-1]
	Sei $X$ ein topologischer Raum. Dann ist $H^p(S^1 \times X;R) \cong H^p(X;R) \oplus H^{p-1}(X;R)$.
\end{proposition}
\begin{beweis}
	Die Paarfolge zu $(S^1 \times X, \set*{\pt} \times X)$ zerfällt in kurze exakte Folgen
	\[
		\begin{tikzcd}
			H^p(S^1 \times X, \set*{\pt} \times X;r) \rar[hook] & H^p(S^1 \times X;R) \rar[two heads,"i^*"] & H^p \enbrace*{\set*{\pt} \times X;R}
		\end{tikzcd}
	\]
	Die Surjektivität folgt, da $\pi(\lambda,x) = (\pt,x)$ einen Spalt induziert. Damit folgt $H^p(S^1 \times X;R) \cong H^p(X;R) \oplus H^p(S^1 \times X, \set*{\pt} \times X;R)$.
	Die Behauptung folgt mittels Ausschneidung:
	\[
		H^p \enbrace*{S^1 \times X, \set*{\pt} \times X;R} \cong H^p \enbrace*{D^1 \times X, \partial D^1 \times X;R} \StackText{\ref{sec3:lem2}}{\cong} H^{p-1}(X;R) \qedhere
	\]
\end{beweis}

\begin{bemerkung}
	Sei $\xi \in H^1(S^1;R)$ ein Erzeuger und $\pi_X \colon S^1 \times X \to X$ die Projektion. Dann ist der Isomorphismus aus \autoref{sec3:prop_Hp-1} gegeben durch
	\[
		H^p \enbrace*{X;R} \oplus H^{p-1}(X;R) \longrightarrow H^p(S^1 \times X;R) \quad , \quad (\eta_1,\eta_2) \longmapsto \enbrace[\big]{\pi_X^*(\eta_1) + \xi\times \eta_2}
	\]
\end{bemerkung}
% section 3 (end)
\newpage

\section{Kommutativität des Cup-Produktes} % (fold)
\label{sec:4}

\begin{definition}[{name=[graduierte Algebra]}]
	Eine \Index{graduierte Algebra} ist ein graduierter Modul $A^* = \bigoplus_p A^p$ zusammen mit einer $R$-Algebrastruktur, für die das Produkt mit der Graduierung verträglich ist: Für $a \in A^p$ und $b \in A^q$ ist $a \cdot b \in A^{p+q}$. Sie heißt \Index{graduiert-kommutativ}, falls 
	\[
		a \cdot b =  (-1)^{p \cdot q} \cdot b \cdot a
	\]
	für alle $a \in A^p$ und $b \in A^q$ gilt.
\end{definition}

\begin{bemerkung}[{name=[Cup-Produkt auf $C^*$ nicht graduiert kommutativ]}]
	Das $\cupp$-Produkt auf $C^*_\sing(X;R)$ ist \emph{nicht} graduiert kommutativ, was man direkt aus der \autoref{sec3:def_cup} ablesen kann. 
\end{bemerkung}

\begin{lemma}[{name=[Spiegelung von singulären Simplizes]}]
	Für $\sigma \colon \abs{\Delta^p} \to X$ sei $\overline{\sigma} \colon \abs{\Delta^n} \to X$ definiert durch 
	\[
		\overline{\sigma}(t_0, \ldots , t_p ) :=  \sigma(t_p, \ldots ,t_0)
	\]
	Sei $\Sp \colon C_p(X,A;R) \to C_p(X,A;R)$ definiert durch $\Sp(\sigma) := (-1)^{\sfrac{p(p+1)}{2}} \cdot \overline{\sigma}$. 
	Dann erhalten wir eine induzierte Abbildung $\Sp^* \colon C^p(X,A;R) \to C^p(X,A;R)$ definiert durch $\Sp^*(\alpha) := \alpha \circ \Sp$. Es gilt
	\begin{enumerate}[a)]
		\item $\Sp$ ist eine Kettenabbildung und $\Sp^*$ ist eine Kokettenabbildung
		\item Ist $R$ kommutativ, so gilt $\beta \cupp \alpha = (-1)^{p \cdot q}\cdot \Sp^*\enbrace[\big]{\Sp^*(\alpha) \cupp \Sp^*(\beta)}$
	\end{enumerate}
\end{lemma}
\begin{beweis}
	\begin{enumerate}[a)]
		\item Sei $\sigma \colon \abs{\Delta^p} \to X$ ein singulärer Simplex. Dann gilt
		\begin{align}
			\Sp(d_p \sigma) = \Sp\enbrace*{\sum_{i=0}^{p} (-1)^i \sigma|_{\benbrace*{0, \ldots , \slashedi, \ldots ,p}}} 
			&= \sum_{i=0}^{p} (-1)^i \Sp \enbrace*{\sigma|_{\benbrace*{0, \ldots , \slashedi, \ldots , p}}} \\
			&= \sum_{i=0}^{p} (-1)^i \cdot (-1)^{\frac{p(p-1)}{2}} \,\sigma|_{\benbrace*{p, \ldots , \slashedi, \ldots , 0}}
		\end{align}
		wobei $\slashedi$ an der $p-i$-ten Stellen von vorn steht. Weiter ist
		\begin{align}
			d_p \enbrace[\big]{\Sp(\sigma)} = d_p \enbrace*{(-1)^{\frac{p(p+1)}{2}} \,\sigma|_{\benbrace*{p, \ldots ,0}}} &= (-1)^{\frac{p(p+1)}{2}} \sum_{i=0}^{p} (-1)^i\, 
			\sigma\big|_{\benbrace*{p, \ldots , \cancel{p-i}, \ldots ,0}} \\ 
			&= (-1)^{\frac{p(p+1)}{2}} \sum_{i=0}^{p} (-1)^{p-i} \,\sigma\big|_{\benbrace*{p, \ldots ,\slashedi, \ldots ,0}}
		\end{align}
		Nun ist $\frac{p(p-1)}{2} \equiv \frac{p(p+1)}{2} + p$ und $x +i \equiv x-i$ und es folgt die Gleichheit.
		\item Es ist $p=\abs*{\alpha}$, $q=\abs*{\beta}$ und $n=p+q$.
		Für einen singulären Simplex $\sigma \colon \abs{\Delta^n} \to X$ gilt 
		\begin{align}
			\enbrace*{\Sp^* \enbrace[\big]{\Sp^*(\alpha) \cupp \Sp^*(\beta)}}(\sigma) &= (-1)^{\frac{n(n+1)}{2}} \cdot \enbrace[\big]{\Sp^*(\alpha) \cupp \Sp^*(\beta)} 
			\enbrace*{\sigma|_{\benbrace*{n, \ldots ,0}}} \\
			&= (-1)^{\frac{n(n+1)}{2}} \cdot (-1)^{p \cdot q}\cdot \Sp^*(\alpha) \enbrace*{\sigma_{\benbrace*{n, \ldots , q}}} \cdot \Sp^*(\beta) 
			\enbrace*{\sigma|_{\benbrace*{q, \ldots ,0}}} \\
			&= (-1)^{\frac{n(n+1)}{2}} (-1)^{p \cdot q} (-1)^{\frac{p(p+1)}{2}} (-1)^{\frac{q(q+1)}{2}} \alpha \enbrace*{\sigma|_{\benbrace*{q, \ldots ,n}}}
			\beta \enbrace*{\sigma|_{\benbrace{0, \ldots ,q}}} \\
			&= (-1)^{\frac{n(n+1)}{2}} (-1)^{p \cdot q} (-1)^{\frac{p(p+1)}{2}} (-1)^{\frac{q(q+1)}{2}} \beta \enbrace*{\sigma|_{\benbrace{0, \ldots ,q}}} 
			\alpha \enbrace*{\sigma|_{\benbrace*{q, \ldots ,n}}} \\
			&\stackrel{!}{=} (-1)^{p \cdot q} \enbrace*{\beta \cupp \alpha}(\sigma) = \cancel{(-1)^{p \cdot q}} \cdot \cancel{(-1)^{p \cdot q}} \,\beta \enbrace*{\sigma|_{\benbrace{0, \ldots ,q}}} 
			\alpha \enbrace*{\sigma|_{\benbrace*{q, \ldots ,n}}}
		\end{align}
		Weiter gilt nun
		\begin{align}
			\sfrac{n(n+1)}{2} + \sfrac{p(p+1)}{2} + \sfrac{q(q+1)}{2} &= (1+ \ldots + n) + (1+ \ldots + p) + (1+ \ldots + q)\\ 
			&\equiv \enbrace[\big]{(p+1) + \ldots + n} + (1+ \ldots + q)\\
			&= (n+1) q = (p+q+1) q = pq + q^2 + q \equiv p \cdot q
		\end{align}
		da $q^2 \equiv q$ und es folgt die Behauptung. \qedhere
	\end{enumerate}
\end{beweis}

\begin{lemma}[{name=[Spiegelung ist kettenhomotop zur Identität]}]
	$\Sp \colon C_*(X,A) \to C_*(X,A)$ ist kettenhomotop zur Identität. 
\end{lemma}
\begin{beweis}
	$\Sp$ ist eine natürliche Transformation mit $\Sp_0 =\id \colon C_0(X) \to C_0(X)$.
	Nach \autoref{prop:H0_kettenhomotop} bzw. \autoref{korr:H0_kettenhomotop} ist $\Sp$ damit kettenhomotop zur Identität.
\end{beweis}

\begin{satz}[{name=[Das $\cupp$-Produkt in Kohomologie ist graduiert kommutativ]}]
	Das $\cupp$-Produkt in Kohomologie ist graduiert kommutativ. 
\end{satz}
\begin{beweis}
	Seien $\xi \in H^p(X,A;R)$, $\eta \in H^q(X,B;R)$. Wähle $\alpha \in C^p(X,A;R)$, $\beta \in C^q(X,B;R)$ mit $\xi= [\alpha]$ und $\eta= [\beta]$. Dann ist, da mit $\Sp$ auch $\Sp^*$ kokettenhomotop zur Identität ist.
	\begin{align}
		\xi \cupp \eta = \benbrace*{\alpha \cupp \beta} = (-1)^{pq} \cdot \benbrace*{\Sp^* \enbrace[\big]{\Sp^*(\beta) \cupp \Sp^*(\alpha)}} &= (-1)^{pq} \cdot 
		\benbrace[\big]{\Sp^*(\beta) \cupp \Sp^*(\alpha)} \\
		&= (-1)^{pq} \cdot \benbrace[\big]{\Sp^*(\beta)} \cupp \benbrace[\big]{\Sp^*(\alpha)}\\ 
		&= (-1)^{pq} \cdot [\beta] \cupp [\alpha] = (-1)^{pq} \cdot \eta \cupp \xi \qedhere
	\end{align}
\end{beweis}

\begin{proposition}[label=prop:H0_kettenhomotop,{name=[Nullkettenhomotopie von natürlichen Transformationen]}]
	Sei $\tau_* \colon C_*^\sing \to C_*^\sing$ eine natürliche Transformation durch Kettenabbildungen. Induziert $\tau$ die triviale Abbildung auf $H_0(\set*{\pt})$, so ist 
	$\tau_*$ kettenhomotop zur Nullabbildung, das heißt es gibt eine natürliche Transformation $\eta_* \colon C_*^\sing \to C_{*+1}^\sing$ mit 
	\[
		d^\sing_{*+1} \circ \eta_* + \eta_{*+1} \circ d_*^\sing = \tau_* 
	\]
\end{proposition}
\begin{beweis}
	Da $\tau$ auf $H_0(\set{\pt})$ trivial ist, ist $\tau \enbrace*{\id_{\abs*{\Delta^0}}}=0$. Es folgt für $\sigma \colon \abs{\Delta^0} \to X$
	\[
		\tau(\sigma) = \tau \enbrace*{\sigma_* \enbrace*{\id_{\abs*{\Delta^0}}}} = \sigma_* \enbrace*{\tau \enbrace*{\id_{\abs*{\Delta^0}}}} = 0
	\]
	Es folgt, dass $\tau_0 \colon C_0(X) \to C_0(X)$ trivial ist. Wir setzen $\eta_0 = 0$ und und definieren $\eta_1, \eta_2, \ldots $ induktiv: Seien $\eta_1, \ldots,\eta_{k-1}$
	schon konstruiert mit $d^\sing_{l+1} \circ \eta_l + \eta_{l-1} \circ d^\sing_{l}=\tau_l$ für $l=0, \ldots , k-1$. Es folgt
	\begin{align}
		d_k^\sing \enbrace*{\tau_k -  \eta_{k-1} \circ d_k^\sing } &= d^\sing_k \circ \tau_k - d^\sing_k \circ \eta_{k-1} \circ d^\sing_k \\
		&= \tau_{k-1} \circ d_k^\sing -  \enbrace*{\tau_{k-1} - \eta_{k-2} \circ d_{k-1}^\sing} \circ d_k^\sing =0
	\end{align}
	Setze nun $\id_{\abs*{\Delta^k}} \in C_k^\sing( \abs{\Delta^k})$ in die Formel ein. 
	Da $\abs*{\Delta^k}$ kontrahierbar ist, gilt $H_k \enbrace*{\abs{\Delta^k}} =0$. Außerdem gilt
	\[
		d_k^\sing \enbrace[\Big]{\tau_k \enbrace*{\id_{\abs*{\Delta^k} }} - \eta_{k-1} \enbrace*{d_k^\sing \enbrace*{\id_{\abs*{\Delta^k} }}}} =0
	\]
	und somit gibt es ein $a_k \in C_{k+1} \enbrace*{\abs{\Delta^k} }$, sodass $d_{k+1}^\sing(a_k)$ das Innere der Klammer ist. 
	Für einen singulären Simplex $\sigma \colon \abs{\Delta^k} \to X$ setzen wir nun $\eta_k(\sigma) := \sigma_*(a_k)$.
	Dann gilt 
	\[
		\enbrace*{d_{k+1}^\sing \circ \eta_k + \eta_{k-1} \circ d_k^\sing - \tau_k}(\sigma) = \sigma_* \enbrace[\big]{d_{k+1}^\sing \circ \eta_k + \eta_{k-1} \circ d_k^\sing - \tau_k} (\id_{\abs*{\Delta^k}}) =0 
	\]
	nach Konstruktion.
	Damit ist $d_{k+1}^\sing \circ \eta_k + \eta_{k-1} \circ d_k^\sing = \tau_k$ wie gewünscht und $\eta_*$ ist die gesuchte Kettenhomotopie.
\end{beweis}

\begin{korollar}[label=korr:H0_kettenhomotop,{name=[Kettenhomotopie von natürlichen Transformationen]}]
	Seien $\tau, \tau'$ zwei natürliche Transformationen $C_*^\sing \to C_*^\sing$ durch Kettenabbildungen, sodass $H_0(\tau) = H_0(\tau')$ für $X= \set*{\pt}$ gilt. 
	Dann sind $\tau$ und $\tau'$ kettenhomotop.
\end{korollar}
\begin{beweis}
	Betrachte $\tau-\tau'$ und wende \autoref{prop:H0_kettenhomotop} an. \todo{Revision 4}
\end{beweis}

\begin{definition}[{name=[Äußere Algebra]}]
	Sei $R$ ein kommutativer Ring. Die \Index{freie graduiert kommutative Algebra}\footnote{$a^2=0$ falls $a$ ungerade} in Erzeugern $a_1, \ldots, a_k$ von Grad $\abs*{a_1}, \ldots , \abs*{a_k}$ über $R$ ist die 
	graduiert kommutative Algebra $\Lambda_R \benbrace*{a_1,\ldots,a_k}$, deren Elemente von Grad $N$ formale Linearkombinationen der Form 
	\[
		\smashoperator[r]{\sum_{\substack{i_1, \ldots ,i_n \ge 0 \\ i_1 \abs*{a_1} + \ldots + i_n \abs*{a_n} = N , i_j \in \set*{0,1} \text{ für } \abs*{a_j} \text{ ungerade}}}} 
		\enspace r_{i_1, \ldots ,i_n} \cdot a_1^{i_1} \cdot a_2^{i_2} \cdots a_n^{i_n}
	\]
	sind.
	Das Produkt ist das eindeutige $R$-lineare assoziative graduiert kommutative Produkt mit $a_\lambda^j \cdot a_{\lambda'}^{j'} = a_\lambda^j a_{\lambda'}^{j'}$ für 
	$\lambda< \lambda'$ und $a_\lambda^j \cdot a_\lambda^{j'} = a_\lambda^{j+j'}$ für $\abs*{a_\lambda}$ ungerade.
\end{definition}

\begin{beispiel}[{name=[Polynomring und äußere Algebra]}]
	Sind alle $\abs*{a_i}$ gerade, so ist $\Lambda_R \benbrace*{a_1,\ldots,a_n}$ der Polynomring über $R$ in den Variablen $a_1, \ldots ,a_n$.
	Sind alle $\abs*{a_i}$ ungerade, so ist $\Lambda_R \benbrace*{a_1, \ldots a_n}$ die \Index{äußere Algebra} in den Variablen $a_1, \ldots , a_n$.
	
	\emph{Definition der äußeren Algbra aus \cite[Ex. 3.15]{Hatcher}:}
	Die \Index{äußere Algebra} $\Lambda_R \benbrace*{a_1,a_2,\ldots }$ über einem kommutativen Ring $R$ ist der freie $R$-Modul, dessen Basis aus endlichen Produkten $a_{i_1} \cdots a_{i_k}$ mit $i_1 < \ldots < i_k$ besteht, versehen mit der Multiplikation, die durch die Forderungen 
	\[
		a_i a_j= -a_j a_i \quad \text{ und } \quad  a_i^2=0
	\]
	eindeutig bestimmt ist. 
	Das leere Produkt der $a_i$ liefert ein Einselement.
\end{beispiel}

\begin{satz}[{name=[Kohomologie des Torus]}]
	Sei $R$ ein kommutativer Ring. Sei $\xi \in H^1 \enbrace*{S^1;R}$ ein Erzeuger. Sei $p_i \colon T^n = (S^1)^n \to S^1$ die Projektion auf die $i$-te Koordinate. Sei 
	$\xi_i := (p_i)^*(\xi) \in H^1(T^n;R)$. Dann ist 
	\[
		H^*(T^n;R) \cong \Lambda_R \benbrace*{\xi_1, \ldots ,\xi_n}
	\]
	als graduiert kommutative Ringe.
\end{satz}
\begin{beweis}
	Per Induktion nach $n$. Für $n=0$ ist nichts zu zeigen. 
	Für den Induktionsschritt erinnern wir uns daran, dass $\xi \times {-} \colon H^p(T^{n-1};R) \to H^p(T^n;R)$ ein 
	Isomorphismus ist. Damit folgt die Behauptung.
\end{beweis}

\begin{beispiel}[{name=[Kohomologieklassen des Torus]}]
	$\xi_1 \cdots \xi_n$ ein Erzeuger ein Erzeuger von $H^n(T^n;R)$.
\end{beispiel}
% section 4 (end)
\newpage

\section{Azyklische Kettenkomplexe} % (fold)
\label{sec:5}

\begin{frage}
	Sei $f  \colon C_* \to D_*$ eine Kettenabbildung, die einen Isomorphismus in Homologie induziert. Ist dann $f$ schon eine Kettenhomotopieäquivalenz?
\end{frage}
% Frage: Sei $f  \colon C_* \to D_*$ eine Kettenabbildung, die einen Isomorphismus in Homologie induziert. Ist dann $f$ schon eine Kettenhomotopieäquivalenz?

\begin{frage}
	Sei $\mathcal{U}$ eine offene Überdeckung von $X$. 
	Dann haben wir schon gesehen, dass die Inklusion $i \colon C_*^\mathcal{U}(X;R) \hookrightarrow C_*^\sing(X;R)$ einen Isomorphismus in Homologie induziert. 
	Ist $i$ eine Kettenhomotopieäquivalenz?
\end{frage}

\begin{definition}[{name=[azyklischer Kettenkomplex]}]
	Ein $R$-Kettenkomplex heißt \bet{azyklisch}\index{azyklischer Kettenkomplex}, falls seine Homologiegruppen trivial sind, also falls der Kettenkomplex eine lange exakte Folge ist.
\end{definition}

\begin{beispiel}[{name=[Beispiel eines azyklischen Kettenkomplex]}]
	Unser altbekanntes Beispiel
	\[
		\begin{tikzcd}
			0 \rar & \mathbb{Z} \rar["\cdot 2",hook] & \mathbb{Z} \rar[two heads] & \sfrac{\mathbb{Z}}{2 \mathbb{Z}} \rar & 0
		\end{tikzcd}
	\]
	ist ein azyklischer Kettenkomplex, der nicht kettenkontraktibel ist, da es keine nicht trivialen Abbildungen $\sfrac{\mathbb{Z}}{2 \mathbb{Z}} \to \mathbb{Z}$ gibt. Damit ist 
	die Antwort auf die erste Frage \enquote{Nein}.
\end{beispiel}

\begin{definition}[{name=[projektiver Modul]},label=def_projektiv]
	Ein $R$-Modul $P$ heißt \bet{projektiv}\index{projektiver Modul}, falls es zu jeder surjektiven, $R$-linearen Abbildung $f \colon M \to N$ und jeder $R$-linearen Abbildung
	$\varphi \colon P \to N$ eine $R$-lineare Abbildung $\hat{\varphi} \colon P \to M$ gibt mit $\varphi= f \circ \hat{\varphi}$. Anders ausgedrückt kommutiert folgendes Diagramm:
	\[
		\begin{tikzcd}
			\mbox{ } & M \dar[two heads,"f"] \\
			P \rar["\varphi"] \urar["\hat{\varphi}",dashed] & N
		\end{tikzcd}
	\]
\end{definition}

\begin{beispiel}[{name=[Freie Moduln sind projektiv]}]
	Freie Moduln sind projektiv, wie man durch Anwendung der universellen Eigenschaft eines freien Moduls sehr einfach sieht (wurde auch in Topologie \RM{1} bewiesen). 
\end{beispiel}

\begin{satz}[{name=[Projektive Moduln als direkte Summanden in freien Moduln]}]
	Für einen $R$-Modul $P$ sind äquivalent:
	\begin{enumerate}[i),itemsep=0pt]
		\item $P$ ist projektiv
		\item $P$ ist ein direkter Summand in einem freien $R$-Modul.
	\end{enumerate}
\end{satz}
\begin{beweis}
	Zur ersten Implikation:
	\[
		\begin{tikzcd}
			\mbox{}&  R[P] \dar[two heads,"p"] \\
			P \rar["\id"] \urar[dashed,"i"]& P
		\end{tikzcd}
	\]
	Mit $R[S] \cong P \oplus \ker p$ folgt nun die gewünschte Aussage.
	Die zweite Implikation ergibt sich wie folgt: Nach Vorraussetzung gibt es $i \colon P \to F$ und $p \colon F \to p$, wobei $F$ frei ist und $p \circ i = \id_P$ gilt.
	Seien nun $M$, $N$, $f$ und $\varphi$ wie in \autoref{def_projektiv}.
	Wir erhalten folgendes Diagramm
	\[
		\begin{tikzcd}[sep=huge]
			F \rar["\hat{\psi}"] \dar[two heads,"p"',xshift=-3pt]  \drar["\psi=\varphi \circ p",sloped,pos=0.16] & M \dar[two heads,"f"] \\
			P \rar["\varphi"] \uar[hook,"i"',xshift=3pt] & N
		\end{tikzcd}
	\]
	Da $F$ frei und insbesondere projektiv ist, existiert ein Homomorphismus $\hat{\psi}$ mit $f \circ \hat{\psi}=\psi$.
	Wir setzen nun $\hat{\varphi} := \hat{\psi} \circ i$. 
	Dann gilt $f \circ \hat{\varphi} = f \circ \hat{\psi} \circ i = \varphi \circ p \circ i = \varphi$ und $P$ ist projektiv. \todo{Revision 5}
\end{beweis}

\begin{satz}
	Sei $C_*$ ein azyklischer $R$-Kettenkomplex. Sind alle Kettenmoduln $C_n$ projektiv, so ist $C_*$ kettenkontrahierbar.
\end{satz}
\begin{beweis}
	Wir konstruieren induktiv $s_n \colon C_n \to C_{n+1}$, sodass $d_{n+1} \circ s_n + s_{n-1} \circ d_n = \id_{C_n}$ gilt. 	
	Für $n=0$ gehen wir wie folgt vor: Da $C_*$ azyklisch ist, ist $d_1 \colon C_1 \to C_0$ surjektiv und es gibt $s_0 \colon C_0 \to C_1$ mit $d_1 \circ s_0 = \id_{C_0}$, weil 
	$C_0$ projektiv ist.
	
	Zum Induktionsschritt $(n-1) \mapsto n$: Da $C_+$ azyklisch ist, ist $d_{n+1} \colon C_{n+1} \to \im d_{n+1}=\ker d_n$ surjektiv. Wegen
	\begin{align}
		d_n \enbrace*{\id_{C_n} - s_{n-1} \circ d_n} &= d_n - \Underbrace{d_n \circ s_{n-1}}{=\id_{C_{n-1}} - s_{n-2} \circ d_{n-1}} \circ d_n \\
		&= d_n - \enbrace*{- s_{n-2} \circ d_{n-1} + \id_{C_{n-1}}} \circ d_n =0
	\end{align}
	liegt das Bild von $\id_{C_n} - s_{n-1} \circ d_n$ im Kern von $d_n$. Da $C_n$ projektiv ist, gibt es $s_n \colon C_n \to C_{n+1}$ mit 
	$d_{n+1} \circ s_n = \id_{C_n} - s_{n-1} \circ d_n$.
	
	\[
		\begin{tikzcd}[column sep=2.5cm,row sep=large]
			\mbox{ } & C_{n+1} \dar["d_{n+1}",two heads] \\
			C_n \urar["s_n", dashed] \rar["\id_{C_n} - s_{n-1} \circ d_n"'] & \im d_{n+1} = \ker d_n
		\end{tikzcd}
	\]
\end{beweis}

\begin{satz}[label=satz:5:kettenhomotopie]
	Sei $f \colon C_* \to D_*$ eine Kettenabbildung. Seien alle Kettenmoduln von $C_*$ und $D_*$ projektiv. Dann sind äquivalent: 
	\begin{enumerate}[(i),itemsep=0pt]
		\item $f$ ist eine Kettenhomotopieäquivalenz
		\item $f$ induziert einen Isomorphismus in Homologie.
	\end{enumerate}
\end{satz}
\begin{beweis}
	Die Implikation (i) $\Rightarrow$ (ii) ist klar (siehe Topologie \RM{1}.) Die zweite Implikation besprechen wir später (siehe Seite~\pageref{bew:satz5}).
\end{beweis}

\begin{definition}[{name=[Kegel und Einhängung]}]
	Sei $f \colon C_* \to D_*$ eine Kettenabbildung. Der \Index{Kegel} $\Keg(f)_*$ von $f$ ist der Kettenkomplex, dessen $n$-ter Kettenmodul $D_n \oplus C_{n-1}$ ist und dessen 
	$n$-te Randabbildung gegeben ist durch
	\[
		d^{\Keg(f)}_n := \begin{pmatrix}
			d^D_n & f\\
			0 & -d^C_{n-1}
		\end{pmatrix}
	\]
	Die \Index{Einhängung} $\Sigma C_*$ von $C_*$ ist der Kettenkomplex, dessen $n$-ter Kettenmodul $C_{n-1}$ ist und dessen $n$-te Randabbildung $-d_{n-1}^C$ ist.
\end{definition}

\begin{lemma}
	$f$ induziert genau dann einen Isomorphismus in Homologie, wenn $\Keg(f_*)$ azyklisch ist.
\end{lemma}
\begin{beweis}
	Betrachte die lange exakte Folge zur kurzen exakten Folge von Kettenkomplexen
	\[
		\begin{tikzcd}[cramped]
			0 \rar & D_* \rar["i"] & \Keg(f)_* \rar["p"] & \Sigma C_* \rar & 0
		\end{tikzcd}
	\]
	Der Verbindungshomomorphismus $H_{n+1}(\Sigma C_*) \to H_n(D_*)$ entspricht unter $H_{n+1}(\Sigma C_*) \cong H_n(C_*)$ genau der von $f$ induzierten Abbildung.
\end{beweis}

\begin{lemma}
	$f$ ist genau dann eine Kettenhomotopieäquivalenz, wenn $\Keg(f)_*$ kettenkontraktibel ist.
\end{lemma}
\begin{beweis}
	Sei $\begin{psmallmatrix} \alpha_n & \beta_{n-1} \\ \gamma_n & \delta_{n-1} \end{psmallmatrix}$ eine Kettenkontraktion für $\Keg(f)$. Es folgt 
	\begin{align}
		f_n \circ \gamma_n + d^D_{n+1} \circ \alpha_n + \alpha_{n-1} \circ d_n^D &= \id_{D_n} \tag*{(1)} \\
		d^D_{n+1} \circ \beta_{n-1} + f_n \circ \delta_{n-1} + d_{n-1} \circ f_{n-1} - \beta_{n-2} \circ d_{n-1}^C &= 0 \tag*{(2)}\\
		- d^C_n \circ \gamma_n + \gamma_n \circ d_n^D &= 0 \tag*{(3)} \\
		- d_n^C \circ \delta_{n-1} + \gamma_{n-1} \circ f_{n-1} - \delta_{n-2} \circ d^C_{n-1} &= \id_{C_{n-1}} \tag*{(4)}
	\end{align}
	Daraus folgt
	\begin{align}
		&\gamma \text{ ist eine Kettenabbildung} \tag*{(1')} \label{sec5:1-} \\
		&\id_D \simeq f \circ g \text{ via } \alpha \tag*{(3')} \label{sec5:3-}\\
		& \gamma \circ f \simeq \id_C \text{ via }  \delta \tag*{(4')} \label{sec5:4-}
	\end{align}
	Es folgt, dass $\gamma$ eine Kettenabbildung ist, $\id_D \simeq f \circ g$ via $\alpha$ und $\gamma \circ f \simeq \id_C$ via $\delta$. Damit ist $f$ eine 
	Kettenhomotopieäquivalenz mit Kettenhomotopieinversem $\gamma$. 
	Gelten umgekehrt \ref{sec5:1-}, \ref{sec5:3-} und \ref{sec5:4-}, so setzt man $\hat{\delta} := \delta + \gamma \circ f - \delta \gamma f$ und 
	$\beta := \alpha f \gamma + \alpha^2 f + f \delta^2$ und 
	\[
		S := \begin{pmatrix}
			\alpha & \beta \\
			\gamma & \hat{\delta}
		\end{pmatrix}
	\]
	eine Kettenkontraktion für $\Keg(f)$.
\end{beweis}

\begin{beweis}[name={von \autoref{satz:5:kettenhomotopie} zweite Implikation},label=bew:satz5]
	Ist $H_n(f)$ ein Isomorphismus, so ist $\Keg(f)$ azyklisch. Die direkte Summen von projektiven Moduln ist wieder projektiv sind, sind die Kettenabbildungen von $\Keg(f)$
	projektiv, falls die von $C_*$ und $D_*$ projektiv sind. Daher ist $\Keg(f)$ kettenkontraktibel und $f$ eine Kettenhomotopieäquivalenz.
\end{beweis}

\begin{korollarB}
	Sei $\mathcal{U}$ eine offene Überdeckung von $X$. Dann ist die Inklusion $C^{\mathcal{U}}_*(X,A) \to C_*(X,A)$ eine Kettenhomotopieäquivalenz.
\end{korollarB}

\begin{korollar}
	Singuläre Kohomologie erfüllt das Ausschneidungsaxiom.
\end{korollar}
\begin{beweis}
	Für eine offene Überdeckung $\mathcal{U}$ von $X$ sei $C^*_{\mathcal{U}}(X,A;V) := \Hom \enbrace*{C_*^{\mathcal{U}}(X,A),V}$. Da die Inklusion $i \colon C^{\mathcal{U}}_*(X,A) \to C_*(X;A)$ 
	eine Kettenhomotopieäquivalenz ist, ist auch $i^* \colon C^*(X,A;V) \to C^*_{\mathcal{U}}(X,A;V)$ eine Kokettenhomotopieäquivalenz. Nun genauso wie für Homologie.
\end{beweis}
% section 5 (end)
\newpage

\section{Das Cap-Produkt und die Poincaré-Dualitätsabbildung} % (fold)
\label{sec:6}

\begin{definition}[{name=[Cap-Produkt]}]
	Sei $X$ ein topologischer Raum und $R$ ein kommutativer Ring. Wir definieren das \bet{Cap-Produkt} ($\capp$-Produkt)\index{Cap-Produkt, $\capp$-Produkt}\todo{Revision 6}
	\[
		\capp \colon C^p(X;R) \otimes_R C_{p+q}(X;R) \longrightarrow C_q(X;R)
	\]
	durch $\alpha \capp \sigma := (-1)^{pq} \cdot \alpha \enbrace*{\sigma|_{\benbrace*{q,\ldots ,p+q}}} \sigma|_{\benbrace*{0,\ldots ,p}}$ für $\alpha \in C^p(X;R)$ und 
	$\sigma \colon \abs{\Delta^{p+q}}\to X$.
\end{definition}

\begin{lemma}[{name=[Produktregel für das Cap-Produkt]}]
	Es gilt $d_q(\alpha \capp \sigma) = d^p \alpha \capp \sigma + (-1)^p \alpha \capp d_{p+q} \sigma$
\end{lemma}
\begin{beweis}
	Sei $n := p+q$. Dann gilt
	\begin{align}
		d_q(\alpha \capp \sigma) &= (-1)^{pq} \cdot \alpha \enbrace*{\sigma|_{\benbrace*{q,\ldots ,p+q}}} d_q \sigma|_{\benbrace*{0,\ldots ,p}} \\
		&= \sum_{i=0}^{q} (-1)^{pq+i}  \alpha \enbrace*{\sigma|_{\benbrace*{q,\ldots ,p+q}}} \sigma|_{\benbrace*{0, \ldots ,\slashedi,\ldots ,q}} 
	\end{align}
	Weiter ist 
	\begin{align}
		d^p(\alpha) \capp \sigma &= (-1)^{(p+1)(q-1)} \cdot (d^p \alpha) \enbrace*{\sigma|_{\benbrace*{q-1,\ldots ,n}}} \sigma|_{\benbrace*{0,\ldots ,q-1}} \\
		&= \Underbrace{(-1)^{pq+p+q+1+1+p}}{=(-1)^{pq+q}} \cdot \alpha \enbrace*{ d_{p-1} \enbrace*{\sigma|_{\benbrace*{q-1,\ldots ,n}}} \sigma|_{\benbrace*{0, \ldots ,q-1}}} \\
		&= \sum_{i=q-1}^{n} (-1)^{pq+q+i+(q+1)} \cdot \alpha \enbrace*{\sigma|_{\benbrace*{q-1,\ldots,\slashedi,\ldots ,n}}} \sigma|_{\benbrace*{0, \ldots ,q-1}}
	\end{align}
	und
	\begin{align}
		\alpha \capp d_{p+q}(\sigma) &= \sum_{i=0}^{n} (-1)^i \alpha \capp \sigma |_{\benbrace*{0, \ldots , \slashedi,\ldots ,n}} \\
		&= \sum_{i=0}^{q-1} (-1)^{i+p(q-1)} \alpha \enbrace*{\sigma|_{\benbrace*{q,\ldots ,n}}} \sigma|_{\benbrace*{0,\ldots , \slashedi,\ldots ,q}}\\
		&\hspace{1em}+ \sum_{i=q}^{n} (-1)^{i+p(q-1)} \alpha \enbrace*{\sigma|_{\benbrace*{q-1,\ldots,\slashedi,\ldots,n}}} \sigma|_{\benbrace*{0, \ldots ,q-1}}
	\end{align}
	Die Vorzeichen werden schon passen \ldots 
\end{beweis}

\begin{lemma}
	Sei $A \subseteq X$. Das $\capp$-Produkt auf $C^*$ induziert folgende $\capp$-Produkte in $H^*$
	\[
		\begin{tikzcd}[row sep=0em]
			\minwidthbox{H^p(X;R) \otimes_R H_n(X;R)}{4.2cm} \rar & \minwidthbox{H_q(X;R)}{2cm} \\
			\minwidthbox{H^p(X,A;R) \otimes_R H_n(X,A;R)}{4.2cm} \rar & \minwidthbox{H_q(X;R)}{2cm} \\
			\minwidthbox{\hfill H^p(X;R) \otimes_R H_n(X,A;R)}{4.2cm} \rar & \minwidthbox{H_q(X,A;R)}{2cm}
		\end{tikzcd}
	\]
\end{lemma}
\begin{beweis}
	Genauso wie beim $\cupp$-Produkt.
\end{beweis}

\begin{bemerkung}[label=bem_augmentation_cap] \leavevmode
	\begin{itemize}
		\item Sind $A,B \subseteq X$ offen, so erhalten wir auch ein $\capp$-Produkt
		\[
			H^p(X,A;R) \otimes_R H_n(X,A\cup B;R) \longrightarrow H_q(X,B;R)
		\]
		\item Sei $\epsilon \colon H_0(X;R) \to R$ die \Index{Augmentation}, also $\epsilon \enbrace*{\sum_{i}r_i \sigma_i} = \sum_i r_i$. Dann ist 
		\[
			\epsilon \enbrace*{\xi \capp x} = \xi(x)
		\]
		für $\xi \in H^p(X;R)$ und $x \in H_p(X;R)$.
	\end{itemize}
\end{bemerkung}

\begin{lemma}
	Für $f \colon X \to Y$ stetig gilt $\xi \capp f_*(x) = f_* \enbrace*{f^* \xi \capp x}$ für $\xi \in H^p(Y;R)$ und $x \in H_{p+q}(X;R)$.
\end{lemma}

\begin{lemma}
	Für $\xi \in H^p(X;R)$ und $\eta \in H^q(X;R)$, $x \in H_{p+q+r}(X;R)$ gilt 
	\[
		\enbrace*{\xi \cupp \eta} \capp x = \xi \capp \enbrace*{\eta \capp x}
	\]
\end{lemma}
\begin{beweis}
	Sei $\alpha \in C^p(X)$, $\beta \in C^q(X)$ und $\sigma \colon \abs{\Delta^{p+q+r}}\to X$ ist
	\begin{align}
		\enbrace*{\alpha \cupp \beta} \capp \sigma &= (-1)^{(p+q)r} \cdot (\alpha \cupp \beta)\enbrace*{\sigma|_{\benbrace*{r,\ldots ,n}}}   \sigma|_{\benbrace*{0,\ldots ,r}} \\
		&= (-1)^{pr+qr+pq} \cdot \alpha \enbrace*{\sigma|_{\benbrace*{r,\ldots ,r+p}}} \beta \enbrace*{\sigma|_{\benbrace*{r+q,\ldots ,n}}} \sigma|_{\benbrace*{0,\ldots ,r}}
	\end{align}
	Für die andere Seite der Formel gilt
	\begin{align}
		\alpha \capp (\beta \capp \sigma) &= (-1)^{q(r+p)} \cdot \alpha \capp \enbrace[\Big]{\beta \enbrace*{\sigma|_{\benbrace*{r+q,\ldots ,n}}} \sigma|_{0,\ldots ,r+p}} \\
		&= (-1)^{qr+qp+rp} \cdot \alpha \enbrace*{\sigma|_{\benbrace*{r,\ldots,r+p}}} \beta \enbrace*{\sigma|_{\benbrace*{r+p,\ldots ,n}}} \sigma|_{\benbrace*{0,\ldots ,r}}
	\end{align}
	Damit folgt die Gleichheit.
\end{beweis}

Ab jetzt: $R$ kommutativ!

\begin{erinnerungA}
	Sei $M$ eine $R$-orientierbare zusammenhängende geschlossene $n$-Mannigfaltigkeit. 
	Dann ist $H_n(M;R) \cong R$.\footnote{das sieht man zB. mit der Paarfolge}
	Ein Erzeuger von $H_n(M;R)$ heißt \Index{Fundamentalklasse}.
	Jeder solcher hat die Eigenschaft, dass $\mu|_x \in H_n(M,M \setminus \set*{x};R) \cong R$ ein Erzeuger ist. 
	Die Wahl einer solchen Fundamentalklasse entspricht genau der Wahl einer $R$-Orientierung für $M$.
	Für $R=\mathbb{F}_2$ ist jede $n$-Mannigfaltigkeit $\mathbb{F}_2$-orientierbar. 
	$M$ ist genau dann nicht $R$-orientierbar, wenn $H_n(M;R)=0$.
\end{erinnerungA}

\begin{satz}[{name={Dualitätssatz von Poincaré für geschlossene Mannigfaltigkeiten}},label=poincare]
	Sei $M$ eine geschlossene zusammenhängende $n$-Mannigfaltigkeit und $\mu \in H_n(M;R)$ eine Fundamentalklasse. 
	Dann ist $\capp \mu \colon H^p(M;R) \to H_{n-p}(M;R)$ ein Isomorphismus.
\end{satz}

\begin{bemerkung}
	Zum Beweis des Satzes werden wir eine Überdeckung von $M$ durch offene Mengen, die homöomorph zu $\mathbb{R}^n$ sind, und ein Mayer-Vietoris-Argument benutzen.
	Dazu werden wir die Aussage zunächst so verallgemeinern, dass sie auch für offene Teilmengen des $\mathbb{R}^n$ richtig ist.
\end{bemerkung}

\begin{bemerkung}
	Ist $K$ ein Körper, so induziert die Paarung
	\mapdef{C^*(X;K) \otimes C_*(X;K)}{K}{\alpha \otimes a}{\alpha(a)}{}
	einen Isomorphismus $H^n(X;K) \cong \Hom_K \enbrace*{H_n(X;K),K}$. 
	Dies folgt einerseits aus dem universellen Koeffizienten-Theorem und andererseits aus Aufgabe $2$ von Blatt $2$.
\end{bemerkung}

\begin{beispiel}[{name=[Kohomologie des reell projektiven Raumes]}]
	Es gilt 
	\[
		H^* \enbrace*{\mathbb{R}P^n;\mathbb{F}_2} = \begin{cases}
			\mathbb{F}_2 &\text{ falls }*=0,\ldots ,n\\
			0 &\text{ sonst }
		\end{cases}
	\]
\end{beispiel}

\begin{bemerkung}
	Ist $H_{n-1}(X)=0$, so induziert die Paarung $C^n(X) \otimes C_n(X) \to \mathbb{Z}$, $\alpha \otimes a \mapsto \alpha(a)$ einen Isomorphismus $H^n(X) \cong \Hom_\mathbb{Z} \enbrace*{H_n (X),\mathbb{Z}}$.
	Dies folgt auch wieder entweder aus universellen Koeffizienten-Theorem oder aus Aufgabe $3$ von Blatt $3$.
\end{bemerkung}

\begin{beispiel}[{name=[Kohomologie des komplex projektiven Raumes]}]
	Es gilt
	\[
		H^* \enbrace*{\mathbb{C}P^n;\mathbb{Z}} = \begin{cases}
			\mathbb{Z} &\text{ falls }*=0,2,\ldots ,2n\\
			0 &\text{ sonst}
		\end{cases}
	\]
\end{beispiel}

\begin{satz}[{name=[Ringstruktur der Kohomologie für die projektiven Räume]}]
	Es gibt Isomorphismen von graduierten Ringen: 
	\begin{enumerate}[(i)]
		\item $H^* \enbrace*{\mathbb{R}P^n,\mathbb{F}_2} \cong \Faktor{\mathbb{F}_2[\xi]}{\xi^{n+1}=0}$ mit $\abs*{\xi}=1$.
		\item $H^* \enbrace*{\mathbb{C}P^n,\mathbb{Z}} \cong \Faktor{\mathbb{Z}[\eta]}{\eta^{n+1}=0}$ mit $\abs{\eta}=2$.
	\end{enumerate}
\end{satz}
\begin{beweis}
	Wir wissen schon, dass es solche Isomorphismen von graduierten Moduln gibt. 
	Es bleibt das $\cupp$-Produkt zu verstehen. Für $n=0,1$ ist nichts zu zeigen. 
	Per Induktion nehmen wir an, dass $H^*(\mathbb{R}P^{n-1};\mathbb{F}_2) \cong \sfrac{\mathbb{F}_2[\xi']}{(\xi')^n=0}$, wobei $\xi' := i^*(\xi)$ für $i \colon \mathbb{R}P^{n-1} \to \mathbb{R}P^n$ die Inklusion ist und $\xi \in H^1 \enbrace*{\mathbb{R}P^n;\mathbb{F}_2}$ ein Erzeuger ist.
	Da $i^*$ mit dem $\cupp$-Produkt verträglich ist und für $0 \le * \le {n-1}$ ein Isomorphismus ist, müssen wir nur noch zeigen, dass $\xi^n\in H^n(\mathbb{R}P^n,\mathbb{F}_2)$ ein Erzeuger ist.
	Sei $\mu \in H_n \enbrace*{\mathbb{R}P^n,\mathbb{F}_2}$ eine Fundamentalklasse. Wir müssen zeigen $\xi^n(\mu)=  \pm 1$.
	Nach Induktionsannahme ist $\xi^{n-1}$ ein Erzeuger von $H^{n-1}(\mathbb{R}P^n;\mathbb{F}_2)$. 
	Nach Poincaré-Dualität ist $\xi^{n-1} \capp \mu$ ein Erzeuger von $H_1(\mathbb{R}P^n,\mathbb{F}_2)$. 
	Da $\xi$ ein Erzeuger von $H^1(\mathbb{R}P^n,\mathbb{F}_2)$ ist, ist also $\xi \enbrace*{\xi^{n-1} \capp \mu} = \pm 1$. 
	Daher ist 
	\[
		\xi^n(\mu) = \epsilon \enbrace*{\xi^n \capp \mu} = \epsilon \enbrace*{ \xi \capp \enbrace*{\xi^{n-1} \capp \mu}} = \xi \enbrace*{\xi^{n-1} \capp \mu} = \pm 1 \qedhere
	\]
\end{beweis}

\begin{satz}
	Sei $M$ eine geschlossene Mannigfaltigkeit, $\dim M = n = 2d +1$ ungerade. 
	Sei $M$ weiter homotopieäquivalent zu einen endlichen CW-Komplex.\footnote{Dies ist keine Einschränkung, siehe zB. \cite[Cor. A.12]{Hatcher}; die Umkehrung ist aber nicht richtig, wie Aufgabe 4 von Blatt 5 zeigt} 
	Dann ist $\chi_M =0$.
\end{satz}
\begin{beweis}
	Ohne Einschränkungen können wir annehmen, dass $M$ zusammenhängend ist, da die Euler-Charakteristik additiv ist. Es ist nach Poincaré-Dualität 
	\begin{align}
		\chi_M = \sum_{k=0}^{n} (-1)^k \dim_{\mathbb{F}_2} H_k(M;\mathbb{F}_2) &= \sum_{k=0}^{n} (-1)^k \dim_{\mathbb{F}_2} H^{n-k}(M;\mathbb{F}_2) \\
		\intertext{Da $M$ homotopieäquivalent zu einem CW-Komplex ist, ist $\dim_{\mathbb{F}_2} H_k(M; \mathbb{F}_2)<\infty $ und insbesondere $\dim_{\mathbb{F}_2} H_k(M;\mathbb{F}_2) = \dim \mathbb{F}_2 \Hom_{\mathbb{F}_2} \enbrace*{H_k(M; \mathbb{F}_2),\mathbb{F}_2} = \dim_{\mathbb{F}_2} H^k(M; \mathbb{F}_2)$ ist, folgt}
		&= \sum_{k=0}^{n} (-1)^k \dim_{\mathbb{F}_2} H_{n-k}(M; \mathbb{F}_2) \\
		&= \sum_{k=0}^{n} (-1)^{n-k} \dim_{\mathbb{F}_2} H_k(M; \mathbb{F}_2) \\
		&= - \sum_{k=0}^{n} (-1)^k \dim_{\mathbb{F}_2} H_k(M;\mathbb{F}_2) \\
		&= - \chi_M \qedhere
	\end{align}
\end{beweis}
% section 6 (end)
\newpage

\section{Kohomologie mit kompaktem Träger} % (fold)
\label{sec:7}

\begin{bemerkung}[{name=[Probleme beim Beweis der Poincaré-Dualität]}]
	Am Beweis der Poincaré-Dualtität hindern uns momentan noch folgende Probleme, die wir in diesem Kapitel durch eine Abwandlung von Kohomologie \enquote{umschiffen} wollen: 
	\begin{itemize}[itemsep=0pt]
		\item Für $n>0$ ist  $H^{n-*}(\mathbb{R}^n) \ncong H_*(\mathbb{R}^n)$
		\item Es gibt keine Fundamentalklasse in $H_n(\mathbb{R}^n)$.
	\end{itemize}
\end{bemerkung}

\begin{erinnerungA}[{name=[Orientierungsklassen für kompakte Teilmengen]}]
	Sei $M$ eine (nicht notwendig kompakte) $n$-Mannigfaltigkeit und $K \subseteq M$ eine kompakte Teilmenge. Dann gilt nach Lemma 16.13 aus Topologie \RM{1}.:
	\begin{enumerate}[(i)]
		\item Ist $(\mu_x)_{x \in M}$ eine $R$-Orientierung von $M$, so gibt es eine eindeutige Orientierungsklasse $\mu_K \in H_n(M,M \setminus K;R)$ mit 
		\[
			\mu_K\big|_x = \mu_x \in H_n(M,M \setminus \set*{x};R)
		\]
		für alle $x \in K$.
		\item $H_l(M, M \setminus K;R)=0$ für $l>n$.
	\end{enumerate}
\end{erinnerungA}

\begin{beispiel}[{name=[Fundamentalklasse eines kompakten Balls]}]
	Sei $B=B_r(x_0)$ ein kompakter Ball im $\mathbb{R}^n$. 
	Dann ist $H_n(\mathbb{R}^n,\mathbb{R}^n \setminus B;R) \cong R$ und jeder Erzeuger $\mu_B \in H_n(\mathbb{R}^n,\mathbb{R}^n \setminus B;R)$ ist eine Fundamentalklasse über $B$, also ist jede Einschränkung $\mu_B|_x \in H_n(\mathbb{R}^n,\mathbb{R}^n \setminus \set*{x};R)$ ein Erzeuger für $x \in B$.
\end{beispiel}

\begin{lemma}[{name=[Isomorphismus von Kohomologie und Homologie für kompakte Bälle]},label=Rn_kompakteBaelle]
	Sei $B \subseteq \mathbb{R}^n$ ein kompakter Ball und $\mu_B \in H_n(\mathbb{R}^n,\mathbb{R}^n \setminus B;R)$ ein Erzeuger. Dann ist 
	\[
		\capp \mu_B \colon H^k \enbrace*{\mathbb{R}^n,\mathbb{R}^n \setminus B;R} \longrightarrow H_{n-k}(\mathbb{R}^n;R)
	\]
	ein Isomorphismus.
\end{lemma}
\begin{beweis}
	Es genügt $k=n$ zu betrachten, denn andernfalls sind beide Moduln trivial. Da die Auswertungsabbildung 
	\[
		H^n(\mathbb{R}^n,\mathbb{R}^n \setminus B;R) \longrightarrow \Hom_R \enbrace[\big]{H_n(\mathbb{R}^n,\mathbb{R}^n \setminus B;R),R}
	\]
	immer surjektiv ist nach \autoref{eig_kohomo_to_hom_homo}, gibt es $\xi \in H^n \enbrace*{\mathbb{R}^n,\mathbb{R}^n \setminus B;R}$ mit $\xi(\mu_B)=1$. 
	Also ist $\epsilon(\xi \capp \mu_B)=1$ (siehe \ref{bem_augmentation_cap}). 
	Folglich ist $\xi \capp \mu_B$ ein Erzeuger von $H^0(\mathbb{R}^n;R)$ und $\capp \mu_B$ ein Isomorphismus.
\end{beweis}

\begin{definition}[{name=[gerichtete Menge]}]
	Eine \Index{gerichtete Menge} ist eine Menge $\Lambda$ zusammen mit einer partiellen Ordnung \enquote{$\le$}, sodass es für alle $\lambda,\lambda' \in \Lambda$ ein $\mu \in \Lambda$ existiert mit $\lambda \le \mu $ und $\lambda' \le \mu$.
\end{definition}

\begin{beispiel}[{name=[für gerichtete Mengen]}]
	Es gibt zahlreiche Beispiele für gerichtete Mengen
	\begin{itemize}[itemsep=0pt] 
		\item $(N,\le)$, $(\mathbb{Z},\le)$, $(\mathbb{Q},\le)$ und $(\mathbb{R},\le)$ mit der üblichen Ordnung.
		\item Sei $X$ ein topologischer Raum. Dann ist $\mathcal{K}(X)= \set*{K \subseteq X \given K \text{ kompakt}}$ gerichtet durch \enquote{$\subseteq$}.
	\end{itemize}
\end{beispiel}

\begin{bemerkung}[{name=[Gerichtete Mengen als kleine Kategorie]}]
	Wir können jede gerichtete Menge $(\Lambda,\le)$ als eine kleine Kategorie auffassen: 
	Die Menge der Objekte ist $\Lambda$ und es gibt genau dann einen Morphismus $i_\lambda^{\lambda'} \colon \lambda \to \lambda'$, falls $\lambda \le \lambda'$ ist. 
	Andernfalls gibt es keine Morphismen.
\end{bemerkung}

\begin{definition}[{name=[gerichteter Kolimes]}]
	Sei $(\Lambda,\le)$ eine gerichtete Menge und $F \colon (\Lambda,\le) \to R\text{-}\MOD$ ein Funktor. Der \bet{gerichtete Kolimes}\index{gerichteter Kolimes} von $F$ ist dann definiert als 
	\[
		\colim_\Lambda F \coloneqq \faktor{\coprod_{\lambda \in \Lambda}F(\lambda)}{ \begin{array}{l}  \scriptstyle v \sim F \enbrace{i_\lambda^{\lambda'}}(v), \\[-3pt]  \scriptstyle\text{für } v \in F(\lambda), \lambda \le \lambda' \end{array}}
	\]
	Für $\lambda_0 \in \Lambda$ sei $s_{\lambda_0} \colon F(\lambda) \to \colim_\Lambda F$ die Komposition von der Inklusion $F(\lambda_0) \hookrightarrow \coprod_{\lambda \in \Lambda}F(\lambda)$ mit der Projektion $\coprod_{\lambda \in \Lambda}F(\lambda) \twoheadrightarrow \colim_\Lambda F$. Durch
	\begin{align}
		r \cdot s_\lambda(v) &\coloneqq s_\lambda(r \cdot v) \quad \text{ und} \\
		s_\lambda(v) + s_{\lambda'}(v') &\coloneqq s_\mu \enbrace*{F(i_\lambda^{\mu})(v) + F(i_{\lambda'}^\mu)(v')} 
	\end{align}
	für $\lambda,\lambda' \le \mu$ erhalten wir eine $R$-Modulstruktur auf $\colim_\Lambda F$. Die Abbildungen sind dann $R$-linear.
\end{definition}

\begin{bemerkung}
	\leavevmode
	\begin{itemize}
		\item Nicht jede kleine Kategorie gehört zu einer gerichteten Menge:
		\begin{enumerate}[(i)]
			\item Die von der Gruppe $\sfrac{\mathbb{Z}}{2}$ erzeugte Kategorie mit einem Objekt und den Gruppenelementen als Morphismen.
			\item Die Kategorie bestehend aus zwei Objekten und keinem Morphismus, der nicht eine Identität ist.
		\end{enumerate}
		\item Der Kolimes über gerichtete Mengen ist eine Spezialfall des Kolimes über kleine Kategorien aus den Übungen. 
		Insbesondere hat der Kolimes über gerichtete Mengen einen universelle Eigenschaft, die ihn charakterisiert.
		\item Sei $v_\lambda \in F(\lambda)$. Dann gilt
		\(
			s_\lambda(v_\lambda) =0 \in \colim_\Lambda F \iff \exists \lambda' : \lambda \le \lambda', F(i_\lambda^{\lambda'})(v_\lambda)=0 \in F(\lambda')
		\).
	\end{itemize}
\end{bemerkung}

\begin{definition}[{name=[kofinal]}]
	Sei $(\Lambda,\le)$ eine gerichtete Menge. $\Lambda_0 \subseteq \Lambda$ heißt \Index{kofinal}, falls es zu jedem $\lambda \in \Lambda$ ein $\lambda_0 \in \Lambda_0$ gibt mit $\lambda\le \lambda_0$.
\end{definition}

\begin{lemma}[{name=[Isomorphismus zum Kolimes einer kofinalen Teilmenge]}]
	Sei $\Lambda_0 \subseteq \Lambda$ kofinal. Dann ist die kanonische Abbildung (induziert von der Inklusion $\coprod_{\lambda_0} F(\lambda_0) \hookrightarrow \coprod_{\lambda} F(\lambda)$)
	\[
		c \colon \colim_{\lambda_0 \in \Lambda_0} F(\lambda_0) \longrightarrow \colim_{\lambda \in \Lambda} F(\lambda)
	\]
	ein Isomorphismus.
\end{lemma}
\begin{beweis}
	Surjektivität: Sei $v \in \colim_\Lambda F(\lambda)$. Wähle $\lambda \in \Lambda$, $v_\lambda \in F(\lambda)$ mit $s_\lambda(v_\lambda)=v$. 
	Da $\Lambda_0 \subseteq \Lambda$ kofinal ist, gibt es $\lambda_0 \in \Lambda_0$ mit $\lambda \le \lambda_0$. 
	Dann ist
	\[
		v = s_{\lambda_0} \enbrace*{F(i_\lambda^{\lambda_0})(v_\lambda)} \in \im c
	\]
	Injektivität: Sei $v \in \ker c$. Dann existieren $\lambda_0 \in \Lambda_0$, $v_{\lambda_0} \in F(\lambda_0)$ mit $v=s_{\lambda_0}(v_{\lambda_0})$. 
	Da $v \in \ker c$ ist, gibt es $\lambda \in \Lambda$ mit $\lambda_0 \le \lambda$ und $F(i_{\lambda_0}^{\lambda})(v_{\lambda_0}) =0$. 
	Da $\Lambda_0 \subseteq \Lambda$ kofinal ist, gibt es $\lambda_0' \in \Lambda_0$ mit $\lambda \le \lambda_0$. 
	Es folgt 
	\begin{align}
		F(i_{\lambda_0}^{\lambda_0'})(v_{\lambda_0}) = F (i_\lambda^{\lambda_0'} \circ i_{\lambda_0}^\lambda)(v_{\lambda_0}) = F(i_{\lambda}^{\lambda_0'}) \enbrace[\Big]{\Underbrace{F(i_{\lambda_0}^\lambda)(v_{\lambda_0})}{=0}} = 0
	\end{align}
	Es folgt $v=0$ und damit die Behauptung.
\end{beweis}

\begin{korollarB}
	Gibt es $\lambda_\infty \in \Lambda$ mit $\lambda \le \lambda_\infty$ für alle $\lambda \in \Lambda$, so ist $\colim_\Lambda F = F(\lambda_\infty)$.
\end{korollarB}

\begin{beispiel}[{name=[Kolimiten]}]\leavevmode
	\begin{itemize}
		\item Sei $\mathcal{P}\mkern-2.5mu_{\mathrm{endl}}$ die Menge aller endlichen Teilmengen der Menge $S$. 
		Bezüglich $\subseteq$ ist $\mathcal{P}\mkern-2.5mu_{\mathrm{endl}}$ gerichtet.
		Sei $F \colon \enbrace*{\mathcal{P}\mkern-2.5mu_{\mathrm{endl}},\subseteq} \to R\text{-}\MOD$ definiert durch $F(S_0)=R[S_0]$ für $S_0 \subseteq S$ endlich. 
		Dann ist $\colim_{\mathcal{P}\mkern-2.5mu_{\mathrm{endl}}} F \cong R[S]$ und $s_{S_0} \colon R[S_0]\to R(S)$ ist die von $S_0 \subseteq S$ induzierte Abbildung.
		\item Sei $F \colon \mathbb{N} \to \mathbb{Z}\text{-}\MOD$ definiert durch $F(n)=\mathbb{Z}$ für alle $n$ und $F(i_n^{n'}) \colon \mathbb{Z} \to \mathbb{Z}$ ist Multiplikation mit $2^{n'-n}$. 
		Es ergibt sich also folgendes Bild:
		\[
			\begin{tikzcd}
				\mathbb{Z} \rar["2"] & \mathbb{Z} \rar["2"] & \mathbb{Z} \rar["2"] &\mathbb{Z} \rar["2"] & \mathbb{Z} \rar["2"] & \ldots 
			\end{tikzcd}
		\]
		Dann liefert $(n,z) \mapsto \frac{z}{2^n}$ eine Abbildung
		\[
			\colim F = \Faktor{\coprod_{\mathbb{N}} \mathbb{Z}}{\sim} = \Faktor{\mathbb{N}\times \mathbb{Z}}{\scriptstyle(n,z) \sim (n',2^{n'-n} z), n < n'} \grenzw{} \mathbb{Z}\benbrace*{\sfrac{1}{2}} = \set*{\sfrac{a}{2^n}\given a \in \mathbb{Z}, n \in \mathbb{N} }
		\]
		\item Ist $F$ konstant -- also $F(i_\lambda^{\lambda'})=\id_{F(\lambda)}$ für alle $\lambda \le \lambda'$ -- so ist 
		\(
			\colim_\Lambda F \cong F(\lambda)
		\)
		für jedes $\lambda \in \Lambda$.
		\item Sei $F(n) := R \benbrace[\big]{\set*{n,n+1,n+2,\ldots }}$ und $F(i_n^{n'}) \colon R \benbrace[\big]{\set*{n,n+1,\ldots }} \to R \benbrace[\big]{\set*{n',n'+1,\ldots }}$ definiert durch 
		\[
			\sum_{k=n}^{\infty} r_k \cdot k \longmapsto \sum_{k=n'}^{\infty} r_k \cdot k
		\]
		Dann ist $\colim_\mathbb{N} F = 0$, denn für jedes $v \in F(n)$ gibt es $n'> n$ mit $F(i_n^{n'})(v)=0$.
	\end{itemize}
\end{beispiel}

\begin{bemerkung}[{name=[gerichteter Limes als Funktor]}]
	Sei $\FUN\enbrace*{(\Lambda,\le), R\text{-}\MOD}$ die Kategorie, deren Objekte Funktoren $F \colon (\Lambda,\le) \to R \text{-}\MOD$ sind und deren Morphismen natürliche Transformationen $\tau \colon F \to G$ sind. Dann ist
	\[
		\colim_\Lambda \colon \FUN\enbrace*{(\Lambda,\le), R\text{-}\MOD} \longrightarrow R\text{-}\MOD
	\]
	ein Funktor: Ist $\tau \colon F \to G$ eine natürliche Transformation, so induziert $\coprod_{\lambda \in \Lambda}\tau_\lambda \colon \coprod_{\lambda \in \Lambda} F(\lambda) \to \coprod_{\lambda \in \Lambda} G(\lambda)$ eine $R$-lineare Abbildung $\tau_* =\colim_\Lambda \tau \colon \colim_\Lambda F \to \colim_\Lambda G$
\end{bemerkung}

\begin{definition}[{name=[Exaktheit in der Funktorkategorie]}]
	Eine Folge 
	\( \begin{tikzcd}[cramped]
		F \rar["\tau"] & F' \rar["\tau'"] & F''
	\end{tikzcd} \) 
	in der Funktorkategorie $\FUN$ heißt \Index{exakt}, wenn 
	\( \begin{tikzcd}[cramped]
		F(\lambda)\rar["\tau_\lambda"] & F'(\lambda) \rar["\tau_\lambda'"] & F''(\lambda)
	\end{tikzcd} \) 
	exakt ist.
\end{definition}


\begin{satz}[name={Exaktheit gerichteter Kolimiten}]
	Sei $(\Lambda,\le)$ eine gerichtete Menge. Ist 
	\[
		\begin{tikzcd}
			0 \rar & F \rar["\tau"] & F' \rar["\tau'"] & F'' \rar & 0
		\end{tikzcd}
	\]
	exakt in $\FUN \enbrace*{(\Lambda,\le), R\text{-}\MOD}$ so ist auch folgende Sequenz exakt
	\[
		\begin{tikzcd}[sep=large]
			0 \rar & \colim_\Lambda F \rar["\colim_\Lambda \tau"] & \colim_\Lambda F' \rar["\colim_\Lambda \tau'"] & \colim_\Lambda F'' \rar & 0
		\end{tikzcd}
	\]
\end{satz}
\begin{beweis}
	Seien $s_\lambda, s'_\lambda, s''_\lambda$ die Strukturabbildungen. 
	Wir zeigen nur, dass $\tau_* =\colim_\Lambda \tau$ injektiv ist. 
	Für $v \in \ker (\tau_*)$ wähle $v_\lambda \in F(\lambda)$ mit $s_\lambda(v_\lambda)=v$.
	Dann ist $\tau_*(v)=s_\lambda'\enbrace[\big]{\tau_\lambda \enbrace*{v_\lambda}}$.
	Da $\tau_*(v)=0$ ist, gibt es $\mu \ge \lambda$ mit 
	\[
		F'(i_\lambda^\mu) \enbrace[\Big]{s_\lambda'\enbrace[\big]{\tau_\lambda \enbrace*{v_\lambda}}}=0
	\]
	Da $\tau$ eine natürliche Transformation ist, ist 
	\begin{align}
		0=F'(i_\lambda^\mu) \enbrace[\big]{ \tau_\lambda(v_\lambda)} = \tau_\mu \enbrace*{F(i_\lambda^\mu)(v_\lambda)}
	\end{align}
	Da $\tau_\mu$ injektiv ist, folgt also $F(i_\lambda^\mu)(v_\lambda)=0$ und damit ist $v=s_\lambda(v_\lambda)=0$.
\end{beweis}

\begin{bemerkung}[{name=[gerichteter Kolimes von Kern und Kokern]}]
	Für eine natürliche Transformation $\tau \colon F \to F'$ in $\FUN \enbrace*{(\Lambda,\le),R \text{-}\MOD}$ können wir $\ker \tau$ und $\coker \tau \in \FUN \enbrace*{(\Lambda,\le), R \text{-}\MOD}$ erklären durch $\ker \tau (\lambda) = \ker \tau_\lambda$ und $\coker \tau(\lambda) = \coker \tau_\lambda$.
	Mit der Exaktheit vom $\colim_\Lambda$ folgt insbesondere\todo{Revision 7}
	\[
		\colim_\Lambda \ker \tau = \ker \colim_\Lambda \tau \qquad \text{und} \qquad \colim_\Lambda \coker \tau = \coker \colim_\Lambda \tau
	\]
\end{bemerkung}

\begin{korollar}[{name=[Homologie vertauscht mit Kolimiten]}]
	Sei $(\Lambda,\le)$ gerichtet und $C_*$ ein Funktor von $(\Lambda,\le)$ in die Kategorie der Kettenkomplexe und Kettenabbildungen. 
	Dann ist $H_*(\colim_\Lambda C_*) = \colim_\Lambda H_*(C_*)$. 
	Die gleiche Aussage gilt für Kokettenkomplexe.
\end{korollar}
\begin{beweis}
	Sei $d^{\lambda}_n \colon C_n(\lambda) \to C_{n-1}(\lambda)$ die Randabbildung von $C_*(\lambda)$.
	Wir fassen $d_n$ als natürliche Transformation $C_n \to C_{n-1}$ in $\FUN \enbrace*{(\Lambda,\le),R\text{-}\MOD}$ auf. 
	Da $(\Lambda,\le)$ gerichtet ist und daher $\colim_\Lambda$ exakt ist, folgen
	\[
		\ker \enbrace*{d_* \colon \colim_{\Lambda} C_n \to \colim_\Lambda C_{n-1}} \cong \colim_\Lambda \enbrace*{\ker d_n \colon C_n \to C_{n-1} }
	\]
	und 
	\[
		\Underbrace{\coker \enbrace*{ d_{*} \colon \colim_\Lambda C_{n+1} \to \colim_\Lambda \ker d_n}}{ \cong H_n(\colim C_*)}
		\cong \colim_\Lambda \enbrace[\Big]{ \Underbrace{\coker d_{n+1} \colon C_{n+1} \to \ker d_n}{\cong H_n(C_*)}} \qedhere
	\]
\end{beweis}

\begin{definition}[{name=[singuläre Kohomologie mit kompaktem Träger]}]
	Sei $X$ ein topologischer Raum. Dann definieren wir die \bet{singuläre Kohomologie von $X$ mit kompaktem Träger}\index{Kohomologie mit kompaktem Träger} durch
	\[
		H^*_c(X;R) := \colim_{K \in \mathcal{K}(X)} H^*(X,X \setminus K;R)
	\]
\end{definition}

\begin{beispiel}[{name=[Kohomologie mit kompaktem Träger des des euklidischen Raumes]}]
	Es gilt $H^k_c(\mathbb{R}^n;R) \cong \mathbb{R}$ für $k=0$ und $H^k_c(\mathbb{R}^n;R)=0$ sonst. 
\end{beispiel}
\begin{beweis}
	Die Menge $\mathcal{B}$ der kompakten Bälle $B_r(0) \subset \mathbb{R}^n$ ist kofinal in der gerichteten Menge aller kompakten Unterräume von $\mathbb{R}^n$.
	Daher ist nach \autoref{Rn_kompakteBaelle}
	\[
		H^*_c(\mathbb{R}^n;R) = \colim_{K \in \mathcal{K}(\mathbb{R}^n)} H^* \enbrace*{\mathbb{R}^n , \mathbb{R}^n \setminus K;R} 
		\cong \colim_{B \in \mathcal{B}} H^* \enbrace*{\mathbb{R}^n,\mathbb{R}^n \setminus B;R} \cong \begin{cases}
			R &\text{ falls }*=n\\
			0 &\text{ sonst }
		\end{cases} \qedhere
	\]
\end{beweis}

\begin{beispiel}
	Ist $X$ kompakt, so ist $H^*_c(X)= H^*(X)$, denn für $X$ kompakt ist $X \in \mathcal{K}(X)$ ein maximal, also $\colim_{K \in \mathcal{K}(X)} H^*(X,X \setminus K) = H^*(X,X\setminus X) = H^*(X)$.
\end{beispiel}

\begin{definition}
	Setze $C_c^*(X) := \colim_{\mathcal{K}(X)} C^*(X,X\setminus K)$
\end{definition}

\begin{korollarB}
	Es gilt $H^*_c(X) = H^*( C^*_c(X))$, da $\colim$ mit $H^*$ vertauscht.
\end{korollarB}

\begin{lemma}
	$C^*_c(X)$ lässt sich mit dem Unterkomplex von $C^*(X)$ identifizieren, der aus allen $\alpha \colon C*(X) \to R$ besteht, für die es $K \subseteq X$ kompakt gibt mit $\alpha(\sigma)=0$ für $\sigma \colon \abs*{\Delta^n} \to X\setminus K$.
\end{lemma}
\begin{beweis}
	Für $K \subseteq L$ ist $C^*(X,X\setminus K)\to C^*(X,X\setminus L)$ injektiv.
	Nachdem wir jedes $C^*(X,X\setminus K)$ mit einem Unterraum von $C^*(X)$ identifizieren, wird der Kolimes über $\mathcal{K}(X)$ einfach zur Vereinigung dieser Unterräume.
\end{beweis}

\begin{bemerkung}
	Sei $M $ eine $n$-dimensionale Mannigfaltigkeit mit einer $R$-Orientierung. Für $K \subseteq M$ kompakt sei $\mu_K \in H_n(M,M \setminus K;R)$ die durch die $R$-Orientierung bestimmte (Orientierungs)-Klasse.
	Für $K \subseteq L$, $L$ auch kompakt ist dann $\mu_L|_K = \mu_K$.
	Daher kommutiert
	\[
		\begin{tikzcd}
			H^k(M,M \setminus K;R)\dar \rar["\capp\mu_K"] &  H_{n-k}(M;R) \dar[equal] \\
			H^k(M,M \setminus L;R) \rar["\capp\mu_L"] & H_{n-k}(M;R)
		\end{tikzcd}
	\]
	Daher induzieren die $\capp$-Produkte mit den $\mu_K$ eine Abbildung 
	\[
		\PD_M \colon H^k_c(M;R) \longrightarrow H_{n-k}(M;R)
	\]
\end{bemerkung}

\begin{definition}
	$\PD_M$ heißt die \Index{Poincaré-Dualitäts-Abbildung} für $M$. 
\end{definition}

\begin{bemerkung}
	Ist $M$ kompakt, so ist $\PD_M= {\capp \mu}$, wobei $\mu$ die durch die Orientierung bestimmte Fundamentalklasse von $M$ ist.
\end{bemerkung}

\begin{lemma}
	Für $M=\mathbb{R}^n$ ist die Poincaré-Dualitäts-Abbildung ein Isomorphismus. 
\end{lemma}
\begin{beweis}
	Ist $B \subset \mathbb{R}^n$ ein kompakter Ball, so ist ${\capp\mu_B} \colon H^k(\mathbb{R}^n,\mathbb{R}^n \setminus B,R) \to H_{n-k}(\mathbb{R}^n)$ ein Isomorphismus.
	Da die Menge der kompakten Bälle $\mathcal{B}$ kofinal in allen kompakten Teilmengen ist, ist auch
	\(
		\PD_M = \colim_{\mathcal{B}} {\capp \mu_B}
	\)
	ein Isomorphismus.
\end{beweis}

\begin{bemerkung}
	Für nicht kompakte orientierte Mannigfaltigkeiten lässt sich die Fundamentalklasse in lokal endlicher Homologie definieren: 
	\[
		C^\lf_{n}(M) = \set*{\sum_{\sigma\in S_n(M)} r_\sigma \cdot \sigma \given \begin{array}{l}
			\text{für jedes $K \subseteq M$ kompakt gibt es nur endlich viele $\sigma \colon \abs*{\Delta^n} \to M$} \\
			\text{mit $r_\sigma\neq 0$ und $\sigma(\abs*{\Delta^n}) \cap K \neq \emptyset$} 
		\end{array}
		  }
	\]
\end{bemerkung}
% section 7 (end)
\newpage

\section{Poincaré-Dualität} % (fold)
\label{sec:8}

\begin{lemma}
	Seien $U,V \subseteq X$ offen. ei $C^*(X,U+V)$ der Unterkomplex von $C^*(X)$, der aus allen $\varphi \colon C_*(X)\to R$ besteht mit $\varphi(\sigma)=0$ für $\sigma \colon \abs*{\Delta^n}\to X$ mit $\im (\sigma) \subseteq U$ \emph{oder} $\im (\sigma) \subseteq V$. 
	Dann induziert die Inklusion $C^*(X, U\cup V) \to C^*(X,U+v)$ einen Isomorphismus in Kohomologie.
\end{lemma}
\begin{beweis}
	Sei $C^*(U+V) := C_*^{\set*{U,V}}(U \cup V)$. 
	Die Inklusion $C_*(U +V)\to C_*(U \cup V)$ ist nach dem Satz über kleine Simplizes eine Kettenhomotopieäquivalenz.
	Sei $C^*(U+V)= \Hom \enbrace*{C_*(U+V),R}$ der zu $C_*(U+V)$ duale Kokettenkomplex. 
	Dann ist auch $C^*(U \cup V) \to C^*(U+V)$ eine Kokettenhomotopieäquivalenz.
	Betrachte nun
	\[
		\begin{tikzcd}
			0 \rar & C^*(X,U\cup V) \dar  \rar & C^*(X)\dar[equal] \rar & C^*(U \cup V) \dar \rar & 0 \\
			0 \rar & C^*(X,U+V) \rar & C^*(X) \rar & C^*(U+V) \rar & 0
		\end{tikzcd}
	\]
	Die zweite und dritte Abbildung induzieren jeweils Isomorphismen in Kohomologie. 
	Da das Diagramm kommutiert erhalten wir eine lange exakte Leiter mit den langen exakten Folgen in Kohomologie.
	Das 5er-Lemma impliziert nun, dass auch die erste Abbildungen einen Isomorphismus in Kohomologie induziert.
\end{beweis}

\begin{satz}
	Seien $U,V \subseteq X$ offen. Dann gibt es eine lange exakte Mayer-Vietoris-Folge 
	\[
		\begin{tikzcd}[sep=1.55em]
			\ldots \rar & H^n(X,U+V) \rar & H^n(X,U) \oplus H^n(X,V) \rar & H^n(X,U\cap V) \rar["\partial"] & H^{n+1}(X,U\cup V) \rar & \ldots 
		\end{tikzcd}
	\]
	Der Verbindungshomomorphismus $\partial$ kann wie folgt beschrieben werden: Sei $\alpha \in C^n(X,U \cap V)$ mit $d^n(\alpha)=0$.
	Dann lässt sich $\alpha$ schreiben als $\alpha= \alpha_U + \alpha_V$ mit $\alpha_U \in C^*(X,U)$ und $\alpha_V \in C^n(X,V)$ und es ist $\partial [\alpha] = \benbrace*{d^n \alpha_U}= \benbrace*{d^n \alpha_V}$
\end{satz}
\begin{beweis}
	Benutze die kurze exakte Folge von Kokettenkomplexen: 
	\[
		\begin{tikzcd}
			0 \rar & C^*(X,U + V) \rar & C^*(X,U) \oplus C^*(X,V) \rar & C^*(X, U \cup V) \rar & 0
		\end{tikzcd}
	\]
	wobei die erste Abbildung $\alpha \mapsto (\alpha,\alpha)$ und die zweite $(\alpha,\beta) \mapsto \alpha-\beta$ ist. 
	
	
	Es ist $\alpha+\beta=0$ in $C^*(X,U \cap V) \subseteq C^*(X)$ für $\alpha \in C^*(X;U) \subseteq C^*(X)$, $\beta \in C^*(X,V) \subseteq C^*(X)$. Also ist
	$\alpha(\sigma)= - \beta(\sigma)$ für alle $\sigma \subseteq X$.
	
	Die zugehörigen lange exakte Folge ist die gesuchte Mayer-Vietoris-Folge. 
\end{beweis}

\begin{satz}
	Seien $U,V \subseteq X$ offen. Sei $X$ lokalkompakt und Hausdorffsch. 
	Dann gibt e seine lange exakte Mayer-Vietoris-Folge 
	\[
		\begin{tikzcd}
			\ldots \rar & H^n_c(U\cap V) \rar & H^n_c(U) \oplus H^n_c(V) \rar & H^n_c(U \cup V) \rar["\partial"] & H^{n+1}_c(U\cap V) \rar & \ldots 
		\end{tikzcd}
	\]
\end{satz}
\begin{beweis}
	Seien $K \subseteq U$, $L \subseteq V$ kompakt. Es ist 
	\begin{align}
		(X\setminus K) \cup (X\setminus L) &= X \setminus (K \cap L) \\
		(X \setminus K) \cap (X \setminus L) &= X \setminus (K \cup L)
	\end{align}
	Wir erhalten eine lange exakte Folge 
	\[
		\begin{tikzcd}[sep=1.9em]
			\ldots \rar & H^n \enbrace*{X,X \setminus (K \cap L)} \rar & H^n \enbrace*{X,X \setminus K} \oplus H^n(X,X \setminus L) \rar & H^n \enbrace*{X,X \setminus (K \cup L)} \rar & \ldots 
		\end{tikzcd}
	\]
	Mit Ausschneidung ergibt sich eine lange exakte Folge 
	\[
		\begin{tikzcd}[sep=small,font=\small]
			\ldots  \rar & H^n \enbrace*{U \cap V, U  \cap V  \setminus K \cap L} \rar & H^k(U,U \setminus K) \oplus H^k(C,V \setminus L) \rar & H^k(U \cup V, U \cup V \setminus K \cup L) \rar & \ldots 
		\end{tikzcd}
	\]
	Sei $\mathcal{K}$ die gerichtete Menge der kompakten Unterräume von $U$, $\mathcal{L}$ die von $V$.
	Sei $\Lambda = \mathcal{K} \times \mathcal{L}$ das Produkt der gerichteten Mengen $\mathcal{K}$ und $\mathcal{L}$.
	Da $\colim_\Lambda$ exakt ist, erhalten wir auch im Kolimes eine lange exakte Folge. Dies ist die gesuchte lange exakte Folge. 
	(Da $X$ lokalkompakt und Hausdorffsch ist, ist jeder kompakter Teilraum $M$ von $U \cup V$ die Vereinigung von $M=K \cup L$ mit $(K,L) \in \Lambda$.)
\end{beweis}

\begin{erinnerung}
	Seien $U,V \subseteq X$ offen. Dann ist $H_*(U \cup V) = H_* \enbrace*{C_*(U+V)}$ und es gibt eine lange exakte Folge 
	\[
		\begin{tikzcd}
			\ldots  \rar & H_n(U \cap V) \rar & H_n(U) \oplus H_n(C) \rar & H_n(U \cup V) \rar["\partial"]& \ldots 
		\end{tikzcd}
	\]
	Der Verbindungshomomorphismus kann wie folgt beschrieben werden: Sei $x \in C_n(U + V)$ mit $d_n x =0$.
	Dann ist $x=x_U - x_V$ mit $x_U \in C_n(U)$, $x_V \in C_n(V)$ und es ist $\partial[x] = \benbrace*{d_n x_U} = \benbrace*{d_n x_V}$.
\end{erinnerung}
\begin{beweis}
	Benutze 
	\[
		\begin{tikzcd}
			0 \rar & C_*(U \cap V) \rar & C_*(U) \oplus C_*(V) \rar & C_*(U+V) \rar & 0
		\end{tikzcd} \qedhere
	\]
\end{beweis}

\begin{proposition}
	Seien $U,V$ offene Teilmengen einer $R$-orientierten $n$-Mannigfaltigkeit $M$. 
	Dann bilden die Poincaré-Dualitätsabbildungen zusammen mit der Mayer-Vietorisfolge ein kommutatives Diagramm
	\[
		\begin{tikzcd}
			\ldots \rar & H^k_c(U \cap V) \rar \dar["\PD_{U \cap V}"] & H^k_c(U) \oplus H^k_c(V) \dar["\PD_{U} \oplus \PD_{V}"] \rar & H^k_c(U \cup V) \rar["\partial"] \dar["\PD_{U \cup V}"]& \ldots \dar \\
			\ldots \rar & H_{n-k}(U \cap V) \rar & H_{n-kk}(U) \oplus H_{n-kk}(V) \rar & H_{n-k}(U \cup V) \rar & \ldots 
		\end{tikzcd}
	\]
\end{proposition}
\begin{beweis}
	Ohne Einschränkungen sei $M=U \cup V$. 
	Die Kommutativität der ersten beiden folgt aus der Natürlichkeit des $\capp$-Produktes.
	Da $\partial \colon H^k_c(M) \to H_c^{k+1}(U \cap V)$ der Kolimes der 
	\[
		H^k(M,M \setminus (K \cup L)) \grenzw{\partial} H^{k+1}(M,M \setminus (K \cap L)) 
	\]
	für $K \subseteq U, L \subseteq V$ kompakt ist, genügt es für die Kommutativität des dritten Quadrats die Kommutativität von folgendem Diagramm nachzurechnen.
	\[
		\begin{tikzcd}[row sep=large,column sep=1.6em]
			H^k \enbrace[\big]{M , (M \setminus K) \cap (M \setminus L)} \rar["\partial"] \dar[equal] & H^{k+1} \enbrace[\big]{M,(M \setminus K) \cup (M \setminus L)} \dar[equal]\\
			H^k \enbrace[\big]{M , M \setminus (K \cup L)} \dar["\capp \mu_{K \cup L}"] \rar["\partial"] & H^{k+1} \enbrace[\big]{M, M \setminus (K \cap L)} \rar["\cong"] & H^{k+1} \enbrace[\big]{U \cap V, U \cap V \setminus (K \cap L)} \dar["\capp \mu_{K \cap L}"] \\
			H_{n-k}(M) \ar[rr,"\partial"] & & H_{n-k-1}(U \cap V)
		\end{tikzcd}
	\]
	Sei $\alpha \in C^{k} \enbrace*{M,M \setminus (K \cup L)}$ mit $d^k\alpha=0$, $\alpha=\alpha_K-\alpha_L$ mit $\alpha_K \in C^k \enbrace*{M,M \setminus K}$ und $\alpha_L \in C^k(M,M \setminus L)$.
	Dann ist $\partial [\alpha]= \benbrace*{d^k \alpha_K} = \benbrace*{d^k \alpha_L}$.\marginnote{$U \cup V = U \setminus L  \cup V \setminus K \cup U \cap V$}
	Wegen Ausschneidung finden wir $x_K \in C_n(V\setminus K)$, $x_L \in C_n(U \setminus L)$ und $x_\cap C_n(U \cap V)$ mit $\mu_{K \cup L} = \benbrace*{x_\cap +x_K+x_L}$.
	Es ist dann auch $\mu_{K \cap L} = \benbrace*{x_\cap}$ und daher ist das Bild von $\partial[\alpha]$ in $H_{n-k-1}(U \cap V)$ gegeben durch $\benbrace*{d^k \alpha_K \capp x_\cap}$.
	Es ist andererseits 
	\begin{align}
		[\alpha] \cap \mu_{K \cup L} &= \benbrace[\big]{(\alpha_K -\alpha_L) \capp (x_\cap + x_K+x_L)} \\
		&= \benbrace[\Big]{\alpha_K \capp x_\cap  + \Underbracket{\alpha_K \capp x_K}{=0} + \alpha_K \capp x_L - \alpha_L \capp x_\cap - \alpha_L \capp x_K - \Underbracket{\alpha_L \capp x_L}{=0}} \\
		&= \benbrace[\Big]{\Underbracket{\alpha_K \capp x_L}{\in C_k(U)} + \enbrace[\big]{\Underbracket{\alpha_K \capp x_\cap - \alpha_L \capp x_\cap}{\in C_K(U \cap V)}} - \Underbracket{\alpha_L \capp x_K}{\in C_k(V)}} 
	\end{align}
	Damit haben wir 
	\(
		\partial \enbrace*{[\alpha] \capp \mu_{K \cup L}} = \benbrace[\big]{d_{n-k} \enbrace*{\alpha_K \capp x_L}} = \benbrace[\Big]{\Underbracket{d^k \alpha_K \capp x_L}{=+d^k \alpha_L \capp x_L=0} + (-1)^k \alpha_K \capp d_n x_L}
	\).
	Mit $d_n x_L = - d_n x_K - d_n x_\cap$ folgt 
	\begin{align}
		\partial \enbrace*{[\alpha]\capp \mu_{K \cup L}} &= (-1)^{k+1} \benbrace[\Big]{\Underbracket{\alpha_K \capp d_n x_K}{=0} + \alpha_K \capp d_n x_\cap} \\
		&= (-1)^{k+1} \benbrace[\big]{\alpha_K \capp d_n x_\cap}
	\end{align}
	Wegen $d_{n-k}(\alpha_K \capp x_\cap)= d^k \alpha_K \capp x_\cap + (-1)^k \alpha_K \capp d_n x_\cap$ folgt $\benbrace*{\alpha_K \capp d_n x_\cap} = (-1)^{k+1} \benbrace*{d^k \alpha_K  \cap x_{\cap}}$.
	Also ist $\partial \enbrace*{[\alpha] \capp \mu_{K \cup L}} = \benbrace*{d^k  \alpha_K \capp x_\cap}$.
\end{beweis}

\begin{satz}
	Sei $M$ eine $R$-orientierbare $n$-Mannigfaltigkeit. 
	Dann ist $\PD_M \colon H^k_c(M) \to H_{n-k}(M)$ ein Isomorphismus.
	Insbesondere ist, falls $M$ sogar geschlossen ist, $\capp \mu_M \colon H^k(M) \to H_{n-k}(M)$ ein Isomorphismus.
\end{satz}
\begin{beweis}
	Für $U \cong \mathbb{R}^n$ haben wir die Behauptung schon bewiesen.\todo{Referenz}
	Da offene konvexe Unterräume im $\mathbb{R}^n$ homöomorph zum $\mathbb{R}^n$ sind, ist die Behauptung auch für solche richtig.
	Schnitte offener konvexer Teilräume des $\mathbb{R}^n$ wieder offen und konvex sind, folgt mit den Mayer-Vietoris-Folgen die Behauptung für endliche Vereinigungen offener und konvexer Unterräume des $\mathbb{R}^n$.
	Mit \autoref{lem:abzUeber} folgt die Behauptung auch für abzählbare Vereinigungen von offenen konvexen Unterräumen im $\mathbb{R}^n$.
	Da jeder offene Unterraum des $\mathbb{R}^n$ die abzählbare Vereinigung offener konvexer Unterräume ist, folgt die Behauptung für beliebige Kartengebiete\footnote{$U \subseteq M$ mit $\varphi \colon U \grenzw{\cong} V \subseteq \mathbb{R}^n$ und $V \subseteq \mathbb{R}^n$ offen} in Mannigfaltigkeiten.
	Mit der Mayer-Vietoris-Folge folgt die Behauptung für endliche Vereinigungen von Kartengebieten in Mannigfaltigkeiten und mit \autoref{lem:abzUeber} für abzählbare Vereinigungen von Kartengebieten.
	Da jede Mannigfaltigkeit die abzählbare Vereinigung von Kartengebieten ist, folgt die Behauptung.
\end{beweis}

\begin{lemma}[label=lem:abzUeber]
	Sei $M$ eine $R$-orientierte $n$-Mannigfaltigkeit. 
	Seien $U_1 \subseteq U_2 \subseteq \ldots \subseteq M$ offen mit $M= \bigcup_{i=1}^\infty U_i$.
	Sind alle $\PD_{U_i}$ Isomorphismen, so ist auch $\PD_M$ ein Isomorphismus.
\end{lemma}
\begin{beweis}
	Jeder singuläre Simplex in $M$ liegt schon in einem $U_i$. 
	Das heißt $C_*(M)=\colim_{\mathbb{N}} C_*(U_i)$ und $H_*(M)= \colim_{i \in \mathbb{N}} H_*(U_i)$.
	Ebenso liegt jede kompakte Teilmenge von $M$ schon in einem $U_i$.
	Daher ist $H^*_c(M) = \colim_{i \in \mathbb{N}} H^*_c(U_i)$. 
	Nun ist $\PD_M = \colim_{i \in \mathbb{N}} \PD_{U_i}$ ein Isomorphismus.
\end{beweis}
% section 8 (end)
\newpage

\section{Mannigfaltigkeiten mit Rand} % (fold)
\label{sec:9}

\begin{definition}[{name=[Mannigfaltigkeit mit Rand]}]
	Eine \bet{$n$-Mannigfaltigkeit mit Rand}\index{Mannigfaltigkeit mit Rand} $M$ ist ein topologischer Hausdorffraum, der das zweite Abzählbarkeitsaxiom erfüllt und folgende lokale Eigenschaft hat: 
	Zu jedem Punkt $x \in M$ gibt es eine Umgebung $U \subseteq M$, die homöomorph ist zu einer offenen Teilmenge in $\mathbb{R}^{n-1} \times [0,\infty)$ ist.
	Die Punkte $x$, für die $U$ homöomorph zu einer offenen Teilmenge des $\mathbb{R}^n$ gewählt werden kann, heißen \Index{innere Punkte}; alle anderen \Index{Randpunkte}.
	Wir schreiben $\mathring{M} \subseteq M$ für die inneren Punkte und $\partial M := M \setminus \mathring{M}$.  \todo{Revision 9}
\end{definition}

\begin{bemerkung}[{name=[Der Rand ist eine $(n-1)$-Mannigfaltigkeit]}]
	$\mathring{M}$ ist eine $n$-Mannigfaltigkeit und $\partial M$ ist eine $(n-1)$-Mannigfaltigkeit.
\end{bemerkung}

\begin{beispiel}[{name=[Mannigfaltigkeit mit Rand]}]
	Sei $D^n = \set*{x \in \mathbb{R}^n \given \norm*{x}\le 1}$. Dann ist $\partial D^n = S^{n-1}$.
	Für $M \times D^n$ ist $\partial (M \times D^n) = M \times S^{n-1}$, wenn $M$ eine Mannigfaltigkeit ohne Rand ist.
\end{beispiel}

\begin{frage}
	Welche geschlossenen $n$-Mannigfaltigkeiten sind Rand einer kompakten $(n+1)$"=Mannigfaltigkeit?
\end{frage}

\begin{satz}[name={Kragensatz},label=kragensatz]
	Sei $M$ eine kompakte Mannigfaltigkeit mit Rand $\partial M$. 
	Dann gibt es eine offene Umgebung $U \subseteq M$ von $\partial M$ und einem Homöomorphismus $\varphi \colon \partial M \times [0,1) \to U$ mit ${\varphi(x,0)=x}$.
\end{satz}
\begin{beweis}[name={Skizze}]
	Sei $M' := M \cup_{\partial M} (\partial M \times [0,1])$.
	Es genügt einen Homöomorphismus $M \cong M'$ mit $x \mapsto (x,0)$ für $x \in \partial M$ zu finden. 
	Da $M$ kompakt ist, gibt es $\varphi_1, \ldots , \varphi_l \colon \partial M \to [0,1]$, sodass die $\varphi_i^{-1}\enbrace*{(0,1]}$ eine offene Überdeckung von $\partial M$ bilden und jedes $V_i$ in einer zu $\mathbb{R}^{n-1} \times [0,\infty)$ homöomorphen offenen Menge in $M$ liegt. 
	Nach Normalisierung können wir annehmen, dass $\varphi_1 + \ldots + \varphi_l =1$ ist.
	Für $0 \le k \le l$ sei $\psi_k := \varphi_1 + \ldots \varphi_k$ und 
	\[
		M_k := M \cup \set[\big]{(x,t) \given x \in \partial M, t \le \psi_k(x)}
	\]
	Nun zeigt man $M_{k-1} \cong M_k$ für $k=1,\ldots ,l$.
\end{beweis}

\begin{bemerkung}[{name=[Orientierung einer Mannigfaltigkeit mit Rand]}]
	Sei $M$ eine Mannigfaltigkeit mit Rand. 
	Eine $R$-Orientierung von $M$ ist eine $R$-Orientierung von $\mathring{M}$.
	Zu einer $R$-Orientierung $\set*{\mu_x}_{x \in \mathring{M}}$ gibt es, falls $M$ kompakt ist, eine eindeutige Orientierungsklasse $\mu \in H_n(M, \partial M)$ mit $\mu|_x = \mu_x$ für alle $x \in \mathring{M}$.
\end{bemerkung}

\begin{lemma}[{name=[Orientierbarkeit des Randes]}]
	Sei $M$ eine $R$-orientierbare kompakte Mannigfaltigkeit mit Rand.
	Dann ist auch $\partial M$ $R$-orientierbar.
\end{lemma}
\begin{beweis}
	Sei $N$ eine Zusammehangskomponente von $\partial M$.
	Mit $M$ ist nach dem \hyperref[kragensatz]{Kragensatz} auch $N \times [0,1]$ orientierbar.
	Nun betrachte einen Teil der Paarfolge zu $\enbrace*{N \times [0,1], \partial (N \times [0,1])}$
	\[
		\begin{tikzcd}[sep=1.8em]
			\underset{=H_n(N)=0}{H_n(N \times [0,1])} \rar & \underset{\neq 0}{H_n \enbrace[\big]{N \times [0,1], N \times \set*{0,1}}} \rar & \underset{= H_{n-1}(N) \oplus H_{n-1}(N) \neq 0}{H_{n-1} \enbrace[\big]{N \times \set*{0,1}}} \rar & H_{n-1} \enbrace*{N \times [0,1]}
		\end{tikzcd}
	\]
	Damit ist $H_n(N)\neq 0$ und $N$ ist $R$-orientierbar.
\end{beweis}
\begin{satz}[{name={Poincaré-Dualität für Mannigfaltigkeiten mit Rand}}]
	Sei $M$ eine kompakte $R$"=orientierbare $n$-Mannigfaltigkeit, deren Rand $\partial M$ die Vereinigung von zwei $(n-1)$-Mannigfaltigkeiten $\partial_0 M$ und $\partial_1 M$ entlang ihres gemeinsamen Randes 
	\[
		\partial \enbrace*{\partial_0 M } = \partial \enbrace*{\partial_1 M} = \partial_0 M \cap \partial_1 M
	\]
	ist.
	Sei $\mu \in H_n(M, \partial M)$ eine Orientierungsklasse für $M$. 
	Dann erhalten wir einen Isomorphismus 
	\[
		\capp \mu \colon H^k \enbrace*{M, \partial_0 M} \grenzw{\cong} H_{n-k} \enbrace*{M, \partial_1 M}
	\]
\end{satz}
\begin{beweis}
	Sei zunächst $\partial_1 M = \emptyset$, also $\partial _0 M = \partial M$.
	Betrachte 
	\[
		\begin{tikzcd}[column sep=3em, row sep=2.8em]
			\displaystyle H_c^k \enbrace*{M \setminus\partial M} = \smashoperator[r]{\colim_{K \subseteq M \setminus \partial M,\text{ kpt.} }}\,\, H^k(M, M \setminus K) \dar["\PD_{M \setminus \partial M}","\cong"']  \rar["\cong","(\star)"'] & H^k(M, \partial M) \dar["\capp \mu"] \\
			H_{n-k}\enbrace*{M\setminus \partial M}\rar["\cong","\text{Kragensatz}"'] & H_{n-k}(M)
		\end{tikzcd}
	\]
	zu $(\star)$: Wähle einen Kragen $\partial M \subseteq U \overset{\varphi}{\cong} \partial M \times [0,1)$ und setze $K_\varepsilon := M \setminus \varphi \enbrace*{\partial M \times [0,\varepsilon)}$.
	Dann sind die $K_\varepsilon$ in der gerichteten Menge aller kompakten Teilmengen von $M \setminus \partial M$ kofinal, also 
	\[
		\colim_{K \subseteq M \setminus \partial M,\text{ kpt.} } H^k(M, M \setminus K) = \colim_{\varepsilon\to 0} H^k(M,M \setminus K_\varepsilon) \cong H^k(M, \partial M) 
	\]
	Für $\partial_1 M \neq \emptyset$ benutzen wir die Paar- sowie die Tripelfolge:
	\[
		\begin{tikzcd}
			\ldots \rar & H^k(M, \partial M) \rar \ar[dd,"\capp \mu","\cong"'] & H^k(M,\partial_0 M) \rar \ar[dd,"\capp \mu","\cong"']& H^k(\partial M,\partial_0 M) \rar  \dar[phantom,sloped,"\cong"]& \ldots  \\
			& & & H^k \enbrace*{\partial_1 M, \partial \partial_1 M} \dar["\PD_{\partial_1 M}","\cong"'] \\
			\ldots \rar & H_{n-k}(M) \rar & H_{n-k}(M,\partial_1 M) \rar & H_{n-k-1}(\partial_1 M) \rar & \ldots 
		\end{tikzcd}
	\]
	wobei der Isomorphismus oben rechts durch Ausschneidung und Kragensatz zustande kommt.
\end{beweis}
\newpage
% section 9 (end)

\section{Die Signatur} % (fold)
\label{sec:10}

\begin{definition}[{name=[Schnittform]}]
	Sei $M$ eine $R$-orientierte geschlossene $n$-Mannigfaltigkeit mit Fundamentalklasse $\mu_M \in H_n(M;R)$.
	Die Abbildung
	\mapdef{\sprod*{\enspace}{\enspace}_R \colon H^k(M) \otimes H^{n-k}(M)}{R}{\xi \otimes \eta}{(\xi \cupp \eta)(\mu_M)}{}
	heißt die \Index{Schnittform} von $M$. \todo{Revision 10}
\end{definition}

\begin{bemerkung}[{name=[Begründung des Namens \enquote{Schnittform}]}]
	Eine kurze Begründung des Namens \enquote{Schnittform}:
	
	Seien für $i=0,1$ $N_i$ geschlossene $K$-orientierte $k_i$-Mannigfaltigkeiten mit Fundamentalklasse $\mu_i \in H_{k_i}(N_i;K)$.
	Seien $f_i \colon N_i \to M$ stetige Abbildungen in eine geschlossene $K$-orientierte $(n=k_0+k_1)$-Mannigfaltigkeit.
	Seien $\xi_i \in H^{n-{k_i}}$ die zu $(f_i)_*(\mu_i)$ Poincaré-dualen Kohomologieklassen. Dann \enquote{gilt}
	\[
		\sprod*{\xi_1}{\xi_2} = \# f_1(N_1) \cap f_2(N_2)
	\]
	nachdem $f_1(N_1)$ und $f_2(N_2)$ in $M$ in \enquote{allgemeiner Lage} sind und $\#$ zählt die Schnittpunkte mit Vorzeichen.
\end{bemerkung}

\begin{lemma}[{name=[Schnittform ist bilinear und graduiert kommutativ]}]
	Es gilt:
	\begin{enumerate}[(i),itemsep=2pt]
		\item Die Schnittform ist bilinear.
		\item $\sprod*{\xi}{\eta}_R = (-1)^{\abs*{\xi}\cdot \abs*{\eta}} \sprod*{\eta}{\xi}_R$
	\end{enumerate}
\end{lemma}
\begin{beweis}
	Folgt aus den entsprechenden Eigenschaften des $\cupp$-Produktes.
\end{beweis}

\begin{satz}[{name=[Für einen Körper sind Schnittformen regulär]}]
	Ist $R=K$ ein Körper, so ist die Schnittform \Index{regulär}, das heißt $\xi \mapsto \enbrace*{\eta \mapsto \sprod*{\xi}{\eta}_K}$ definiert einen Isomorphismus $H^k(M;K) \to \Hom_K \enbrace*{H^{n-k}(M;K),K}$.
\end{satz}
\begin{beweis}
	Wir haben einen Isomorphismen
	\begin{align}
		H^k(M;K) &\cong \Hom_K \enbrace*{H_k(M;K),K}  & \xi &\longmapsto  \enbrace[\big]{x \mapsto \xi(x)} \\
		\Hom_K \enbrace*{H_k(M;K),K} &\cong \Hom_K \enbrace*{H^{n-k}(M;K),K}   &  \varphi &\longmapsto  \enbrace[\big]{\eta \mapsto \varphi \enbrace*{\eta \capp \mu}}
	\end{align}
	Ihre Komposition bildet $\xi \in H^k(M;K)$ auf die Abbildung
	\[
		\eta \longmapsto \xi \enbrace*{\eta \capp \mu} \StackText{\ref{bem_augmentation_cap}}{=} \epsilon \enbrace[\big]{\xi \capp (\eta \capp \mu)} = \epsilon \enbrace[\big]{(\xi \cupp \eta) \capp \mu} = (\xi \cupp \eta)(\mu) = \sprod*{\xi}{\eta}_K
	\]
	ab und es folgt die Behauptung.
\end{beweis}

\begin{korollar}[{name=[Kohomologiemoduln über $K$ sind endlichdimensional]}]
	Für $K$-orientierbare geschlossene Mannigfaltigkeiten sind die Kohomologiemoduln $H^*(M;K)$ (und auch die Homologiemoduln $H_*(M;K)$) endlich dimensional.
\end{korollar}
\begin{beweis}
	Wir benutzen, dass Vektorräume genau dann endlich dimensional sind, wenn sie isomorph sind zu ihrem Bidualraum sind.
	Wegen 
	\[
		H^k(M;K) \cong \Hom_K \enbrace*{H^{n-k}(M;K),K} \cong \Hom_K \enbrace*{\Hom_K \enbrace*{H^k(M;K),K}} 
	\]
	folgt $\dim_K H^k(M;K) < \infty$.
\end{beweis}

\begin{lemma}
	Sei $\varphi \colon V \times V \to \mathbb{R}$ eine symmetrische reguläre Bilinearform auf einem endlich dimensionalen $\mathbb{R}$-Vektorraum $V$.
	Dann gibt es Unterräume $V_+$ und $V_-$ mit $\varphi|_{V_\mp \times V_\mp}$ negativ/positiv definit und $\varphi|_{V_- \times V_+}=0$ und $V= V_- \oplus V_+$.
	Die Zahlen $n_- := \dim_\mathbb{R} V_-$ und $n_+ := \dim_\mathbb{R} V_+$ hängen nur von $\varphi$ ab.
\end{lemma}
\begin{beweis}
	Sei $B$ eine Basis von $V$ und $f \colon V \to V$ bestimmt durch $f(b)= \sum_{b' \in B} \varphi(b,b') \cdot b'$.
	Dann ist die zugehörige Matrix $A= \enbrace*{\varphi(b,b')}_{b,b' \in B}$ symmetrisch und $V$ hat eine Basis aus Eigenvektoren von $f$.
	Sei $V_- \le V$ von Eigenvektoren zu negativen Eigenwerten erzeugt und $V_+ \le V$ von Eigenvektoren zu positiven Eigenwerten erzeugt.
	Da $\varphi$ regulär ist, ist $0$ kein Eigenwert von $f$ und daher ist $V= V_+ \oplus V_-$.
	Seien $v,v'$ Eigenvektoren mit Eigenwerten $\lambda,\lambda' \in \mathbb{R}$.
	Dann ist
	\begin{align}
		(\lambda')^{-1} \sum \varphi(v,b) \varphi(v',b)&= \varphi \enbrace*{v, (\lambda')^{-1} \sum \varphi(v',b) b} \\ &= \varphi(v,v')\\
		&= \varphi \enbrace*{\lambda^{-1} \sum \varphi(v,b)b , v'} = \lambda^{-1} \sum \varphi(v,b) \varphi(b,v')
	\end{align}
	Also ist $\lambda=\lambda'$ oder $\varphi(v,v')=0$. 
	Für die Eindeutigkeit der Dimensionen sei $V=W_- \oplus W_+$ eine zweite solche Zerlegung, so gilt $V_ \cap W_+ =0$ also $\dim V_- + \dim W_+ \le n$, bzw. $V_+ \cap W_- =0$ also $\dim V_+ + \dim W_- \le n$.
	Mit $\dim V_- + \dim V_+ = \dim W_- + \dim W_+ = n$ folgt $\dim V_{\mp} = \dim W_{\mp}$.
\end{beweis}

\begin{definition}[{name=[Signatur]}]
	Die \bet{Signatur}\index{Signatur!symmetrische reguläre Bilinearform} von $\varphi$ ist definiert als $\sig(\varphi) := n_+ - n_-$. 
\end{definition}

\begin{korollar}
	Es gibt eine Basis $v_1, \ldots , v_{n_-}, w_1 , \ldots , w_{n_+}$ von $V$ mit 
	\[
		\varphi(v_i,v_j) = - \delta_{ij} \quad , \quad \varphi(w_i,w_j) =\delta_{ij} \quad , \quad \varphi(v_i,w_j)=0
	\]
\end{korollar}
\begin{beweis}
	Wähle eine Orthonormalbasis von $V_-$ bezüglich $-\varphi$ und eine Orthonormalbasis bezüglich $\varphi$.
\end{beweis}

\begin{lemma}[label=sig_null]
	Die Signatur von $\varphi$ ist genau dann $0$, wenn es $L \le V$ gibt mit $\varphi|_{L \times L}=0$ und $\dim L = \frac{1}{2} \dim V$.
\end{lemma}
\begin{beweis}
	Sei $\sig(\varphi)=0$. Dann wählen wir eine Basis $v_1,\ldots ,v_l,w_1,\ldots ,w_l$ wie im Korollar\todo{oder doch Lemma?}.
	Dann hat $L := \langle v_1+ w_1, \ldots ,v_l+w_l$ die gesuchte Eigenschaft.
	
	Sei umgekehrt $L \le V$ mit $\varphi|_{L \times L}=0$ und $\dim L = \frac{1}{2} \dim V$.
	Dann ist $L \cap V_- =0$, also $\dim L + \dim V_- \le \dim V$ und genauso $L \cap V_+ =0$, also $\dim L + \dim V_+ \le \dim V$.
	Damit folgt $\dim L = \dim V_- = \dim V_+ = \frac{1}{2} \dim V$.
\end{beweis}

\begin{definition}[{name=[Signatur]}]
	Sei $M$ eine $4k$-dimensionale geschlossene orientierte Mannigfaltigkeit.
	Dann ist $\sprod*{\enspace}{\enspace}_\mathbb{R} \colon H^{2k}(M;\mathbb{R}) \otimes H^{2k}(M;\mathbb{R}) \to \mathbb{R}$ eine reguläre symmetrische Bilinearform.
	Wir definieren die \bet{Signatur}\index{Signatur!Mannigfaltigkeit} $\sig(M)$ von $M$ als die Signatur dieser Bilinearform.
\end{definition}

\begin{bemerkung}
	\leavevmode
	\begin{itemize}
		\item Es ist $H_n(M;\mathbb{R}) \cong H_n(M;\mathbb{Z})\otimes \mathbb{R}$.
	Insbesondere bestimmt jede $\mathbb{Z}$-Orientierung von $M$ eine $\mathbb{R}$-Orientierung von $M$.
	\item Für eine orientierte Mannigfaltigkeit $M$ bezeichnen wir mit $-M$ die Mannigfaltigkeit $M$ mit der umgekehrten Orientierung (also $\mu_{-M}=-\mu_M$).
	Ist $M$ geschlossen mit $\dim M = 4k$, so gilt $\sig(M)=- \sig(-M)$.
	\end{itemize}
\end{bemerkung}

\begin{beispiel}
	Sei $\alpha \in H^2(\mathbb{CP}^{2k})$ ein Erzeuger.
	Wir wählen die Orientierung $\mu \in H_{4k} \enbrace{\mathbb{CP}^{2k}}$ für die $\alpha^{2k} (\mu)=1$.
	Mit dieser Orientierung ist $\sig \enbrace{\mathbb{CP}^{2k}}=1$.
\end{beispiel}

\begin{satz}
	Sei $W$ eine kompakte orientierte Mannigfaltigkeit mit Rand.
	Falls $\dim W = 4k+1$ ist, so gilt $\sig(\partial W)=0$.
\end{satz}
\begin{beweis}
	Sei $L :=  \im \enbrace[\big]{i^* \colon H^{2k}(W;\mathbb{R}) \hookrightarrow H^{2k}(\partial W;\mathbb{R})}$. Es gilt
	\begin{enumerate}[a)]
		\item $(i^* \xi \cupp i^* \eta) (\mu_{\partial W}) = \enbrace*{i^*(\xi \cupp \eta)}(\mu_{\partial W})= (\xi \cupp \eta) \enbrace*{i_* \mu_{\partial W}} =0$, da $\mu_{\partial W}=\partial (\mu_{W,\partial W})$, wobei $\mu_{W,\partial W} \in H^{4k+1}(W,\partial W)$ die Orientierungsklasse von $W$ relativ zu $\partial W$ ist.
		\item Zu zeigen: $\dim L = \frac{1}{2} \dim H^{2k}(\partial W;\mathbb{R})$.
		
		Betrachte 
		\[
			\begin{tikzcd}[sep=large]
				\ldots  \rar & H^{2k}(W) \dar["\capp \mu_{W,\partial W}","\cong"'] \rar["i^*"] & H^{2k}(\partial W)  \dar["\capp \mu_{\partial W}","\cong"']\rar["\partial"] & H^{2k+1}(W,\partial W) \rar \dar["\capp \mu_{W,\partial W}","\cong"']& \ldots \\
				\ldots \rar & H_{2k+1} \rar["\partial"] & H_{2k}(\partial W) \rar & H_{2k}(W) \rar & \ldots 
			\end{tikzcd}
		\]
		Die Kommutativität benutzt wieder $\partial (\mu_{W,\partial W})=\mu_{\partial W}$.
		Es gilt
		\[
			\dim \im i^* = \dim \ker i_* = \dim \ker \enbrace[\big]{\Hom(i_*,\mathbb{R})} = \dim \ker i^*
		\]
		Da $\dim H^{2k}(\partial W) = \dim \ker i^* + \dim \im i^*$, folgt $\dim \im i^* = \frac{1}{2} \dim H^{2k}(\partial W)$. 
	\end{enumerate}
	Mit dem \autoref{sig_null} folgt die Behauptung.
\end{beweis}

\begin{korollarB}
	$\mathbb{CP}^{2k}$ ist nicht Rand einer orientierten kompakten Mannigfaltigkeit.
\end{korollarB}

\begin{ausblick}
	Nach dem Satz von \textsc{Freedman}\footnote{\url{https://de.wikipedia.org/wiki/Michael_Freedman}} werden einfach zusammenhängende geschlossene orientierte $4$-Mannigfaltigkeiten bis auf Homöomorphismus durch die ganzzahlige Schnittform auf $H^2(M;\mathbb{Z})$ und eine $\sfrac{\mathbb{Z}}{2 \mathbb{Z}}$-wertige Invariante klassifiziert.
	Jede ganzzahlige reguläre Schnittform wird realisiert.
\end{ausblick}
% section 10 (end)
\newpage

\section{Projektive Auflösungen} % (fold)
\label{sec:11}

\begin{definition}[{name=[projektive Auflösung]}]
	Sei $M$ ein $R$-Modul.
	Eine \Index{projektive Auflösung} von $M$ ist eine lange exakte Folge 
	\[
		\begin{tikzcd}
			0 & M \lar & P_0 \lar["\epsilon"'] & P_1 \lar["d_0"'] & P_2 \lar["d_1"'] & \ldots \lar
		\end{tikzcd}
	\]
	wobei die $P_i$ projektive $R$-Moduln sind.
	Wir schreiben oft $(P_*,\epsilon)$ für eine solche Auflösung, wobei $P_*$ ein Kettenkomplex projektiver Moduln ist, $H_k(P_*)=0$ für $k>0$ und $H_0(P_*)\cong M$ via $\epsilon$ ist.\todo{Revision 11}
\end{definition}

\begin{satz}[{name=[Existenz projektiver Auflösungen]}]
	Sei $M$ ein $R$-Modul. Dann gilt
	\begin{enumerate}[a),itemsep=0pt]
		\item $M$ besitzt eine projektive Auflösung.
		\item Sind $(P_*,\epsilon)$ und $(P_*',\epsilon')$ zwei projektive Auflösungen, so gibt es eine Kettenhomotopieäquivalenz $f \colon P_* \to P_*'$ mit $H_0(f)= (\epsilon') ^{-1} \circ \epsilon$.
	\end{enumerate}
\end{satz}
\begin{beweis}
	\begin{enumerate}[a),itemsep=0pt]
		\item Die $P_n$ werden induktiv definiert. 
		Sei 
		$\begin{tikzcd}[cramped,sep=1.5em]
			0 & M \lar & P_0 \lar["\epsilon"'] & P_1 \lar["d_0"'] & P_2 \lar["d_1"'] & \ldots \lar & P_n \lar["d_n"']
		\end{tikzcd}$
		bereits definiert und $S$ ein Erzeugendensystem von $\ker d_n$.
		Wir setzen $P_{n+1} := R[S]$ und definieren $d_{n+1} \colon P_{n+1} \to P_n$ als die Fortsetzung der Inklusion $S \subseteq P_n$ zu einer $R$-linearen Abbildung $R[S]\to P_n$.
		\item Folgt aus der folgenden \autoref{prop:113}.\qedhere
	\end{enumerate}
\end{beweis}

\begin{proposition}[label=prop:113]
	Seien $A_*$ und $P_*$ zwei Kettenkomplexe mit 
	\begin{itemize}[itemsep=0pt]
		\item Die Kettenmoduln von $P_*$ sind projektiv.
		\item $H_*(A_*)=0$ für $*\ge 1$.
	\end{itemize}
	Dann definiert $f \mapsto H_0(f)$ einen Isomorphismus zwischen Kettenhomotopieklassen von Kettenabbildungen $f \colon P_* \to A_*$ und $R$-linearen Abbildungen $\varphi \colon H_0(P_*) \to H_0(A_*)$.
\end{proposition}
\begin{beweis}
	Sei $\varphi \colon H_0(P_*) \to H_0(A_*)$ gegeben. 
	Wir konstruieren $f_n \colon P_n \to A_n$ induktiv: 
	Seien entsprechende Abbildungen $f_0,\ldots ,f_n$ bereits konstruiert, sodass folgendes Diagramm kommutiert
	\[
		\begin{tikzcd}[sep=large]
			H_0(P_*) \dar["\varphi"] & P_0 \lar[two heads]  \dar["f_0"]& P_1 \lar  \dar["f_1"]& \ldots \lar & P_n \lar["d_n"']\dar["f_n"] & P_{n+1} \lar["d_{n+1}"'] \dar[dashed,"f_{n+1}"] \\
			H_0(A_*) & A_0 \lar[two heads] & A_1 \lar["d_1'"'] & \ldots \lar & A_n \lar["d_n'"'] & A_{n+1} \lar["d_{n+1}'"']
		\end{tikzcd}
	\]
	Wir erhalten 
	\[
		\begin{tikzcd}[row sep=large, column sep=huge]
			\mbox{ } & A_{n+1} \dar[two heads,"d_{n+1}'"]\\
			P_{n+1} \ar[ur,dashed] \rar["f_n \circ d_{n+1}"] & \ker d_{n}'
		\end{tikzcd}
	\]
	Da $P_{n+1}$ projektiv ist, gibt es $f_{n+1} \colon P_{n+1} \to A_{n+1}$ mit $d_{n+1}' \circ f_{n+1} = f_n \circ d_{n+1}$.
	
	Für die Injektivität seien $f,g \colon P_* \to A_*$ Kettenabbildung mit $H_0(f) = H_0(g)$. 
	Wir konstruieren eine Kettenhomotopie $s_n \colon P_n \to A_{n+1}$ induktiv. Seien $s_0, \ldots ,s_{n-1}$ schon konstruiert mit $f_i -g_i= s_{i-1} d_i + d_{i+1} s_i$ für $i=0, \ldots , n$.
	\[
		\begin{tikzcd}[row sep=large, column sep=3.15em]
			H_0(P_*) \dar["0"] & P_0 \lar \dar["f_0-g_0"'] \drar[sloped,pos=.4,"s_0"] & P_1 \lar \dar[pos=.7,"f_1-g_1"] \drar[sloped,pos=.4,"s_1"] & P_2 \lar \dar["f_2 - g_2"] & \ldots \lar & P_{n-1} \lar\dar["f_{n-1}- g_{n-1}"'] \drar[sloped,pos=.3,"s_{n-1}"] & P_n \lar \dar[pos=.7,"f_n - g_n"]  \drar[dashed,sloped,pos=.4,"s_n"] & P_{n+1} \lar \\
			H_0(A_*) & A_0 \lar & A_1 \lar & A_2 \lar & \ldots \lar & A_{n-1} \lar & A_n \lar & A_{n+1} \lar 
		\end{tikzcd}
	\]
	Wir erhalten
	\[
		\begin{tikzcd}[row sep=large, column sep=7.5em]
			\mbox{} &  A_{n+1} \dar["d_{n+1}"] \\
			P_n \rar["(f_n-g_n)- s_{n-1} d_n"'] \urar[dashed] &\ker d_n
		\end{tikzcd}
	\]
	Da $P_n$ projektiv ist, gibt es $s_n \colon P_n \to A_{n+1}$ mit $d'_{n+1} \circ s_n = (f_n -g_n) - s_{n-1} \circ d_n$.
\end{beweis}

\begin{definition}[{name=[projektive Dimension]}]
	Sei $R$ ein Ring.
	Die \Index{projektive Dimension} von $R$ ist das kleinste $n$, sodass jeder $R$-Modul eine projektive Auflösung der Länge $n$ hat, also eine Auflösung der Form
	\[
		\begin{tikzcd}
			M & P_0 \lar & P_1 \lar & P_2 \lar & \ldots \lar & P_n \lar & 0 \lar
		\end{tikzcd}
	\] 
	Gibt es kein solches $n \in \mathbb{N}$, so sagen wir die projektive Dimension von $R$ ist $\infty$.
\end{definition}

\begin{beispiel}[{name=[Körper haben projektive Dimension $0$]}]
	Ist $R=K$ ein Körper, so ist jeder $K$-Modul, also jeder $K$-Vektorraum, frei.
	Daher ist die projektive Dimension $0$.
\end{beispiel}

\begin{satz}[{name=[projektive Dimension von Hauptidealringen]}]
	Für Hauptidealringe ist die projektive Dimension immer $\le 1$.
\end{satz}
\begin{beweis}
	Ist $M$ ein $R$-Modul, so gibt es einen freien $R$-Modul $P_0$ mit einer surjektiven Abbildung $\epsilon \colon P_0 \to M$.
	Ist $R$ ein Hauptidealring, so ist $P_1 := \ker \epsilon$ als Untermodul eines freien Moduls wieder frei\todo{Referenz} und wir erhalten wir eine projektive Auflösung
	\(
		\begin{tikzcd}[cramped,sep=small]
			M & P_0 \lar & P_1 \lar
		\end{tikzcd}
	\)
	der Länge $1$.
\end{beweis}

\begin{lemma}
	Sei 
	\[
		\begin{tikzcd}
			0 \rar & M \rar["f"] & M' \rar["f'"] & M'' \rar & 0
		\end{tikzcd}
	\]
	eine kurze exakte Sequenz von $R$-Moduln. Dan gilt
	\begin{enumerate}[a)]
		\item Sind $\epsilon \colon P \to M$, $\epsilon'' \colon P'' \to M''$ surjektiv, wobei $P''$ projektiv, so gibt es $\epsilon' \colon P \oplus P'' \to M'$ surjektiv, sodass folgendes Diagramm kommutiert:
		\[
			\begin{tikzcd}
				0 \rar & P \rar[hook] \dar["\epsilon"] & P \oplus P'' \rar[two heads] \dar["\epsilon'"] & P'' \rar  \dar["\epsilon''"] & 0 \\
				0 \rar & M \rar["f"] & M' \rar["f'"] & M'' \rar & 0
			\end{tikzcd}
		\]
		\item Es gibt projektive Auflösungen $(P_*,\epsilon)$, $(P_*',\epsilon')$ und $(P_*'',\epsilon'')$ von $M$, $M'$ und $M''$ und eine kurze exakte Folge von Kettenkomplexen
		\[
			\begin{tikzcd}
				0 \rar & P_* \rar["F"] & P_*' \rar["F'"] & P_*'' \rar & 0
			\end{tikzcd}
		\]
		sodass folgendes Diagramm kommutiert:
		\[
			\begin{tikzcd}[sep=large]
				H_0(P_*) \rar["H_0(F)"] \dar["\epsilon","\cong"'] & H_0(P_*') \dar["\epsilon'","\cong"'] \rar["H_0(F')"] & H_0(P_*'') \dar["\epsilon''","\cong"'] \\
				M \rar & M' \rar & M''
			\end{tikzcd}
		\]
	\end{enumerate}
\end{lemma}
\begin{beweis}
	\begin{enumerate}[a)]
		\item Da $f' \colon M' \to M''$ surjektiv und $P''$ projektiv ist, gibt es $\hat{\epsilon}'' \colon P'' \to M'$ mit $\epsilon'' = f' \circ \hat{\epsilon}''$.
		Nun können wir $\epsilon' \colon P \oplus P'' \to M'$ durch $\epsilon' (p,p'') := f \enbrace*{\epsilon(p)} + \hat{\epsilon}''(p'')$ definieren.
		Noch zu zeigen: $\epsilon'$ ist surjektiv. 
		Sei $v' \in M'$. Da $\epsilon''$ surjektiv ist, gibt es $x'' \in P''$ mit $\epsilon''(x'')=f'(v')$. Dann gilt
		\[
			f'\enbrace*{v' - \hat{\epsilon}''(x'')} = f'(v') - \Underbracket{f' \enbrace*{\hat{\epsilon}''(x'')}}{=\epsilon''(x'')} = 0
		\]
		Daher gibt es $v \in M$ mit $f(v)=v' - \hat{\epsilon}''(x'')$.
		Sei nun $x \in P$ nut $\epsilon(x)=v$.
		Dann ist $\epsilon'(x,x'')=v'$.
		\item Wähle projektive Auflösungen $(P_*,\epsilon)$ von $M$ und $(P_*'',\epsilon'')$ von $M''$.
		Sei $P_n' := P_n \oplus P_n''$.
		Durch induktive Anwendung von a) erhalten wir Abbildungen $\epsilon' \colon P_0' \to M'$, $d_i' \colon P_i' \to P_{i-1}'$, sodass $\begin{tikzcd}[cramped,sep=small]
			P_* \rar & P_*' \rar & P_*''
		\end{tikzcd}$ eine kurze exakte Sequenz von Kettenkomplexen ist und das Diagramm aus der Behauptung kommutiert.\qedhere
	\end{enumerate}
\end{beweis}
\newpage
% section 11 (end)

\section[Tor und Ext]{$\Tor$ und $\Ext$} % (fold)
\label{sec:12}

\todo[inline]{Revision 12}
\begin{frage}
	Sei $V$ ein $\mathbb{Z}$-Modul und $V''$ ein $R$-Modul. Ist $H_*(X;V) \cong H_*(X) \otimes V$? 
	\[
		H_* \enbrace*{V' \otimes_R C_*(X;R)}= H_*(X;V') \cong V' \otimes_R H_*(X;R)
	\]
	Sei $W$ ein $R$-Modul. Gilt dann
	\[
		H^*(X;V) \cong \Hom_\mathbb{Z}(H_*(X),V) \qquad H^*(X;W) \cong \Hom_R \enbrace*{H_*(X;R),W}
	\]
	Genauso für Kettenkomplexe: Sei $C_*$ ein $\mathbb{Z}$-Kettenkomplex ud $C_*'$ ein $R$-Kettenkomplex. Gilt dann
	\[
		H_* \enbrace*{V \otimes_\mathbb{Z} C_*} \cong V \otimes _\mathbb{Z} H_*(C_*) \qquad H_* \enbrace*{V' \otimes_R C_*'} \cong V' \otimes_R H_*(C_*)
	\]
	und wie sieht es hiermit aus:
	\[
		H^* \enbrace[\big]{\Hom_\mathbb{Z}(C_*,V)} \cong \Hom_\mathbb{Z} \enbrace[\big]{H_*(C_*),V} \qquad H^*\enbrace[\big]{\Hom_R(C_*',W)} \cong \Hom_R \enbrace[\big]{H_*(C_*'),W}
	\]
\end{frage}

\begin{konvention}[{name=[Links- und Rechtsmoduln]}]
	Moduln sind -- wie gewohnt -- immer Linksmoduln, außer wir sprechen explizit von Rechtsmoduln
\end{konvention}

\begin{definition}[{name=[{lineare und exakte Funktoren}]}]
	Seien $R$ und $S$ Ringe.
	Ein Funktor $F\colon \Mod{R} \to \Mod{S}$ heißt \bet{$\mathbb{Z}$-linear}\index{Funktor!$\mathbb{Z}$-linearer}, wenn für $f,g \colon M \to M'$ gilt $F(f + g)= F(f) + F(g)$.
	
	Ein solcher Funktor heißt \bet{exakt}\index{Funktor!exakter}, wenn er kurze exakte Folgen von $R$-Moduln auf kurze exakte Folgen von $S$-Moduln abbildet.
\end{definition}

\begin{bemerkung}[{name=[exakte Funktoren und (Ko-)Homologie]}]
	Sei $C_*$ ein $R$-Kettenkomplex.
	Ist $F \colon \Mod{R} \to \Mod{S}$ exakt, so gilt
	\[
		\begin{array}{ll}
			H_* \enbrace[\big]{F(C_*)} = F_* \enbrace[\big]{H_*(C_*)} & \text{ falls $F$ kovariant ist} \\
			H^* \enbrace[\big]{F(C_*)} = F \enbrace[\big]{H_*(C_*)} & \text{ falls $F$ kontravariant ist} 
		\end{array}
	\]
\end{bemerkung}

\begin{beispiel}
	\begin{enumerate}[(i)]
		\item Sei $V$ ein $R$-Rechtsmodul.
		Dann erhalten wir einen $\mathbb{Z}$-linearen Funktor 
		\[
			V \otimes_R \colon \Mod{R} \to \Mod{\mathbb{Z}}
		\]
		\item Sei $W$ ein $R$-Modul. 
		Dann erhalten wir einen $\mathbb{Z}$-linearen Funktor 
		\[
			\Hom_R(-,W) \colon \Mod{R} \to \Mod{\mathbb{Z}}
		\]
		\item Für $R=\mathbb{Z}$ und $V= \sfrac{\mathbb{Z}}{2 \mathbb{Z}}$ sind weder $\sfrac{\mathbb{Z}}{2 \mathbb{Z}} \otimes $ noch $\Hom_\mathbb{Z} \enbrace*{-, \sfrac{\mathbb{Z}}{2 \mathbb{Z}}}$ exakt:
		\[
			\begin{tikzcd}
				0 \rar & \mathbb{Z} \rar["\cdot 2"] & \mathbb{Z} \rar & \sfrac{\mathbb{Z}}{2 \mathbb{Z}} \rar & 0
			\end{tikzcd}
		\]
		wird nach Tensorieren zu 
		\[
			\begin{tikzcd}
				\sfrac{\mathbb{Z}}{2 \mathbb{Z}} \rar["0"] & \sfrac{\mathbb{Z}}{2 \mathbb{Z}} \rar["\cong"] & \sfrac{\mathbb{Z}}{2 \mathbb{Z}} \rar & 0
			\end{tikzcd}
		\]
		und nach Anwenden von $\Hom(-,\sfrac{\mathbb{Z}}{2 \mathbb{Z}})$ zu
		\[
			\begin{tikzcd}
				\sfrac{\mathbb{Z}}{2 \mathbb{Z}}  & \sfrac{\mathbb{Z}}{2 \mathbb{Z}} \lar["0"']  & \sfrac{\mathbb{Z}}{2 \mathbb{Z}}\lar["\cong"']  & 0 \lar
			\end{tikzcd}
		\]
	\end{enumerate}
\end{beispiel}

\begin{lemma}
	Sei $V$ ein $R$-Rechtsmodul und $W$ ein $R$-Modul.
	Sei $\begin{tikzcd}[cramped,sep=small]
		0 \rar& M \rar &M' \rar & M'' \rar & 0
	\end{tikzcd}$ 
	eine kurze exakte Folge von $R$-Moduln.
	Dann sind folgende Sequenzen exakt:
	\begin{enumerate}[(i)]
		\item \( \begin{tikzcd}[cramped,sep=2.3em]
					V \otimes_R M \rar & V \otimes_R M' \rar & V \otimes_R M'' \rar & 0
				\end{tikzcd} \)
		\item $\begin{tikzcd}[cramped,sep=1.8em]
			\Hom_R(M,W)  & \Hom_R(M',W) \lar & \Hom_R(M'',W) \lar & 0 \lar
		\end{tikzcd}$
	\end{enumerate}
\end{lemma}
\begin{beweis}
	\emph{Übungsaufgabe 2 von Blatt 5.}
\end{beweis}

\begin{lemma}
	Sei $V$ ein $R$-Rechtsmodul, $W$ und $M$  $R$-Moduln.
	Sei $\epsilon \colon P_* \to M$ eine projektive Auflösung.
	Dann sind die (Ko)"=Kettenkomplexe $V \otimes_R P_*$ und $\Hom_R(P_*,W)$ bis auf (Ko)"=Kettenhomotopieäquivalenz unabhängig von der Wahl der Auflösung $P_*$ und es gilt
	\[
		H_0(V \otimes_R P_*) = V \otimes_R M \qquad \text{ und } \qquad H^0 \enbrace[\big]{\Hom_R(P_*,W)} = \Hom_R(M,W)
	\]
\end{lemma}

\noindent\todo[inline]{Hier fehlt noch was \ldots insbesondere kann sich die Nummerierung noch ändern}

\begin{definition}[{name=[$\Tor$ und $\Ext$]}]
	Wir setzen
	\[
		\Tor_n^R(V,M) := H_n(V \otimes_R P_*) \qquad \text{ und } \qquad \Ext_R^n(M,W) := H^n \enbrace*{\Hom_R(P_*,W)} 
	\]
\end{definition}

\begin{bemerkung}
	Es gilt $\Tor_0^R(V,M) = V \otimes_R M$ und $\Ext_R^0(M,W)=\Hom_R(M,W)$.
	Unter einigen zusätzlichen Annahmen, haben wir folgende einfache Schlussfolgerungen für $\Tor$ und $\Ext$:
	\begin{enumerate}[(i)]
		\item Ist $M$ projektiv, so ist $\Tor_n^R(V,M)=0$ und $\Ext_R^n(M,W)=0$ für $n>0$.
		\item Ist $K$ ein Körper, so ist $\Tor_n^K(V,M)=0$ und $\Ext_K^n(M,W)=0$ für $n>0$.
		\item Ist $R$ ein Hauptidealring, so ist $\Tor_n^R(V,M)=0$ und $\Ext_R^n(M,W)=0$ für $n>1$.
	\end{enumerate}
	$\Tor^R_n$ und $\Ext_R^n$ sind Funktoren in $V,M$ für $\Tor$ (kovariant) und $M$ (kontravariant), $W$ (kovariant) für $\Ext$.
\end{bemerkung}

\begin{satz}
	Sei 
	\(
		\begin{tikzcd}[cramped]
			0 \rar & M \rar & M' \rar & M'' \rar & 0
		\end{tikzcd}
	\)
	eine kurze exakte Folge von $R$-Moduln.
	Dann gibt es lange exakte Folgen \todo{Verbindungsho\-momorphismen hinzufügen}
	\begin{enumerate}[(i)]
		\item \[
			\begin{tikzcd}
				\ldots \rar & \Tor^R_n(V,M) \rar & \Tor^R_n(V,M') \rar & \Tor_n^R(V,M'') & \\
				&\Tor^R_{n-1}(V,M) \rar & \ldots \rar & \Tor_1^R(V,M'') \\
				&V \otimes_R M \rar & V \otimes_R M' \rar & V \otimes_R M'' \rar & 0 
			\end{tikzcd}
		\]
		\item \[
			\begin{tikzcd}
				\ldots & \Ext_R^n(M,W) \lar & \Ext^n_R(M',W) \lar & \Ext_R^n(M'',W) \lar \\
				&\Ext_R^{n-1}(M,W) & \lar \ldots \\
				&\Hom_R(M,W) & \Hom_R(M',W) \lar & \Hom_R(M'',W) \lar & 0 \lar
			\end{tikzcd}
		\]
	\end{enumerate}
\end{satz}
\begin{beweis}
	Wir können Auflösungen $(P_*,\epsilon)$, $(P_*',\epsilon')$ und $(P_*'',\epsilon'')$ von $M$, $M'$ und $M''$ wählen, sodass wir eine kurze exakte Folge
	\(
		\begin{tikzcd}[cramped,sep=small]
			0 \rar & P_* \rar & P_*' \rar & P_*'' \rar & 0
		\end{tikzcd}
	\)
	erhalten, für die folgendes Diagramm kommutiert: 
	\[
		\begin{tikzcd}
			P_* \rar \dar["\epsilon"] & P_*' \dar["\epsilon'"] \rar & P_*'' \dar["\epsilon''"] \\
			M \rar & M' \rar & M''
		\end{tikzcd}
	\]
	Da $P_n''$ projektiv ist, spaltet 
	\(
		\begin{tikzcd}[cramped,sep=small]
			0 \rar & P_n \rar & P_n' \rar & P_n'' \rar & 0
		\end{tikzcd}
	\)
	für jedes $n$ und
	\[
		\begin{tikzcd}
			0 \rar & V \otimes_R P_* \rar & V \otimes_R P_*' \rar & V \otimes_R P_*'' \rar &0
		\end{tikzcd}
	\]
	sowie
	\[
		\begin{tikzcd}
			0 & \Hom_R(P_*,W) \lar & \Hom_R(P_*',W) \lar & \Hom_R(P_*'',W) \lar & 0 \lar
		\end{tikzcd}
	\]
	sind immer noch kurz exakt. 
	Die lange exakte Folge in (i) und (ii) sind die zugehörigen langen exakten Folgen in Homologie bzw. Kohomologie.
\end{beweis}

\begin{bemerkung}
	Es gibt natürliche Isomorphismen für $\Tor$:
	\[
		\bigoplus_{i \in I} \Tor_n^R(V_i,M) \cong \Tor_n^R \enbrace*{\bigoplus_{i \in I} V_i, M} \qquad \text{ und } \qquad 
		\bigoplus_{i \in I} \Tor_n^R(V,M_i) \cong \Tor_n^R \enbrace*{V, \bigoplus_{i \in I} M} 
	\]
	und für $\Ext$:
	\[
		\Ext_R^n \enbrace*{\bigoplus_{i \in I} M,W} \cong \prod_{i \in I} \Ext_R^n(M_i,W) \qquad \text{ und } \qquad \Ext_R^n \enbrace*{M, \prod_{i \in I} W_i} \cong \prod_{i \in I} \Ext_R^n(M,W_i) 
	\]
\end{bemerkung}

\begin{bemerkung}
	Sei $R$ ein Hauptidealring. 
	Seien $a,b \in R\setminus \set*{0}$ und $d$ ein größter gemeinsamer Teiler von $a$ und $b$.
	Sei $M$ ein $R$-Modul.
	Dann gilt 
	\begin{multicols}{2}
		\begin{enumerate}[(i)]
			\item $\sfrac{R}{(a)} \otimes_R \sfrac{R}{(b)} \cong \sfrac{R}{(d)}$
			\item $\Hom_R \enbrace*{\sfrac{R}{(a)},\sfrac{R}{(b)}} \cong \sfrac{R}{(d)}$
			\item $\Tor_1^R \enbrace*{\sfrac{R}{(a)}, \sfrac{R}{(b)}} \cong \sfrac{R}{(d)}$
			\item $\Ext^1_R \enbrace*{\sfrac{R}{(a)}, \sfrac{R}{(b)}} \cong \sfrac{R}{(d)}$
			\item $M \otimes_R R \cong R \otimes_R M \cong M$
			\item $\Hom_R(R,M) \cong M$
			\item $\Tor^R_1(R,M) \cong \Tor_1^R(M,R)=0$
			\item $\Ext^1_R(R,M) =0$
			\item $\Hom_R \enbrace*{\sfrac{R}{(a)},R}=0$
			\item $\Ext_R^1 \enbrace*{\sfrac{R}{(a)},R}=0$
		\end{enumerate}
	\end{multicols}
\end{bemerkung}

\begin{bemerkung}
	Sei $(P_*,\epsilon)$ ein projektive Auflösung von $V$ durch projektive $R$-Rechtsmoduln.
	Dann gilt 
	\[
		\Tor_n^R(V,M) \cong H_n(P_* \otimes M)
	\]
	Beweis kommt später noch.\todo{Referenz wenn soweit}
\end{bemerkung}

\begin{korollarB}
	Für $R$ kommutativ gilt: $\Tor_n^R(V,M) \cong \Tor_n^R(M,V)$.
\end{korollarB}

\begin{definition}
	Ein $R$-Modul $I$ heißt \bet{injektiv}\index{injektiver Modul}, falls es zu jedem $j \colon M \hookrightarrow N$ und $f \colon M \to I$ ein $\hat{f} \colon N \to I$ gibt mit $f = \hat{f} \circ j$, also
	\[
		\begin{tikzcd}
			\mbox{ } & N \dlar[dashed,"\hat{f}"'] \\
			I & M \lar["f"'] \uar["j"',hook]
		\end{tikzcd}
	\]
\end{definition}

\begin{bemerkung}
	Jede kurze exakte Folge \(
		\begin{tikzcd}[cramped,sep=small]
			0 \rar & M \rar & M' \rar & M'' \rar & 0
		\end{tikzcd}
	\)
	mit $M$ injektiv spaltet.
\end{bemerkung}

\begin{satz}[{name=[Existenz von injektiven Moduln]}]
	Sei $M$ ein $R$-Modul.
	Dann gibt es einen injektiven $R$-Modul $I$ mit einer injektiven Abbildung $M \hookrightarrow I$.
\end{satz}

\begin{definition}
	Eine injektive Auflösung einer $R$-Moduls $M$ ist eine lange exakte Folge 
	\[
		\begin{tikzcd}
			0 \rar & M \rar[hook] & I_0 \rar["\partial_0"] & I_1 \rar["\partial_1"] & \ldots 
		\end{tikzcd}
	\]
	mit $I_i$ injektiv.
\end{definition}

\begin{korollarB}
	Zu jedem $M$ gibt es eine injektive Auflösung.
\end{korollarB}

\begin{bemerkung}
	Sei $(I_*,j)$ eine injektive Auflösung von $W$. 
	Dann gilt
	\[
		\Ext_R^n(M,W) \cong H^n \enbrace*{\Hom_R(M,I_*)}
	\]
\end{bemerkung}
\newpage
% section 12 (end)

\section{Das Tensorprodukt und $\Hom$ als Kettenkomplexe} % (fold)
\label{sec:das_tensorprodukt_und_hom_als_kettenkomplexe}

\begin{definition}
	Sei $R$ ein Ring, $C_*$ ein Kettenkomplex von $R$-Rechtsmoduln und $D_*$ ein Kettenkomplex von $R$-Linksmoduln.
	Dann definieren wir den $\mathbb{Z}$-Kettenkomplex $(C \otimes_R D)_*$ durch 
	\[
		(C \otimes_R D)_n := \bigoplus_{p+q=n} C_p \otimes_R D_q
	\]
	mit Randabbildung definiert durch
	\(
		\partial^{C \otimes D} (c \otimes d) := \partial^C(c) \otimes d + (-1)^{\abs*{c}} c \otimes \partial^D(d)
	\).
\end{definition}

\begin{bemerkung}
	\leavevmode
	\begin{itemize}[itemsep=0pt]
		\item Sind $C$ und $D$ $\mathbb{N}$-graduiert, so ist $C \otimes_R D$ auch $\mathbb{N}$-graduiert. 
		\item Ist $R$ kommutativ, so ist $C \otimes_R D$ ein $R$-Kettenkomplex.
	\end{itemize}
\end{bemerkung}

\begin{definition}
	Sei $C_*$ ein $R$-Kettenkomplex und $\tilde{D}^*$ ein $R$-Kokettenkomplex.
	Wir definieren den $\mathbb{Z}$-Kokettenkomplex $\Hom_R(C,\tilde{D})^*$ durch 
	\[
		\enbrace*{\Hom_R(C,\tilde{D})}^n := \bigoplus_{p+q=n} \Hom_R \enbrace*{C_p, \tilde{D}^q}
	\]
	und die Randabbildung durch
	\(
		\partial f = \partial^{\tilde{D}} \circ f - (-1)^{\abs*{f}} f \circ \partial^C
	\).
\end{definition}

\begin{bemerkung}
	\leavevmode
	\begin{itemize}[itemsep=0pt]
		\item Sind $C$ und $\tilde{D}$ $\mathbb{N}$-graduiert, so ist  $\Hom_R(C,\tilde{D})$ auch $\mathbb{N}$-graduiert.
		\item Ist $R$ kommutativ, so ist $\Hom_R(C,\tilde{D})$ ein $R$-Kokettenkomplex.
	\end{itemize}
\end{bemerkung}

\begin{lemma}[label=projaufl:iso:hom]
	Sei $C_*$ ein $R$-Kettenkomplex ($\mathbb{N}$-graduiert), dessen Kettenmodule projektiv sind.
	Sei $V$ ein $R$-Rechtsmodul und $(P_*,\epsilon)$ eine projektive Auflösung von $V$ durch $R$-Rechtsmoduln.
	Dann induziert $\epsilon \otimes \id \colon P_* \otimes_R C_* \to V \otimes_R C_*$ einen Isomorphismus in Homologie.
\end{lemma}
\begin{beweis}
	Sei $v \in (P \otimes C)_n$.
	Dann schreiben wir $v=\sum_{i=0}^{n} v_i$ mit $v_i \in P_i \otimes C_{n-i}$
	Die $v_i$ sind durch $v$ eindeutig bestimmt.
	Es ist
	\[
		\partial^{P \otimes C}(v)= \sum_{i=0}^{n-1} \Underbracket{(-1)^i \id \otimes \partial^C(v_i) + (\partial^P \otimes \id)(v_{i+1})}{\in P_i \otimes C_{n-1-i}}
	\]
	die eindeutige Darstellung von $\partial^{P \otimes C}(v)$.
	\begin{figure}[bt]
		\centering{
		\begin{tikzcd}[row sep=2.5em]
			\vdots \dar & \vdots \dar & \vdots \dar & \vdots \dar \\
			V \otimes C_2 \dar["\id \otimes \partial^C"'] & P_0 \otimes C_2 \lar["\epsilon \otimes \id"']\dar & P_1 \otimes C_2\lar\dar & P_2 \otimes C_2 \lar \dar & \cdots \lar\\
			V \otimes C_1 \dar["\id \otimes \partial^C"'] & P_0 \otimes C_1 \lar["\epsilon \otimes \id"']\dar & P_1 \otimes C_1\lar\dar & P_2 \otimes C_1 \lar\dar & \cdots \lar\\
			V \otimes C_0 & P_0 \otimes C_0 \lar["\epsilon \otimes \id"']& P_1 \otimes C_0\lar & P_2 \otimes C_0 \lar & \cdots \lar
		\end{tikzcd}
		\caption{Diagramm zum Beweis von \cref{projaufl:iso:hom}\label{fig:tensor:ketten}}
		}
	\end{figure}%
	Insbesondere gilt: 
	\begin{equation}
		\partial^{P \otimes C}v =0 \iff (-1)^{i+1} (\id \otimes \partial^C)(v_i)= (\partial^P \otimes \id)(v_{i+1}) \label{eq:13:1}\tag{\#}
	\end{equation}
	Da die $C_i$ projektiv sind, ist $\otimes_R C_i$ exakt und für jedes $i$ ist 
	\begin{equation}
		\begin{tikzcd}[sep=large]
			V \otimes_R C_i & P_0 \otimes_R C_i \lar["\epsilon \otimes \id"'] & P_1 \otimes C_i \lar["\partial^P \otimes \id"'] & \ldots \lar["\partial^P \otimes \id"']
		\end{tikzcd}
		\label{eq:13:2}\tag{\S}
	\end{equation}
	exakt.
	Sei $x \in H_n(V \otimes C_*)$.
	Sei $a \in V \otimes C_n$ mit $\id \otimes \partial^C(a)=0$ und $x=[a]$.
	Wähle $v_0 \in P_0 \otimes C_n$ mit $\epsilon \otimes \id(v_0)=a$.
	Behauptung: Es gibt $v_i \in P_i \otimes C_{n-i}$ mit 
	\(
		\partial^{P \otimes C} \enbrace[\big]{\sum_{i=0}^{n} v_i } = 0
	\).
	Wir beweisen dies per Induktion: Seien $v_0, \ldots ,v_j$ bereits konstruiert mit \eqref{eq:13:1} für $i<j$.
	Dann ist $(\partial^P \otimes \id)(\id \otimes \partial^C)(v_j) = \pm \enbrace*{\partial^P \otimes \id} \enbrace*{\partial^P \otimes \id}(v_{j-1})=0$.
	Wegen \eqref{eq:13:2} gibt es $v_{j+1} \in P_{j+1} \otimes C_{n-j-1}$ mit 
	\[
		\enbrace*{\partial^P \otimes \id}(v_{j+1}) = \enbrace*{\id \otimes \partial^C}(v_j)
	\]
	Damit ist die Behauptung bewiesen.
	Es folgt $(\epsilon \otimes \id)_* \enbrace*{\benbrace*{\sum_{i=0} v_i}} = [a]$.
	Daher ist $\enbrace*{\epsilon \otimes \id}_*$ surjektiv. 
	Die Injektivität erhält man genauso.
\end{beweis}

\begin{lemma}
	Sei $C_*$ ein $R$-Kettenkomplex ($\mathbb{N}$-graduiert), dessen Kettenmodule projektiv sind.
	Sei $(I^*,j)$ eine injektive Auflösung des $R$-Moduls $W$.
	Dann induziert $j$ einen Isomorphismus 
	\[
		H^*\enbrace*{\Hom_R(C_*,W)} \to H^* \enbrace*{\Hom_R(C_*,I^*)}
	\]
\end{lemma}

\begin{bemerkung}
	Für den Isomorphismus $H_*(P \otimes_R C_*) \to H_*(V \otimes_R C_*)$ war es wichtig, dass die Homologie von $\enbrace*{P_* \otimes C_i,\partial^P \otimes \id}$ in Grad $0$ konzentriert ist und damit genau $V \otimes C_i$ ist (dies entspricht den Spalten in \cref{fig:tensor:ketten}).
	Umgekehrt ist die Homologie von $\enbrace*{P_p \otimes C_*, \id \otimes \partial^C}$ isomorph zu $P_p \otimes H_*(C_*)$.
	Wir erhalten also Komplexe $\enbrace*{P_* \otimes H_n(C_*), \partial^P \otimes \id}$ deren Homologie genau $\Tor^R_* \enbrace*{V, H_n(C_*)}$ ist.
	Dies stellt einen Zusammehang zwischen $H_*(V \otimes C_*)$ und $\Tor_*^R \enbrace*{V, H_*(C_*)}$ her.
\end{bemerkung}
% section das_tensorprodukt_und_hom_als_kettenkomplexe (end)





\cleardoubleoddemptypage
\appendix
\pagenumbering{Alph}
\setcounter{page}{1}

\printindex
\printbibliography
\listoffigures
\todototoc
\listoftodos[To-do's und andere Baustellen]
\end{document}
